% Trigonométrie — Mathématiques 1re Tech — Cours signé OnBuch+
% Programme officiel MINESEC. Compiler avec : tectonic N10-trigonometrie.tex
\newcommand{\DOCMATIERE}{Mathématiques}
\newcommand{\DOCNIVEAU}{1re Tech}
\newcommand{\DOCMODULE}{Mathématiques – classe de Première technique}
\newcommand{\DOCLECON}{N10}
\newcommand{\DOCTITRE}{Trigonométrie}
% OnBuch+ — préambule « cours signés » ENRICHI (compilé avec tectonic / XeLaTeX)
% Système visuel « Pop » d'OnBuch+ : fond crème (jamais blanc), bordures encre
% épaisses, ombres « dures » décalées, titres Archivo Black en capitales.
% Version MATHÉMATIQUES : graphes (pgfplots/tikz), boîtes pédagogiques
% (définition, propriété, méthode, attention, le savais-tu, exercice résolu,
% exercice, TP, bilan), page de garde et signature OnBuch+ ancrée en bas.
%
% Macros attendues dans main.tex AVANT \input{preamble} :
%   \DOCMATIERE  \DOCNIVEAU  \DOCMODULE  \DOCLECON (numéro)  \DOCTITRE
\documentclass[11pt]{article}

\usepackage{fontspec}
\usepackage[french]{babel}
\usepackage[a4paper,margin=2cm,top=3.4cm,bottom=2.8cm,headheight=1.4cm,footskip=1.3cm]{geometry}
\usepackage{amsmath,amssymb,amsfonts}
\usepackage{mathtools} % \xmapsto, \coloneqq... souvent utilisés par les modèles
\usepackage{cancel} % simplifications barrées
% \abs{z} (module, valeur absolue) et \norm{u} (norme), utilisés par les modèles
\providecommand{\abs}{}\let\abs\relax\DeclarePairedDelimiter{\abs}{\lvert}{\rvert}
\providecommand{\norm}{}\let\norm\relax\DeclarePairedDelimiter{\norm}{\lVert}{\rVert}
\usepackage[table]{xcolor} % [table] : fournit aussi \rowcolors (alternance de lignes)
\usepackage{pagecolor}
\usepackage{tikz}
\usetikzlibrary{arrows.meta,positioning,calc,shapes.geometric,shapes.misc,shapes.arrows,decorations.pathmorphing,decorations.pathreplacing,decorations.markings,patterns,shadows,mindmap,trees,fit,backgrounds,matrix,intersections,angles,quotes}
\usepackage{pgfplots}
\usepackage{pgf-pie} % diagrammes circulaires \pie (statistiques)
\pgfplotsset{compat=1.18}
\usepgfplotslibrary{fillbetween,groupplots} % aires entre courbes (calcul intégral)
\usepackage{enumitem}
\usepackage{fancyhdr}
% [most] charge toutes les bibliothèques tcolorbox (listings, minted, raster,
% documentation...) et gonfle inutilement la mémoire TeX sur un document long
% (~300 boîtes) : on ne charge que ce qui est réellement utilisé.
\usepackage[skins,breakable]{tcolorbox}
\usepackage{graphicx}
\usepackage{colortbl}
\usepackage{booktabs}
\usepackage{tabularx}
\usepackage{diagbox} % case d'en-tête diagonale des tableaux à double entrée
% Colonnes extensibles centrée / gauche / droite (C, L, R), souvent utilisées par les modèles
\newcolumntype{C}{>{\centering\arraybackslash}X}
\newcolumntype{L}{>{\raggedright\arraybackslash}X}
\newcolumntype{R}{>{\raggedleft\arraybackslash}X}
\usepackage{array}
\usepackage{multicol}
\usepackage{siunitx}
% parse-numbers=false : accepte aussi \SI{2,5 \times 10^{-3}}{...}
\sisetup{locale=FR,output-decimal-marker={,},group-minimum-digits=4,parse-numbers=false}
\DeclareSIUnit\ohm{\ensuremath{\symup{\Omega}}} % U+2126 absent de la police
\usepackage{qrcode}
% \degree : commande fréquemment utilisée par les modèles pour un angle ou une
% température en dehors de \SI{}{} (non fournie par les paquets chargés).
\providecommand{\degree}{^{\circ}}
% \qed (fin de démonstration), utilisé par les modèles sans amsthm
\providecommand{\qed}{\hfill$\square$}
% \barre : double barre des valeurs interdites dans un tableau de variations (tkz-tab)
\providecommand{\barre}{\ensuremath{\|}}
% environnement proof (amsthm non chargé) utilisé par les modèles
\ifdefined\proof\else\newenvironment{proof}[1][Démonstration]{\par\noindent\textit{#1.}\ }{\qed\par}\fi
\usepackage{lastpage}
\usepackage{hyperref}
\hypersetup{hidelinks}

% ============================================================
% Palette « Pop » OnBuch+ — mêmes hex que la classe P (plus_ui.dart)
% ============================================================
\definecolor{poporange}{HTML}{F59321}
\definecolor{popdark}{HTML}{E07A0C}
\definecolor{popcream}{HTML}{FFF6E8}
\definecolor{popink}{HTML}{1C1714}
\definecolor{poppurple}{HTML}{7A5AE0}
\definecolor{popgreen}{HTML}{2E9E6A}
\definecolor{poppink}{HTML}{E0457B}
\definecolor{popblue}{HTML}{2F7DE1}
\definecolor{popgold}{HTML}{C98A12}
\definecolor{popmuted}{HTML}{8A7F76}
\definecolor{popline}{HTML}{EADFD2}
\definecolor{popgray}{HTML}{8A7F76}
\colorlet{popgrey}{popgray}
% Alias tolérants (le modèle nomme parfois les couleurs différemment)
\colorlet{popwhite}{white}
\colorlet{popblack}{popink}
\colorlet{popred}{poppink}
\colorlet{popyellow}{popgold}
% Teintes claires pour fonds de tableaux / zones
\colorlet{popblueL}{popblue!10!white}
\colorlet{poporangeL}{poporange!14!white}
\colorlet{popgreenL}{popgreen!12!white}
\colorlet{poppurpleL}{poppurple!12!white}
\colorlet{poppinkL}{poppink!12!white}
\colorlet{popgoldL}{popgold!14!white}

\pagecolor{popcream}

% ============================================================
% Typographie
% ============================================================
\defaultfontfeatures{Ligatures=TeX}
\setmainfont{PlusJakartaSans-Medium.ttf}[Path=../../../fonts/,BoldFont=PlusJakartaSans-Bold.ttf]
% Maths en sans-serif (Fira Math) avec chiffres et lettres droites de Plus
% Jakarta, pour que les formules se fondent dans le texte.
\usepackage[math-style=ISO,bold-style=ISO]{unicode-math}
\setmathfont{FiraMath-Regular.otf}[Path=../../../fonts/]
\setmathfont{PlusJakartaSans-Medium.ttf}[Path=../../../fonts/,range={up/num,up/{Latin,latin},\mathup}]
\usepackage{icomma} % virgule décimale sans espace en mode math
% Caractères mathématiques Unicode absents des polices de texte (Plus Jakarta,
% Archivo Black) : rendus par leur équivalent mathématique, en texte comme en
% formule — sinon ils s'affichent en carré vide (ex. « Arithmétique dans ℤ »).
\usepackage{newunicodechar}
\newunicodechar{ℝ}{\ensuremath{\mathbb{R}}}
\newunicodechar{ℤ}{\ensuremath{\mathbb{Z}}}
\newunicodechar{ℂ}{\ensuremath{\mathbb{C}}}
\newunicodechar{ℕ}{\ensuremath{\mathbb{N}}}
\newunicodechar{ℚ}{\ensuremath{\mathbb{Q}}}
\newunicodechar{𝔻}{\ensuremath{\mathbb{D}}}
\newunicodechar{α}{\ensuremath{\alpha}}
\newunicodechar{β}{\ensuremath{\beta}}
\newunicodechar{γ}{\ensuremath{\gamma}}
\newunicodechar{δ}{\ensuremath{\delta}}
\newunicodechar{ε}{\ensuremath{\varepsilon}}
\newunicodechar{θ}{\ensuremath{\theta}}
\newunicodechar{λ}{\ensuremath{\lambda}}
\newunicodechar{μ}{\ensuremath{\mu}}
\newunicodechar{σ}{\ensuremath{\sigma}}
\newunicodechar{τ}{\ensuremath{\tau}}
\newunicodechar{φ}{\ensuremath{\varphi}}
\newunicodechar{ψ}{\ensuremath{\psi}}
\newunicodechar{ω}{\ensuremath{\omega}}
\newunicodechar{Σ}{\ensuremath{\Sigma}}
% majuscules produites par \MakeUppercase dans les titres (α-aminés -> Α)
\newunicodechar{Α}{\ensuremath{\alpha}}
\newunicodechar{Β}{\ensuremath{\beta}}
\newunicodechar{Γ}{\ensuremath{\Gamma}}
\newunicodechar{Δ}{\ensuremath{\Delta}}
\newunicodechar{Ε}{\ensuremath{\varepsilon}}
\newunicodechar{Θ}{\ensuremath{\Theta}}
\newunicodechar{Λ}{\ensuremath{\Lambda}}
\newunicodechar{Μ}{\ensuremath{\mu}}
\newunicodechar{Τ}{\ensuremath{\tau}}
\newunicodechar{Φ}{\ensuremath{\Phi}}
\newunicodechar{Ψ}{\ensuremath{\Psi}}
\newunicodechar{Ω}{\ensuremath{\Omega}}
\newunicodechar{↦}{\ensuremath{\mapsto}}
\newunicodechar{∈}{\ensuremath{\in}}
\newunicodechar{∉}{\ensuremath{\notin}}
\newunicodechar{∧}{\ensuremath{\wedge}}
\newunicodechar{⇔}{\ensuremath{\Leftrightarrow}}
\newunicodechar{⟺}{\ensuremath{\Longleftrightarrow}}
\newunicodechar{⇒}{\ensuremath{\Rightarrow}}
\newunicodechar{⟹}{\ensuremath{\Longrightarrow}}
\newunicodechar{∘}{\ensuremath{\circ}}
\newunicodechar{∪}{\ensuremath{\cup}}
\newunicodechar{∩}{\ensuremath{\cap}}
\newunicodechar{≡}{\ensuremath{\equiv}}
\newunicodechar{⊂}{\ensuremath{\subset}}
\newunicodechar{⊥}{\ensuremath{\perp}}
\newunicodechar{∃}{\ensuremath{\exists}}
\newunicodechar{∀}{\ensuremath{\forall}}
\newunicodechar{∛}{\ensuremath{\sqrt[3]{\,}}}
\newunicodechar{∜}{\ensuremath{\sqrt[4]{\,}}}
\newunicodechar{⃗}{\ensuremath{{}^{\to}}}
\newunicodechar{ⁿ}{\ifmmode^{n}\else\textsuperscript{\ensuremath{n}}\fi}
\newunicodechar{ˣ}{\ifmmode^{x}\else\textsuperscript{\ensuremath{x}}\fi}
\newunicodechar{ᵇ}{\ifmmode^{b}\else\textsuperscript{\ensuremath{b}}\fi}
\newunicodechar{ʳ}{\ifmmode^{r}\else\textsuperscript{\ensuremath{r}}\fi}
\newunicodechar{ᵘ}{\ifmmode^{u}\else\textsuperscript{\ensuremath{u}}\fi}
\newunicodechar{ʸ}{\ifmmode^{y}\else\textsuperscript{\ensuremath{y}}\fi}
\newunicodechar{ᵃ}{\ifmmode^{a}\else\textsuperscript{\ensuremath{a}}\fi}
\newunicodechar{ᵉ}{\ifmmode^{e}\else\textsuperscript{\ensuremath{e}}\fi}
\newunicodechar{ᵏ}{\ifmmode^{k}\else\textsuperscript{\ensuremath{k}}\fi}
\newunicodechar{ᵖ}{\ifmmode^{p}\else\textsuperscript{\ensuremath{p}}\fi}
\newunicodechar{ᴺ}{\ifmmode^{N}\else\textsuperscript{\ensuremath{N}}\fi}
\newunicodechar{ᶜ}{\ifmmode^{c}\else\textsuperscript{\ensuremath{c}}\fi}
\newunicodechar{ʰ}{\ifmmode^{h}\else\textsuperscript{\ensuremath{h}}\fi}
\newunicodechar{ᵠ}{\ifmmode^{q}\else\textsuperscript{\ensuremath{q}}\fi}
\newunicodechar{ₙ}{\ifmmode_{n}\else\textsubscript{\ensuremath{n}}\fi}
\newunicodechar{ₐ}{\ifmmode_{a}\else\textsubscript{\ensuremath{a}}\fi}
\newunicodechar{ᵢ}{\ifmmode_{i}\else\textsubscript{\ensuremath{i}}\fi}
\newunicodechar{ᵣ}{\ifmmode_{r}\else\textsubscript{\ensuremath{r}}\fi}
\newunicodechar{ₖ}{\ifmmode_{k}\else\textsubscript{\ensuremath{k}}\fi}
\newunicodechar{ᵦ}{\ifmmode_{\beta}\else\textsubscript{\ensuremath{\beta}}\fi}
\newunicodechar{ₓ}{\ifmmode_{x}\else\textsubscript{\ensuremath{x}}\fi}
\newfontfamily\popdisplay{ArchivoBlack-Regular.ttf}[Path=../../../fonts/]
\newfontfamily\poplogo{SpaceGrotesk-Bold.ttf}[Path=../../../fonts/]
\newfontfamily\popsemi{PlusJakartaSans-SemiBold.ttf}[Path=../../../fonts/]
\newfontfamily\popxbold{PlusJakartaSans-ExtraBold.ttf}[Path=../../../fonts/]
\newfontfamily\popxboldls{PlusJakartaSans-ExtraBold.ttf}[Path=../../../fonts/,LetterSpace=3.0]

\setlist[enumerate]{leftmargin=1.7em,itemsep=5pt,topsep=4pt}
\setlist[itemize]{leftmargin=1.7em,itemsep=4pt,topsep=4pt,label={\color{poporange}\textbullet}}
\setlength{\parindent}{0pt}
\setlength{\parskip}{5pt}
\renewcommand{\arraystretch}{1.25}

% pgfplots : style Pop par défaut
\pgfplotsset{
  every axis/.append style={axis line style={popink,line width=1pt},tick style={popink},
    grid style={popline,line width=0.5pt},label style={font=\small},tick label style={font=\footnotesize}},
  popcurve/.style={poporange,line width=1.8pt,no markers,smooth},
}
% Même style disponible dans un \draw TikZ simple (hors pgfplots)
\tikzset{popcurve/.style={poporange,line width=1.8pt}}

% ============================================================
% Pill / squiggle
% ============================================================
\newcommand{\poppill}[3]{%
  \tikz[baseline=-0.55ex]\node[draw=popink,line width=1.2pt,rounded corners=2.6pt,%
    fill=#2,inner xsep=6.5pt,inner ysep=3pt]{%
    \popxboldls\fontsize{7}{7}\selectfont\color{#3}\MakeUppercase{#1}};%
}
\newcommand{\popsquiggle}[1]{%
  \tikz[baseline=-0.4ex]{
    \draw[#1,line width=2.6pt,line cap=round]
      plot [smooth] coordinates {(0,0.09) (0.34,0.01) (0.68,0.21) (1.02,0.01) (1.36,0.10)};
  }%
}
% Mot-clé surligné (marqueur orange sous le texte)
\newcommand{\cle}[1]{{\popxbold\color{popdark}#1}}

% ============================================================
% En-tête / pied
% ============================================================
\pagestyle{fancy}
\fancyhf{}
\renewcommand{\headrulewidth}{0pt}
\renewcommand{\footrulewidth}{0pt}
\lhead{\raisebox{-0.15ex}{\poplogo\fontsize{13}{13}\selectfont\color{popink}OnBuch\color{poporange}+}}
\rhead{\poppill{\DOCMATIERE\ \textbullet\ \DOCNIVEAU\ \textbullet\ Leçon \DOCLECON}{white}{popink}}
\lfoot{\popxbold\fontsize{7.6}{7.6}\selectfont\color{popmuted}\MakeUppercase{Cours signé OnBuch+}\ \textbullet\ {\popsemi\DOCTITRE}}
\rfoot{\tikz[baseline=-0.6ex]\node[rounded corners=8.5pt,fill=popink,minimum height=17pt,minimum width=17pt,inner xsep=5pt,inner sep=0]{\popxbold\fontsize{8}{8}\selectfont\color{popcream}\thepage\,/\,\pageref*{LastPage}};}

% ============================================================
% Titres
% ============================================================
\newcommand{\chaptitre}[2]{%
  \begin{tcolorbox}[enhanced,colback=poporange,colframe=popink,boxrule=2.6pt,arc=11pt,
      drop shadow=popink,left=15pt,right=15pt,top=12pt,bottom=11pt,before skip=3pt,after skip=18pt]
    \poppill{#2}{popink}{popcream}\par\vspace{9pt}
    {\raggedright\hyphenpenalty=10000\exhyphenpenalty=10000\popdisplay\fontsize{22}{24}\selectfont\color{popcream}\MakeUppercase{#1}\par}
  \end{tcolorbox}
}
% Section numérotée : pastille orange avec le numéro + titre Archivo
\newcounter{popsec}
\newcommand{\coursec}[1]{%
  \stepcounter{popsec}\setcounter{popsub}{0}%
  \par\vspace{16pt}\noindent
  \begin{minipage}{\linewidth}
  \tikz[baseline=-0.7ex]\node[draw=popink,line width=1.6pt,fill=poporange,rounded corners=4pt,minimum size=22pt,inner sep=0]{\popdisplay\fontsize{12}{12}\selectfont\color{popcream}\thepopsec};%
  \hspace{8pt}{\popdisplay\fontsize{15}{17}\selectfont\color{popink}\MakeUppercase{#1}}\par
  \vspace{4pt}\noindent{\color{poporange}\rule{36pt}{3.6pt}}
  \end{minipage}\par\nopagebreak\vspace{8pt}
}
\newcounter{popsub}
\newcommand{\courssub}[1]{%
  \stepcounter{popsub}%
  \par\vspace{9pt}\noindent{\popxbold\fontsize{11.5}{13}\selectfont\color{popink}{\color{poporange}\thepopsec.\thepopsub}\hspace{6pt}#1}\par\nopagebreak\vspace{4pt}
}

% ============================================================
% Boîtes pédagogiques — même recette : fond blanc, bordure épaisse colorée,
% ombre dure encre, badge plein. Titre optionnel [#1].
% ============================================================
\tcbset{popbase/.style={breakable,enhanced,colback=white,boxrule=2.3pt,arc=9pt,
  shadow={3pt}{-3pt}{0pt}{popink},left=14pt,right=14pt,top=11pt,bottom=11pt,before skip=11pt,after skip=15pt,
  fontupper=\popsemi\fontsize{10.6}{14.5}\selectfont\color{popink},
  fontlower=\popsemi\fontsize{10.6}{14.5}\selectfont\color{popink}}}
% Coche dessinée (le glyphe ✓ n'existe pas dans les polices du document)
\newcommand{\popcheck}[1]{\tikz[baseline=-0.5ex]\draw[#1,line width=1.6pt,line cap=round,line join=round](-0.13,0.0)--(-0.04,-0.09)--(0.13,0.1);}
\newcommand{\popboxhead}[3]{\poppill{#1}{#2}{white}\ifx\relax#3\relax\else\hspace{6pt}{\popxbold\fontsize{10.6}{12}\selectfont\color{popink}#3}\fi\par\vspace{7pt}}

\newtcolorbox{aretenir}[1][]{popbase,colframe=popgreen,before upper={\popboxhead{À retenir}{popgreen}{#1}}}
\newtcolorbox{exemplebox}[1][]{popbase,colframe=poporange,before upper={\popboxhead{Exemple}{poporange}{#1}}}
\newtcolorbox{definition}[1][]{popbase,colframe=popblue,before upper={\popboxhead{Définition}{popblue}{#1}}}
\newtcolorbox{propriete}[1][]{popbase,colframe=poppurple,before upper={\popboxhead{Propriété}{poppurple}{#1}}}
\newtcolorbox{methode}[1][]{popbase,colframe=popgold,before upper={\popboxhead{Méthode}{popgold}{#1}}}
\newtcolorbox{attention}[1][]{popbase,colframe=poppink,before upper={\popboxhead{Attention}{poppink}{#1}}}
\newtcolorbox{experience}[1][]{popbase,colframe=popblue,colback=popblueL,before upper={\popboxhead{Expérience}{popblue}{#1}}}
% « Le savais-tu ? » : carte violette pleine (contexte, culture, Cameroun)
\newtcolorbox{savaistu}[1][]{popbase,colback=poppurple,colframe=popink,
  fontupper=\popsemi\fontsize{10.4}{14.2}\selectfont\color{white},
  before upper={\poppill{Le savais-tu\,?}{white}{poppurple}\ifx\relax#1\relax\else\hspace{6pt}{\popxbold\color{white}#1}\fi\par\vspace{7pt}}}
% Exercice résolu : énoncé en haut, \tcblower puis la solution sur fond crème
\newcounter{popexo}\newcounter{popexr}
\newtcolorbox{exoresolu}[1][]{popbase,colframe=poporange,colbacklower=popcream,
  segmentation style={popink,line width=1.2pt,dashed},
  before upper={\stepcounter{popexr}\popboxhead{Exercice résolu \thepopexr}{poporange}{#1}},
  before lower={\poppill{Solution}{popink}{popcream}\par\vspace{6pt}}}
% Exercice (non résolu en place) — la correction est dans la partie Corrigés
\newtcolorbox{exercice}[1][]{popbase,colframe=popink,
  before upper={\stepcounter{popexo}\popboxhead{Exercice \thepopexo}{popink}{#1}}}
% Corrigé
\newtcolorbox{corrige}[1][]{popbase,colframe=popgreen,colback=popgreenL,
  before upper={\popboxhead{Corrigé}{popgreen}{#1}}}
% Objectifs / prérequis / bilan
\newtcolorbox{objectifs}{popbase,colframe=popink,colback=poporangeL,
  before upper={\popboxhead{Objectifs de la leçon}{popink}{}}}
\newtcolorbox{prerequis}{popbase,colframe=popink,colback=popblueL,
  before upper={\popboxhead{Prérequis}{popblue}{}}}
\newtcolorbox{bilan}{popbase,colframe=popink,colback=white,boxrule=2.8pt,
  before upper={\popboxhead{Fiche bilan}{poporange}{L'essentiel à connaître pour l'examen}}}
% Figure encadrée avec légende
\newtcolorbox{popfigure}[1][]{enhanced,breakable=false,colback=white,colframe=popline,boxrule=1.4pt,arc=8pt,
  left=10pt,right=10pt,top=10pt,bottom=8pt,before skip=10pt,after skip=12pt,halign=center,
  lower separated=false,halign lower=center,
  fontlower=\popsemi\fontsize{9}{11.5}\selectfont\color{popmuted},
  #1}
\newcommand{\legende}[1]{\par\vspace{6pt}{\popsemi\fontsize{9}{11.5}\selectfont\color{popmuted}#1}}

% ============================================================
% Signature OnBuch+
% ============================================================
\newcommand{\poplogoinline}[1]{{\poplogo\fontsize{#1}{#1}\selectfont\color{popink}OnBuch\color{poporange}+}}
\newcommand{\signatureonbuch}{%
  \begin{tcolorbox}[enhanced,colback=popink,colframe=popink,boxrule=2.2pt,arc=10pt,
      drop shadow=poporange,left=14pt,right=14pt,top=10pt,bottom=10pt,before skip=0pt,after skip=0pt]
    \tikz[baseline=-0.7ex]\node[circle,fill=popgreen,draw=popcream,line width=1.3pt,minimum size=22pt,inner sep=0]{\popcheck{white}};%
    \hspace{9pt}\begin{minipage}[c]{0.62\linewidth}
      {\popxbold\fontsize{12}{13}\selectfont\color{popcream}Signé\ \textbullet\ {\poplogo OnBuch\color{poporange}+}}\par\vspace{2pt}
      {\raggedright\popsemi\fontsize{8}{10.5}\selectfont\color{popline}\MakeUppercase{Équipe pédagogique OnBuch+}\\\MakeUppercase{Conforme au programme officiel MINESEC}\par}
    \end{minipage}\hfill
    {\poplogo\fontsize{22}{22}\selectfont\color{popcream}OnBuch\color{poporange}+}
  \end{tcolorbox}%
}

% ============================================================
% Page de garde
% #1 titre · #2 matière · #3 niveau · #4 module · #5 n° leçon · #6 durée ·
% #7 description / objectifs (texte ou itemize)
% ============================================================
\newcommand{\pagegarde}[7]{%
  \thispagestyle{empty}
  \newgeometry{a4paper,margin=1.7cm,top=1.6cm,bottom=1.4cm}%
  \begin{tikzpicture}[remember picture,overlay]
    \fill[poporange!18] ([xshift=-3.2cm,yshift=-3.6cm]current page.north east) circle (4.6cm);
    \fill[poppurple!14] ([xshift=2.4cm,yshift=7.6cm]current page.south west) circle (3.8cm);
  \end{tikzpicture}%
  {\poplogo\fontsize{30}{30}\selectfont\color{popink}OnBuch\color{poporange}+}\hfill
  \raisebox{0.6ex}{\poppill{Cours signé \textbullet\ Premium}{popink}{popcream}}\par
  \vspace{1pt}\popsquiggle{poppurple}\par
  \vspace{8pt}
  \begin{tcolorbox}[enhanced,colback=poporange,colframe=popink,boxrule=2.8pt,arc=13pt,
      drop shadow=popink,left=18pt,right=18pt,top=12pt,bottom=13pt,after skip=10pt]
    \poppill{#2 \textbullet\ #3}{popink}{popcream}\hspace{5pt}\poppill{#4}{white}{popink}\par\vspace{8pt}
    {\popdisplay\fontsize{14}{14}\selectfont\color{popink}LEÇON #5}\par\vspace{4pt}
    {\raggedright\hyphenpenalty=10000\exhyphenpenalty=10000\popdisplay\fontsize{25}{27}\selectfont\color{popcream}\MakeUppercase{#1}\par}
    \vspace{8pt}
    {\popxbold\fontsize{9.5}{9.5}\selectfont\color{popink}\MakeUppercase{Durée conseillée : #6}}
  \end{tcolorbox}
  \begin{tcolorbox}[enhanced,colback=white,colframe=popink,boxrule=2.2pt,arc=10pt,
      drop shadow=popink,left=16pt,right=16pt,top=10pt,bottom=10pt,after skip=8pt]
    \poppill{Dans cette leçon}{poppurple}{white}\par\vspace{5pt}
    \popsemi\fontsize{9.3}{12.6}\selectfont\color{popink}#7
  \end{tcolorbox}
  \noindent\begin{minipage}[t]{0.487\linewidth}
    \begin{tcolorbox}[enhanced,colback=popgreen,colframe=popink,boxrule=2.2pt,arc=10pt,
        drop shadow=popink,left=11pt,right=11pt,top=10pt,bottom=10pt,height=3.1cm,valign=center]
      \tcbox[colback=white,colframe=popink,boxrule=1.3pt,arc=5pt,boxsep=0pt,nobeforeafter,
        left=3pt,right=3pt,top=3pt,bottom=3pt,valign=center]{\qrcode[height=1.9cm]{https://play.google.com/store/apps/details?id=com.luvvix.onbuch&hl=fr}}%
      \hspace{8pt}\begin{minipage}[c]{0.5\linewidth}
        {\popdisplay\fontsize{11}{12.5}\selectfont\color{white}\MakeUppercase{Télécharge OnBuch}}\par\vspace{4pt}
        {\raggedright\popsemi\fontsize{8.4}{11}\selectfont\color{white}Gratuit \textbullet\ Google Play\\Tuteur IA, annales, concours, résultats\par}
      \end{minipage}
    \end{tcolorbox}
  \end{minipage}\hfill
  \begin{minipage}[t]{0.487\linewidth}
    \begin{tcolorbox}[enhanced,colback=poppurple,colframe=popink,boxrule=2.2pt,arc=10pt,
        drop shadow=popink,left=11pt,right=11pt,top=10pt,bottom=10pt,height=3.1cm,valign=center]
      \tikz[baseline=-0.7ex]\node[circle,fill=white,draw=popink,line width=1.3pt,minimum size=1.6cm,inner sep=0]{\popdisplay\fontsize{24}{24}\selectfont\color{poppurple}+};%
      \hspace{8pt}\begin{minipage}[c]{0.6\linewidth}
        {\popdisplay\fontsize{11}{12.5}\selectfont\color{white}\MakeUppercase{Passe à OnBuch+}}\par\vspace{4pt}
        {\raggedright\popsemi\fontsize{8.4}{11}\selectfont\color{white}Tous les cours signés, Tuteur IA illimité, fiches PDF — onglet OnBuch+ de l'app\par}
      \end{minipage}
    \end{tcolorbox}
  \end{minipage}
  \vspace{8pt}
  \popcontenu
  \vfill
  \signatureonbuch
  \restoregeometry
  \newpage
}

% Rangée « Ce cours contient » (page de garde)
\newcommand{\poptile}[3]{% #1 couleur #2 gros texte #3 légende
  \begin{tcolorbox}[enhanced,colback=#1,colframe=popink,boxrule=2pt,arc=9pt,shadow={2.5pt}{-2.5pt}{0pt}{popink},
      left=6pt,right=6pt,top=9pt,bottom=9pt,halign=center,valign=center,height=1.75cm,width=\linewidth,nobeforeafter]
    {\popdisplay\fontsize{12.5}{13}\selectfont\color{white}\MakeUppercase{#2}}\par\vspace{3pt}
    {\centering\popsemi\fontsize{7.8}{9.5}\selectfont\color{white}#3\par}
  \end{tcolorbox}}
\newcommand{\popcontenu}{%
  \par\noindent{\popxboldls\fontsize{7.5}{7.5}\selectfont\color{popmuted}CE COURS SIGNÉ CONTIENT}\par\vspace{7pt}
  \noindent\begin{minipage}[t]{0.235\linewidth}\poptile{popblue}{Cours}{complet, illustré, pas à pas}\end{minipage}\hfill
  \begin{minipage}[t]{0.235\linewidth}\poptile{poporange}{Méthodes}{et exercices résolus}\end{minipage}\hfill
  \begin{minipage}[t]{0.235\linewidth}\poptile{poppink}{Exercices}{type examen + corrigés}\end{minipage}\hfill
  \begin{minipage}[t]{0.235\linewidth}\poptile{popgreen}{Fiche bilan}{l'essentiel à retenir}\end{minipage}\par}

% Sommaire
\newtcolorbox{sommaire}{enhanced,colback=white,colframe=popink,boxrule=2.2pt,arc=10pt,
  drop shadow=popink,left=18pt,right=18pt,top=13pt,bottom=13pt,before skip=6pt,after skip=20pt,
  before upper={\poppill{Sommaire}{poppurple}{white}\par\vspace{9pt}\popsemi\fontsize{10.6}{14.5}\selectfont\color{popink}}}

% Signature de fin de cours — poussée tout en bas de la dernière page
\newcommand{\findecours}{%
  \par\vfill\nopagebreak
  \begin{center}\popsquiggle{poporange}\end{center}\vspace{4pt}
  \signatureonbuch
}
\AtBeginDocument{\def\llbracket{\mathopen{[\mkern-3.5mu[}}\def\rrbracket{\mathclose{]\mkern-3.5mu]}}} % intervalles d'entiers (glyphe ⟦ absent de la police)
% Glyphes absents de Fira Math : redéfinis à partir de glyphes disponibles
\AtBeginDocument{%
  \def\setminus{\mathbin{\backslash}}\let\smallsetminus\setminus
  \def\vdots{\mathord{\vbox{\baselineskip3.2pt\lineskiplimit0pt\kern4pt\hbox{.}\hbox{.}\hbox{.}}}}%
  \def\ddots{\mathinner{\mkern1mu\raise6.5pt\hbox{.}\mkern2mu\raise3.6pt\hbox{.}\mkern2mu\raise0.7pt\hbox{.}\mkern1mu}}%
  \def\triangle{\mathord{\scalebox{0.9}{$\symup{\Delta}$}}}%
  \def\nabla{\mathord{\raisebox{\depth}{\scalebox{1}[-1]{$\symup{\Delta}$}}}}%
  \def\checkmark{\popcheck{popdark}}%
  \def\star{\mathbin{\ast}}\def\bigstar{\mathord{\ast}}%
  \def\diamond{\mathbin{\scalebox{0.6}{$\Diamond$}}}%
  \def\bigcup{\mathop{\mathchoice{\raisebox{-0.3em}{\scalebox{1.7}{$\cup$}}}{\scalebox{1.25}{$\cup$}}{\cup}{\cup}}\displaylimits}%
  \def\bigcap{\mathop{\mathchoice{\raisebox{-0.3em}{\scalebox{1.7}{$\cap$}}}{\scalebox{1.25}{$\cap$}}{\cap}{\cap}}\displaylimits}%
  \def\lesseqqgtr{\mathrel{\lessgtr}}\def\bigoplus{\mathop{\oplus}\displaylimits}%
}
\providecommand{\ssi}{si et seulement si} % abréviation fréquente
\AtBeginDocument{\providecommand{\wideparen}{\overparen}} % arc de cercle

\begin{document}

\pagegarde{Trigonométrie}{Mathématiques}{1re Tech}{Mathématiques – classe de Première technique}{N10}{1 séance (fiche de progression harmonisée)}{Cette leçon complète l'étude de la trigonométrie par les formules de transformation (sommes en produits, produits en sommes) et la résolution d'équations et d'inéquations trigonométriques. Elle fournit des outils indispensables pour la modélisation de phénomènes périodiques et pour la classe de Terminale.
\vspace{4pt}
\begin{multicols}{2}\small
\begin{itemize}
\item À la fin de cette leçon, je sais rappeler et utiliser les formules d'addition et de duplication
\item transformer une somme de sinus ou de cosinus en produit (et réciproquement)
\item transformer un produit de fonctions trigonométriques en somme
\item résoudre les équations de type cos x = a, sin x = a et tan x = a sur un intervalle donné
\item résoudre des équations se ramenant aux formes précédentes par factorisation ou transformation
\item résoudre des inéquations trigonométriques simples à l'aide du cercle trigonométrique
\end{itemize}\end{multicols}}

\begin{sommaire}
\begin{multicols}{2}
\begin{enumerate}
\item Rappels : formules d'addition et de duplication
\item Transformer des produits en sommes
\item Transformer des sommes en produits
\item Équations trigonométriques élémentaires
\item Équations se ramenant aux formes élémentaires
\item Inéquations trigonométriques
\item Activité d'intégration
\item Méthodes et exercices résolus
\item Exercices
\item Corrigés des exercices
\item Fiche bilan
\end{enumerate}
\end{multicols}
\end{sommaire}

\begin{objectifs}
\begin{itemize}
\item À la fin de cette leçon, je sais rappeler et utiliser les formules d'addition et de duplication
\item transformer une somme de sinus ou de cosinus en produit (et réciproquement)
\item transformer un produit de fonctions trigonométriques en somme
\item résoudre les équations de type cos x = a, sin x = a et tan x = a sur un intervalle donné
\item résoudre des équations se ramenant aux formes précédentes par factorisation ou transformation
\item résoudre des inéquations trigonométriques simples à l'aide du cercle trigonométrique
\item représenter les solutions d'une équation ou d'une inéquation sur le cercle trigonométrique
\item appliquer ces notions à des situations concrètes de la vie courante (phénomènes périodiques)
\item rédiger et justifier une démarche de résolution complète et rigoureuse
\end{itemize}
\end{objectifs}

\begin{prerequis}
\begin{itemize}
\item cercle trigonométrique, mesures d'angles en radians et en degrés
\item valeurs remarquables de cos, sin et tan (0, π/6, π/4, π/3, π/2)
\item formules d'addition : cos(a±b), sin(a±b) vues en début d'année
\item résolution d'équations et d'inéquations du premier et du second degré
\item notion de périodicité des fonctions sinus et cosinus
\end{itemize}
\end{prerequis}

\newpage

\coursec{Rappels : formules d'addition et de duplication}

Un technicien d'Eneo doit régler deux alternateurs dont les tensions sinusoïdales se superposent : pour prévoir la tension résultante, il doit transformer une somme de cosinus en produit. De même, un technicien en topographie à Yaoundé cherche à quels instants de la journée l'angle d'élévation du soleil dépasse une valeur donnée : il doit résoudre une inéquation trigonométrique. Comment additionner, transformer et résoudre des expressions trigonométriques ? C'est l'objet de cette leçon.

Avant de transformer et de résoudre, il faut disposer d'une boîte à outils : les \cle{formules d'addition} et les \cle{formules de duplication}. C'est ce que nous rappelons ici, en les démontrant pour les plus importantes.

\courssub{Pourquoi des formules d'addition ?}

Sur le cercle trigonométrique, on connaît par cœur les valeurs de $\cos$ et $\sin$ pour les angles remarquables : $0$, $\dfrac{\pi}{6}$ ($30^{\circ}$), $\dfrac{\pi}{4}$ ($45^{\circ}$), $\dfrac{\pi}{3}$ ($60^{\circ}$), $\dfrac{\pi}{2}$ ($90^{\circ}$). Mais comment calculer $\cos 75^{\circ}$ ou $\sin 15^{\circ}$ ? L'idée est simple : $75^{\circ} = 45^{\circ} + 30^{\circ}$ et $15^{\circ} = 45^{\circ} - 30^{\circ}$. Si l'on sait exprimer $\cos(a+b)$ et $\sin(a-b)$ en fonction de $\cos a$, $\cos b$, $\sin a$, $\sin b$, le tour est joué.

\begin{attention}[Le piège fondamental]
Pour des nombres réels $a$ et $b$, on a en général :
\[ \cos(a+b) \neq \cos a + \cos b \qquad \text{et} \qquad \sin(a+b) \neq \sin a + \sin b. \]
Par exemple : $\cos\left(\dfrac{\pi}{2} + \dfrac{\pi}{2}\right) = \cos \pi = -1$, alors que $\cos\dfrac{\pi}{2} + \cos\dfrac{\pi}{2} = 0 + 0 = 0$. Le cosinus n'est \textbf{pas} une fonction additive : il faut obligatoirement passer par les formules d'addition.
\end{attention}

\courssub{Formules d'addition pour le cosinus et le sinus}

\begin{propriete}[Formules d'addition]
Pour tous réels $a$ et $b$ :
\begin{align*}
\cos(a+b) &= \cos a \cos b - \sin a \sin b \\
\cos(a-b) &= \cos a \cos b + \sin a \sin b \\
\sin(a+b) &= \sin a \cos b + \cos a \sin b \\
\sin(a-b) &= \sin a \cos b - \cos a \sin b
\end{align*}
La démonstration de $\cos(a-b)$ est guidée ci-dessous ; les trois autres formules s'en déduisent et sont admises comme conséquences.
\end{propriete}

\begin{experience}[Démonstration guidée de $\cos(a-b)$ par le produit scalaire]
Plaçons-nous dans un repère orthonormé $(O\,;\vec{i},\vec{j})$ et considérons le cercle trigonométrique (cercle de centre $O$ et de rayon $1$).

\textbf{Étape 1.} Soit $A$ le point du cercle associé à l'angle $a$ et $B$ le point associé à l'angle $b$. Par définition du cosinus et du sinus :
\[ \overrightarrow{OA} = \cos a \, \vec{i} + \sin a \, \vec{j} \qquad \text{et} \qquad \overrightarrow{OB} = \cos b \, \vec{i} + \sin b \, \vec{j}. \]
Ces deux vecteurs sont \cle{unitaires} : $\|\overrightarrow{OA}\| = \|\overrightarrow{OB}\| = 1$.

\textbf{Étape 2.} L'angle géométrique entre $\overrightarrow{OA}$ et $\overrightarrow{OB}$ a pour mesure (à $2\pi$ près) $a - b$. Or le produit scalaire de deux vecteurs s'écrit :
\[ \overrightarrow{OA} \cdot \overrightarrow{OB} = \|\overrightarrow{OA}\| \times \|\overrightarrow{OB}\| \times \cos(a-b) = \cos(a-b). \]

\textbf{Étape 3.} Calculons le même produit scalaire avec les coordonnées :
\[ \overrightarrow{OA} \cdot \overrightarrow{OB} = \cos a \cos b + \sin a \sin b. \]

\textbf{Étape 4.} En égalisant les deux expressions :
\[ \boxed{\cos(a-b) = \cos a \cos b + \sin a \sin b.} \]

\textbf{Conséquences (admises).} En remplaçant $b$ par $-b$ et en utilisant la parité ($\cos(-b) = \cos b$, $\sin(-b) = -\sin b$), on obtient $\cos(a+b)$. En écrivant $\sin x = \cos\left(\dfrac{\pi}{2} - x\right)$, on obtient les formules pour $\sin(a+b)$ et $\sin(a-b)$.
\end{experience}

\begin{popfigure}
\begin{center}
\begin{tikzpicture}[scale=2.6]
% axes
\draw[->,popmuted] (-1.35,0) -- (1.45,0) node[below right] {$\cos$};
\draw[->,popmuted] (0,-1.25) -- (0,1.35) node[above left] {$\sin$};
% cercle
\draw[popline,thick] (0,0) circle (1);
% points et vecteurs
\coordinate (O) at (0,0);
\coordinate (A) at (50:1);
\coordinate (B) at (20:1);
\draw[->,popblue,very thick] (O) -- (A) node[above right] {$A(\cos a\,;\sin a)$};
\draw[->,poporange,very thick] (O) -- (B) node[right] {$B(\cos b\,;\sin b)$};
% angles
\draw[poporange,->] (0.55,0) arc (0:20:0.55) node[midway,right] {\small $b$};
\draw[popblue,->] (0.4,0) arc (0:50:0.4) node[midway,above right=-2pt] {\small $a$};
\draw[popgreen,->] (20:0.75) arc (20:50:0.75) node[midway,above right=-2pt] {\small $a-b$};
\fill[popdark] (O) circle (0.015) node[below left] {$O$};
\fill[popblue] (A) circle (0.015);
\fill[poporange] (B) circle (0.015);
\end{tikzpicture}
\end{center}
\legende{Cercle trigonométrique : les vecteurs unitaires $\overrightarrow{OA}$ et $\overrightarrow{OB}$, d'angles $a$ et $b$. Leur produit scalaire donne $\cos(a-b)$.}
\end{popfigure}

\begin{aretenir}[Moyen mnémotechnique]
Dans $\cos(a \pm b)$, on mélange cosinus et sinus, et le \textbf{signe est inversé} : $\cos(a+b)$ porte un $-$, $\cos(a-b)$ porte un $+$. Dans $\sin(a \pm b)$, le \textbf{signe est conservé}.
\end{aretenir}

\courssub{Formules d'addition pour la tangente}

\begin{propriete}[Tangente d'une somme et d'une différence]
Pour tous réels $a$ et $b$ tels que $\tan a$, $\tan b$ et le dénominateur existent :
\[ \tan(a+b) = \dfrac{\tan a + \tan b}{1 - \tan a \tan b} \qquad \text{et} \qquad \tan(a-b) = \dfrac{\tan a - \tan b}{1 + \tan a \tan b}. \]
\textbf{Conditions d'existence :} il faut $a \neq \dfrac{\pi}{2} + k\pi$, $b \neq \dfrac{\pi}{2} + k\pi$ ($k \in \mathbb{Z}$), et pour $\tan(a+b)$ : $a+b \neq \dfrac{\pi}{2} + k\pi$, c'est-à-dire $\tan a \tan b \neq 1$.
\end{propriete}

Ces formules s'obtiennent en écrivant $\tan(a+b) = \dfrac{\sin(a+b)}{\cos(a+b)}$, puis en divisant numérateur et dénominateur par $\cos a \cos b$.

\courssub{Formules de duplication}

En posant $b = a$ dans les formules d'addition, on obtient les formules dites de \cle{duplication} (angle double).

\begin{propriete}[Formules de duplication]
Pour tout réel $a$ :
\[ \sin 2a = 2 \sin a \cos a \]
\[ \cos 2a = \cos^2 a - \sin^2 a = 2\cos^2 a - 1 = 1 - 2\sin^2 a \]
et, lorsque $\tan a$ et $\tan 2a$ existent :
\[ \tan 2a = \dfrac{2\tan a}{1 - \tan^2 a}. \]
Les trois écritures de $\cos 2a$ se déduisent les unes des autres grâce à l'identité $\cos^2 a + \sin^2 a = 1$.
\end{propriete}

\begin{exemplebox}[Vérification sur un angle connu]
Prenons $a = \dfrac{\pi}{6}$, donc $2a = \dfrac{\pi}{3}$. On sait que $\sin\dfrac{\pi}{6} = \dfrac{1}{2}$ et $\cos\dfrac{\pi}{6} = \dfrac{\sqrt{3}}{2}$.
\[ \sin 2a = 2 \times \dfrac{1}{2} \times \dfrac{\sqrt{3}}{2} = \dfrac{\sqrt{3}}{2} = \sin \dfrac{\pi}{3}. \quad \checkmark \]
\[ \cos 2a = 2\cos^2 a - 1 = 2 \times \dfrac{3}{4} - 1 = \dfrac{1}{2} = \cos \dfrac{\pi}{3}. \quad \checkmark \]
\end{exemplebox}

\courssub{Formules de linéarisation}

En isolant $\cos^2 a$ et $\sin^2 a$ dans les formules de duplication, on obtient des formules très utiles pour « abaisser le degré » d'une expression, notamment en électricité (puissance d'un signal sinusoïdal) et pour les calculs d'intégrales en Terminale.

\begin{propriete}[Formules de linéarisation]
Pour tout réel $a$ :
\[ \cos^2 a = \dfrac{1 + \cos 2a}{2} \qquad \text{et} \qquad \sin^2 a = \dfrac{1 - \cos 2a}{2}. \]
\end{propriete}

\begin{exemplebox}[Application]
Calculons $\cos^2 \dfrac{\pi}{8}$. Avec $a = \dfrac{\pi}{8}$, on a $2a = \dfrac{\pi}{4}$ :
\[ \cos^2 \dfrac{\pi}{8} = \dfrac{1 + \cos \frac{\pi}{4}}{2} = \dfrac{1 + \frac{\sqrt{2}}{2}}{2} = \dfrac{2 + \sqrt{2}}{4}. \]
Comme $\cos \dfrac{\pi}{8} > 0$, on en déduit $\cos \dfrac{\pi}{8} = \dfrac{\sqrt{2 + \sqrt{2}}}{2}$.
\end{exemplebox}

\courssub{Méthode : calculer une valeur non remarquable}

\begin{methode}[Décomposer l'angle en somme ou différence d'angles remarquables]
Pour calculer une valeur exacte comme $\cos 75^{\circ}$ ou $\sin 15^{\circ}$ :
\begin{enumerate}
\item Écrire l'angle comme somme ou différence de deux angles remarquables : $75^{\circ} = 45^{\circ} + 30^{\circ}$ ; $15^{\circ} = 45^{\circ} - 30^{\circ}$ ou $60^{\circ} - 45^{\circ}$.
\item Appliquer la formule d'addition correspondante.
\item Remplacer par les valeurs remarquables connues.
\item Simplifier l'expression obtenue (factoriser, réduire au même dénominateur).
\end{enumerate}
\end{methode}

\begin{center}
\begin{tabularx}{\linewidth}{|c|X|X|X|X|X|}
\hline
\rowcolor{popblueL} $x$ & $0$ & $\dfrac{\pi}{6}$ ($30^{\circ}$) & $\dfrac{\pi}{4}$ ($45^{\circ}$) & $\dfrac{\pi}{3}$ ($60^{\circ}$) & $\dfrac{\pi}{2}$ ($90^{\circ}$) \\
\hline
$\cos x$ & $1$ & $\dfrac{\sqrt{3}}{2}$ & $\dfrac{\sqrt{2}}{2}$ & $\dfrac{1}{2}$ & $0$ \\
\hline
$\sin x$ & $0$ & $\dfrac{1}{2}$ & $\dfrac{\sqrt{2}}{2}$ & $\dfrac{\sqrt{3}}{2}$ & $1$ \\
\hline
$\tan x$ & $0$ & $\dfrac{\sqrt{3}}{3}$ & $1$ & $\sqrt{3}$ & indéfini \\
\hline
\end{tabularx}
\end{center}

\begin{exoresolu}[Calcul exact de $\cos 75^{\circ}$]
Énoncé. Calculer la valeur exacte de $\cos 75^{\circ}$, puis en déduire $\sin 15^{\circ}$.
\tcblower
Solution détaillée. On écrit $75^{\circ} = 45^{\circ} + 30^{\circ}$ et on applique la formule d'addition :
\begin{align*}
\cos 75^{\circ} &= \cos(45^{\circ} + 30^{\circ}) \\
&= \cos 45^{\circ} \cos 30^{\circ} - \sin 45^{\circ} \sin 30^{\circ} \\
&= \dfrac{\sqrt{2}}{2} \times \dfrac{\sqrt{3}}{2} - \dfrac{\sqrt{2}}{2} \times \dfrac{1}{2} \\
&= \dfrac{\sqrt{6}}{4} - \dfrac{\sqrt{2}}{4} = \dfrac{\sqrt{6} - \sqrt{2}}{4}.
\end{align*}
Pour $\sin 15^{\circ}$, on remarque que $15^{\circ} = 90^{\circ} - 75^{\circ}$, donc $\sin 15^{\circ} = \cos 75^{\circ} = \dfrac{\sqrt{6} - \sqrt{2}}{4}$. Vérification numérique : $\dfrac{\sqrt{6}-\sqrt{2}}{4} \approx 0{,}259$ et $\sin 15^{\circ} \approx 0{,}259$. $\checkmark$
\end{exoresolu}

\begin{exoresolu}[Calcul exact de $\sin 75^{\circ}$]
Énoncé. Calculer la valeur exacte de $\sin 75^{\circ}$.
\tcblower
Solution détaillée. On applique la formule $\sin(a+b) = \sin a \cos b + \cos a \sin b$ :
\begin{align*}
\sin 75^{\circ} &= \sin 45^{\circ} \cos 30^{\circ} + \cos 45^{\circ} \sin 30^{\circ} \\
&= \dfrac{\sqrt{2}}{2} \times \dfrac{\sqrt{3}}{2} + \dfrac{\sqrt{2}}{2} \times \dfrac{1}{2} \\
&= \dfrac{\sqrt{6} + \sqrt{2}}{4}.
\end{align*}
On remarque que $\sin 75^{\circ} = \cos 15^{\circ}$, ce qui est cohérent avec les angles complémentaires.
\end{exoresolu}

\begin{attention}[Erreurs fréquentes]
\begin{itemize}
\item \textbf{Inverser les signes :} $\cos(a+b) = \cos a \cos b \textcolor{popdark}{-} \sin a \sin b$ (signe inversé), mais $\sin(a+b) = \sin a \cos b \textcolor{popdark}{+} \cos a \sin b$ (signe conservé).
\item \textbf{Oublier les conditions d'existence} pour la tangente : $\tan a$ n'existe pas si $a = \dfrac{\pi}{2} + k\pi$.
\item \textbf{Confondre $\cos^2 a$ et $\cos 2a$} : $\cos^2 a = (\cos a)^2$ alors que $\cos 2a = \cos(2a)$. Ce sont deux objets différents, reliés par $\cos 2a = 2\cos^2 a - 1$.
\item \textbf{Écrire $\sin 2a = 2 \sin a$} : c'est faux ! La formule correcte est $\sin 2a = 2 \sin a \cos a$.
\end{itemize}
\end{attention}

\begin{savaistu}[Trigonométrie et électricité au Cameroun]
Les tensions alternatives distribuées par Eneo sont des signaux sinusoïdaux de fréquence \SI{50}{Hz}. Lorsque deux sources se superposent, la tension résultante fait intervenir des expressions du type $\cos a \cos b$ ou $\sin a + \sin b$ : les formules de cette leçon permettent de prévoir l'amplitude et le déphasage du signal résultant. C'est aussi grâce à ces formules que les topographes calculent des distances inaccessibles par triangulation, une technique utilisée pour cartographier le territoire camerounais.
\end{savaistu}

\begin{exercice}[À toi de jouer]
\begin{enumerate}
\item Calculer la valeur exacte de $\cos 15^{\circ}$ en écrivant $15^{\circ} = 60^{\circ} - 45^{\circ}$.
\item Calculer la valeur exacte de $\tan 75^{\circ}$ à l'aide de la formule d'addition des tangentes.
\item Montrer que $\cos^2 15^{\circ} = \dfrac{2 + \sqrt{3}}{4}$.
\end{enumerate}
\end{exercice}

\begin{corrige}[Exercice]
\begin{enumerate}
\item $\cos 15^{\circ} = \cos 60^{\circ}\cos 45^{\circ} + \sin 60^{\circ}\sin 45^{\circ} = \dfrac{1}{2}\times\dfrac{\sqrt{2}}{2} + \dfrac{\sqrt{3}}{2}\times\dfrac{\sqrt{2}}{2} = \dfrac{\sqrt{2}+\sqrt{6}}{4}$.
\item $\tan 75^{\circ} = \dfrac{\tan 45^{\circ} + \tan 30^{\circ}}{1 - \tan 45^{\circ}\tan 30^{\circ}} = \dfrac{1 + \frac{\sqrt{3}}{3}}{1 - \frac{\sqrt{3}}{3}} = \dfrac{3 + \sqrt{3}}{3 - \sqrt{3}} = 2 + \sqrt{3}$ après multiplication par le conjugué.
\item $\cos^2 15^{\circ} = \dfrac{1 + \cos 30^{\circ}}{2} = \dfrac{1 + \frac{\sqrt{3}}{2}}{2} = \dfrac{2 + \sqrt{3}}{4}$.
\end{enumerate}
\end{corrige}

\coursec{Transformer des produits en sommes}

Dans la section précédente, nous avons rappelé les formules d'addition et de duplication qui permettent de \emph{décomposer} le cosinus ou le sinus d'une somme ou d'un double angle. Nous allons maintenant faire le chemin inverse : à partir de ces mêmes formules, nous allons apprendre à \emph{recomposer} un produit de fonctions trigonométriques en une somme (ou une différence). Cette technique, appelée \textbf{linéarisation}, est un outil puissant pour simplifier des expressions, résoudre des équations et, plus tard en Terminale, calculer des primitives d'expressions trigonométriques.

\courssub{Obtention des formules : addition et soustraction des formules d'addition}

Rappelons les deux formules d'addition pour le cosinus, qui sont notre point de départ :
\[
\cos(a+b) = \cos a \cos b - \sin a \sin b \qquad \text{et} \qquad \cos(a-b) = \cos a \cos b + \sin a \sin b.
\]
Si nous les additionnons membre à membre, les termes en $\sin a \sin b$ s'annulent. Si nous les soustrayons, ce sont les termes en $\cos a \cos b$ qui disparaissent. C'est l'idée centrale de la démonstration guidée ci-dessous.

\begin{experience}[Démonstration guidée : du produit à la somme]
Partons des deux égalités fondamentales :
\begin{align*}
(1) \quad & \cos(a+b) = \cos a \cos b - \sin a \sin b \\
(2) \quad & \cos(a-b) = \cos a \cos b + \sin a \sin b
\end{align*}

\textbf{Étape 1 : Additionner (1) et (2).}
\emph{Que deviennent les termes en $\sin a \sin b$ ? Écrivez le résultat et isolez $\cos a \cos b$.}

\textbf{Étape 2 : Soustraire (1) de (2) (c'est-à-dire faire (2) $-$ (1)).}
\emph{Que deviennent les termes en $\cos a \cos b$ ? Écrivez le résultat et isolez $\sin a \sin b$.}

\textbf{Étape 3 : Utiliser les formules du sinus.}
Rappelons $\sin(a+b) = \sin a \cos b + \cos a \sin b$ et $\sin(a-b) = \sin a \cos b - \cos a \sin b$.
\emph{Additionnez ces deux égalités pour isoler $\sin a \cos b$.}

\begin{center}
\begin{tikzpicture}[node distance=0.8cm and 1.2cm, every node/.style={font=\small}]
    % Formules de base
    \node (c1) {$\cos(a+b) = \cos a\cos b - \sin a\sin b$};
    \node (c2) [below=of c1] {$\cos(a-b) = \cos a\cos b + \sin a\sin b$};
    
    % Flèche addition
    \draw[popblue, thick, ->] ($(c1.east)+(0.5,0)$) -- ++(1,0) node[midway, above] {Addition} coordinate (mid1);
    \draw[popblue, thick, ->] ($(c2.east)+(0.5,0)$) -- ++(1,0) coordinate (mid2);
    
    % Résultat addition
    \node (res1) [right=of mid1, anchor=west, popblue] {$\cos(a+b)+\cos(a-b) = 2\cos a\cos b$};
    
    % Flèche division
    \draw[poporange, thick, ->] ($(res1.east)+(0.5,0)$) -- ++(1,0) node[midway, above] {$\div 2$};
    \node (final1) [right=of res1, anchor=west, poporange] {$\cos a\cos b = \frac{1}{2}[\cos(a+b)+\cos(a-b)]$};

    % Deuxième bloc : Soustraction
    \node (c3) [below=2.5cm of c2] {$\cos(a-b) = \cos a\cos b + \sin a\sin b$};
    \node (c4) [below=of c3] {$\cos(a+b) = \cos a\cos b - \sin a\sin b$};
    
    \draw[popgreen, thick, ->] ($(c3.east)+(0.5,0)$) -- ++(1,0) node[midway, above] {Soustraction} coordinate (mid3);
    \draw[popgreen, thick, ->] ($(c4.east)+(0.5,0)$) -- ++(1,0) coordinate (mid4);
    
    \node (res2) [right=of mid3, anchor=west, popgreen] {$\cos(a-b)-\cos(a+b) = 2\sin a\sin b$};
    \draw[poppurple, thick, ->] ($(res2.east)+(0.5,0)$) -- ++(1,0) node[midway, above] {$\div 2$};
    \node (final2) [right=of res2, anchor=west, poppurple] {$\sin a\sin b = \frac{1}{2}[\cos(a-b)-\cos(a+b)]$};

    % Troisième bloc : Sinus
    \node (s1) [below=2.5cm of c4] {$\sin(a+b) = \sin a\cos b + \cos a\sin b$};
    \node (s2) [below=of s1] {$\sin(a-b) = \sin a\cos b - \cos a\sin b$};
    
    \draw[poppink, thick, ->] ($(s1.east)+(0.5,0)$) -- ++(1,0) node[midway, above] {Addition} coordinate (mid5);
    \draw[poppink, thick, ->] ($(s2.east)+(0.5,0)$) -- ++(1,0) coordinate (mid6);
    
    \node (res3) [right=of mid5, anchor=west, poppink] {$\sin(a+b)+\sin(a-b) = 2\sin a\cos b$};
    \draw[popgold, thick, ->] ($(res3.east)+(0.5,0)$) -- ++(1,0) node[midway, above] {$\div 2$};
    \node (final3) [right=of res3, anchor=west, popgold] {$\sin a\cos b = \frac{1}{2}[\sin(a+b)+\sin(a-b)]$};
\end{tikzpicture}
\legende{Construction visuelle des formules produit $\to$ somme par combinaisons linéaires des formules d'addition.}
\end{center}
\end{experience}

\begin{propriete}[Formules de transformation Produit $\to$ Somme]
Pour tous réels $a$ et $b$, on a les trois formules fondamentales :
\begin{align*}
\cos a \cos b &= \frac{1}{2}\bigl[\cos(a+b) + \cos(a-b)\bigr] \tag{P1}\\[4pt]
\sin a \sin b &= \frac{1}{2}\bigl[\cos(a-b) - \cos(a+b)\bigr] \tag{P2}\\[4pt]
\sin a \cos b &= \frac{1}{2}\bigl[\sin(a+b) + \sin(a-b)\bigr] \tag{P3}
\end{align*}
\emph{Remarque :} La formule (P2) s'obtient par soustraction des formules d'addition du cosinus (étape 2 de l'activité guidée). La formule (P3) est symétrique : $\cos b \sin a$ donne le même résultat.
\end{propriete}

\begin{attention}[Pièges classiques : signes et ordre]
\begin{itemize}
    \item \textbf{Formule (P2) :} C'est \textbf{$\cos(a-b) - \cos(a+b)$} (le cosinus de la différence \emph{moins} le cosinus de la somme). L'ordre compte car le cosinus est pair : $\cos(a-b)=\cos(b-a)$, mais l'écriture standard privilégie $a-b$.
    \item \textbf{Facteur $\frac{1}{2}$ :} Ne l'oubliez jamais ! C'est l'erreur n°1 aux examens. Vérifiez : si $a=b=0$, $\cos 0 \cos 0 = 1$ et $\frac{1}{2}[\cos 0 + \cos 0] = \frac{1}{2}[1+1]=1$. Sans le $\frac{1}{2}$, on trouve $2$.
    \item \textbf{Arguments :} Dans les parenthèses, ce sont toujours \textbf{la somme} et \textbf{la différence} des angles initiaux ($a+b$ et $a-b$).
\end{itemize}
\end{attention}

\courssub{Applications immédiates : linéarisation des carrés et du produit $\sin x \cos x$}

Un cas particulier extrêmement fréquent (et au programme du Probatoire) est celui où les deux angles sont égaux : $a = b = x$. Les formules deviennent des formules de \textbf{linéarisation} (on « linéarise » les puissances 2).

\begin{propriete}[Formules de linéarisation (cas $a=b=x$)]
\begin{align*}
\cos^2 x &= \frac{1 + \cos 2x}{2} \\
\sin^2 x &= \frac{1 - \cos 2x}{2} \\
\sin x \cos x &= \frac{\sin 2x}{2}
\end{align*}
\emph{Démonstration :} Dans (P1), posez $a=b=x$ : $\cos^2 x = \frac{1}{2}[\cos 2x + \cos 0] = \frac{1+\cos 2x}{2}$.
Dans (P2), posez $a=b=x$ : $\sin^2 x = \frac{1}{2}[\cos 0 - \cos 2x] = \frac{1-\cos 2x}{2}$.
Dans (P3), posez $a=b=x$ : $\sin x \cos x = \frac{1}{2}[\sin 2x + \sin 0] = \frac{\sin 2x}{2}$.
\end{propriete}

\begin{methode}[Linéariser une expression trigonométrique]
Pour transformer un produit ou une puissance en somme :
\begin{enumerate}
    \item \textbf{Identifiez} la structure : $\cos^2$, $\sin^2$, $\sin\cos$, $\cos\cos$, $\sin\sin$ avec des angles identiques ou différents.
    \item \textbf{Choisissez} la formule adaptée (P1), (P2), (P3) ou les formules de linéarisation ci-dessus.
    \item \textbf{Substituez} les angles ($a$ et $b$ ou $x$) soigneusement.
    \item \textbf{Simplifiez} les angles obtenus ($a+b$, $a-b$, $2x$) et réduisez les fractions.
\end{enumerate}
\end{methode}

\begin{exemplebox}[Linéariser $\cos^2 3x$ et $\sin 5x \cos 2x$]
\begin{enumerate}
    \item $\cos^2 3x$ : On utilise $\cos^2 X = \frac{1+\cos 2X}{2}$ avec $X = 3x$.
    \[
    \cos^2 3x = \frac{1 + \cos(2\times 3x)}{2} = \frac{1 + \cos 6x}{2}.
    \]
    \item $\sin 5x \cos 2x$ : On utilise (P3) $\sin a \cos b = \frac{1}{2}[\sin(a+b) + \sin(a-b)]$ avec $a=5x$, $b=2x$.
    \[
    \sin 5x \cos 2x = \frac{1}{2}\bigl[\sin(5x+2x) + \sin(5x-2x)\bigr] = \frac{1}{2}[\sin 7x + \sin 3x].
    \]
\end{enumerate}
\end{exemplebox}

\begin{exoresolu}[Écrire $\sin 3x \cos 2x$ comme somme de sinus]
\textbf{Énoncé :} Linéariser l'expression $E = \sin 3x \cos 2x$.

\textbf{Solution détaillée :}
\begin{enumerate}
    \item \textbf{Identification :} Produit d'un sinus et d'un cosinus $\rightarrow$ Formule (P3) : $\sin a \cos b = \frac{1}{2}[\sin(a+b) + \sin(a-b)]$.
    \item \textbf{Repérage des angles :} $a = 3x$ et $b = 2x$.
    \item \textbf{Application :}
    \[
    E = \frac{1}{2}\bigl[\sin(3x+2x) + \sin(3x-2x)\bigr]
    \]
    \item \textbf{Calcul des sommes/différences :}
    \[
    3x+2x = 5x \quad \text{et} \quad 3x-2x = x.
    \]
    \item \textbf{Résultat final :}
    \[
    \boxed{\sin 3x \cos 2x = \frac{1}{2}\sin 5x + \frac{1}{2}\sin x}
    \]
\end{enumerate}
\textbf{Vérification rapide (optionnelle) :} Pour $x = \frac{\pi}{2}$, $\sin\frac{3\pi}{2}\cos\pi = (-1)(-1)=1$. La formule donne $\frac{1}{2}\sin\frac{5\pi}{2} + \frac{1}{2}\sin\frac{\pi}{2} = \frac{1}{2}(1) + \frac{1}{2}(1) = 1$. C'est cohérent.
\end{exoresolu}

\courssub{Autres cas de linéarisation : angles distincts}

Lorsque les angles sont différents ($a \neq b$), on applique directement les formules (P1), (P2), (P3). L'astuce consiste parfois à \emph{reconnaître} un produit caché.

\begin{exemplebox}[Produits cosinus-cosinus et sinus-sinus]
\begin{enumerate}
    \item \textbf{$\cos 40^\circ \cos 20^\circ$} (angles en degrés, typique en géométrie/arpentage).
    Formule (P1) : $\cos a \cos b = \frac{1}{2}[\cos(a+b) + \cos(a-b)]$ avec $a=40^\circ, b=20^\circ$.
    \[
    \cos 40^\circ \cos 20^\circ = \frac{1}{2}[\cos 60^\circ + \cos 20^\circ] = \frac{1}{2}\left[\frac{1}{2} + \cos 20^\circ\right] = \frac{1}{4} + \frac{1}{2}\cos 20^\circ.
    \]
    \item \textbf{$\sin 7x \sin 3x$}.
    Formule (P2) : $\sin a \sin b = \frac{1}{2}[\cos(a-b) - \cos(a+b)]$ avec $a=7x, b=3x$.
    \[
    \sin 7x \sin 3x = \frac{1}{2}[\cos(4x) - \cos(10x)].
    \]
    \emph{Attention :} Le résultat ne contient que des \textbf{cosinus} (la formule (P2) fait « sauter » le sinus).
\end{enumerate}
\end{exemplebox}

\begin{savaistu}[Télécommunications : la modulation d'amplitude (AM)]
Ces formules ne sont pas de simples exercices scolaires ! Elles sont au cœur du fonctionnement de la \textbf{radio AM} (Amplitude Modulation) et de la téléphonie mobile.
Un signal audio (voix, musique) est une onde de basse fréquence $s(t) = A \sin(\omega_s t)$. Pour la transporter par antenne, on la « module » sur une porteuse de haute fréquence $c(t) = \sin(\omega_c t)$ ($\omega_c \gg \omega_s$). Le signal émis est le produit :
\[
e(t) = s(t) \cdot c(t) = A \sin(\omega_s t) \sin(\omega_c t).
\]
Grâce à la formule (P2) : $\sin a \sin b = \frac{1}{2}[\cos(a-b) - \cos(a+b)]$, ce produit devient une \textbf{somme} de deux ondes pures :
\[
e(t) = \frac{A}{2}\cos\bigl((\omega_c-\omega_s)t\bigr) - \frac{A}{2}\cos\bigl((\omega_c+\omega_s)t\bigr).
\]
Le spectre du signal émis contient donc deux fréquences : $\omega_c - \omega_s$ (bande latérale inférieure) et $\omega_c + \omega_s$ (bande latérale supérieure). Le récepteur radio utilise l'opération inverse (détection d'enveloppe ou mélangeur) pour retrouver $\omega_s$. \textbf{Linéariser, c'est passer du domaine temporel (produit) au domaine fréquentiel (somme) — la base du traitement du signal numérique !}
\end{savaistu}

\begin{attention}[Erreurs fréquentes en situation d'examen]
\begin{itemize}
    \item \textbf{Confondre (P1) et (P2) :} Retenir que « Cosinus $\times$ Cosinus $\to$ Somme de Cosinus » (P1) et « Sinus $\times$ Sinus $\to$ Différence de Cosinus » (P2). Mnémotechnique : \emph{« Même fonctions $\to$ Cosinus ; Différentes fonctions $\to$ Sinus »} (pour P3).
    \item \textbf{Oublier le facteur $\frac{1}{2}$ :} Systématiquement vérifier par une valeur simple ($x=0$ ou $x=\pi/2$).
    \item \textbf{Mauvaise gestion du signe moins dans (P2) :} Écrire $\cos(a+b) - \cos(a-b)$ au lieu de $\cos(a-b) - \cos(a+b)$.
    \item \textbf{Ne pas simplifier l'angle final :} Laisser $\sin(3x-2x)$ au lieu de $\sin x$.
\end{itemize}
\end{attention}

\begin{exercice}[Entraînement : Linéarisez les expressions suivantes]
\begin{enumerate}
    \item $\cos 5x \cos 2x$
    \item $\sin 4x \sin x$
    \item $\cos^2 2x$
    \item $\sin 3x \cos 5x$
    \item $\sin^2 \frac{\pi}{8}$ (donner une valeur exacte)
    \item $2 \sin 10^\circ \cos 10^\circ$ (simplifier au maximum)
\end{enumerate}
\end{exercice}

\begin{corrige}[Exercice n°1]
\begin{enumerate}
    \item $\cos 5x \cos 2x = \frac{1}{2}[\cos 7x + \cos 3x]$ \quad (P1)
    \item $\sin 4x \sin x = \frac{1}{2}[\cos 3x - \cos 5x]$ \quad (P2 : $\cos(4x-x) - \cos(4x+x)$)
    \item $\cos^2 2x = \frac{1+\cos 4x}{2}$ \quad (Linéarisation cos²)
    \item $\sin 3x \cos 5x = \frac{1}{2}[\sin 8x + \sin(-2x)] = \frac{1}{2}\sin 8x - \frac{1}{2}\sin 2x$ \quad (P3 + parité du sinus)
    \item $\sin^2 \frac{\pi}{8} = \frac{1-\cos\frac{\pi}{4}}{2} = \frac{1-\frac{\sqrt{2}}{2}}{2} = \frac{2-\sqrt{2}}{4}$.
    \item $2 \sin 10^\circ \cos 10^\circ = \sin 20^\circ$ \quad (Formule de duplication $\sin 2x = 2\sin x \cos x$, ou P3 multipliée par 2).
\end{enumerate}
\end{corrige}

\coursec{Transformer des sommes en produits}

Dans la section précédente, nous avons vu comment \emph{transformer des produits en sommes} pour simplifier des intégrales, calculer des primitives ou linéariser des expressions. Nous abordons maintenant l'opération inverse : \emph{transformer des sommes en produits}. Cette technique, appelée \textbf{factorisation trigonométrique}, est l'outil principal pour résoudre les équations de la forme $\sin p \pm \sin q = 0$ ou $\cos p \pm \cos q = 0$, et pour étudier le signe d'expressions trigonométriques. Elle permet de passer d'une \emph{somme} (difficile à annuler) à un \emph{produit} (facile à annuler grâce à la règle du signe : un produit est nul si et seulement si l'un de ses facteurs est nul).

\courssub{Formules de transformation : sommes $\to$ produits}

\courssub{Énoncé des quatre formules fondamentales}

\begin{propriete}[Formules de transformation sommes en produits]
Pour tous réels $p$ et $q$ :
\begin{align*}
\cos p + \cos q &= 2 \cos\left(\frac{p+q}{2}\right) \cos\left(\frac{p-q}{2}\right) \tag{1}\\
\cos p - \cos q &= -2 \sin\left(\frac{p+q}{2}\right) \sin\left(\frac{p-q}{2}\right) \tag{2}\\
\sin p + \sin q &= 2 \sin\left(\frac{p+q}{2}\right) \cos\left(\frac{p-q}{2}\right) \tag{3}\\
\sin p - \sin q &= 2 \cos\left(\frac{p+q}{2}\right) \sin\left(\frac{p-q}{2}\right) \tag{4}
\end{align*}
\end{propriete}

\begin{aretenir}[Mémorisation efficace : la règle des « demi-somme / demi-différence »]
Toutes les formules ont la structure : \textbf{2 $\times$ f(demi-somme) $\times$ g(demi-différence)}.
\begin{itemize}
    \item \textbf{Somme de cosinus} $\to$ produit de \textbf{cosinus} : $2 \cos\frac{p+q}{2} \cos\frac{p-q}{2}$.
    \item \textbf{Différence de cosinus} $\to$ produit de \textbf{sinus} (avec un signe $-$ devant) : $-2 \sin\frac{p+q}{2} \sin\frac{p-q}{2}$.
    \item \textbf{Somme de sinus} $\to$ \textbf{sinus} de la demi-somme $\times$ \textbf{cosinus} de la demi-différence : $2 \sin\frac{p+q}{2} \cos\frac{p-q}{2}$.
    \item \textbf{Différence de sinus} $\to$ \textbf{cosinus} de la demi-somme $\times$ \textbf{sinus} de la demi-différence : $2 \cos\frac{p+q}{2} \sin\frac{p-q}{2}$.
\end{itemize}
\emph{Astuce :} Pour les deux dernières, l'ordre des fonctions (sin puis cos, ou cos puis sin) suit l'ordre alphabétique \textbf{cos} < \textbf{sin} pour la différence, et l'ordre naturel \textbf{sin} puis \textbf{cos} pour la somme.
\end{aretenir}

\courssub{Démonstration guidée par changement de variables}

La démonstration est exigible au Probatoire. Elle repose sur les formules d'addition et le changement de variables $a = \frac{p+q}{2}$, $b = \frac{p-q}{2}$.

\begin{experience}[Démonstration des formules sommes $\to$ produits]
On part des formules d'addition (rappelées section précédente) :
\begin{align*}
\cos(a+b) &= \cos a \cos b - \sin a \sin b \\
\cos(a-b) &= \cos a \cos b + \sin a \sin b \\
\sin(a+b) &= \sin a \cos b + \cos a \sin b \\
\sin(a-b) &= \sin a \cos b - \cos a \sin b
\end{align*}

\noindent\textbf{Étape 1 : Changement de variables.}
On pose :
\[
a = \frac{p+q}{2} \quad \text{(demi-somme)} \qquad b = \frac{p-q}{2} \quad \text{(demi-différence)}
\]
Alors :
\[
p = a+b \quad \text{et} \quad q = a-b
\]

\noindent\textbf{Étape 2 : Somme de cosinus (Formule 1).}
On additionne $\cos(a+b)$ et $\cos(a-b)$ :
\[
\cos(a+b) + \cos(a-b) = (\cos a \cos b - \sin a \sin b) + (\cos a \cos b + \sin a \sin b) = 2 \cos a \cos b
\]
En revenant à $p, q$ : $\cos p + \cos q = 2 \cos\frac{p+q}{2} \cos\frac{p-q}{2}$. \quad \textcolor{popgreen}{\textbf{OK}}

\noindent\textbf{Étape 3 : Différence de cosinus (Formule 2).}
On soustrait $\cos(a-b)$ de $\cos(a+b)$ :
\[
\cos(a+b) - \cos(a-b) = (\cos a \cos b - \sin a \sin b) - (\cos a \cos b + \sin a \sin b) = -2 \sin a \sin b
\]
Donc $\cos p - \cos q = -2 \sin\frac{p+q}{2} \sin\frac{p-q}{2}$. \quad \textcolor{popgreen}{\textbf{OK}}

\noindent\textbf{Étape 4 : Somme de sinus (Formule 3).}
On additionne $\sin(a+b)$ et $\sin(a-b)$ :
\[
\sin(a+b) + \sin(a-b) = (\sin a \cos b + \cos a \sin b) + (\sin a \cos b - \cos a \sin b) = 2 \sin a \cos b
\]
Donc $\sin p + \sin q = 2 \sin\frac{p+q}{2} \cos\frac{p-q}{2}$. \quad \textcolor{popgreen}{\textbf{OK}}

\noindent\textbf{Étape 5 : Différence de sinus (Formule 4).}
On soustrait $\sin(a-b)$ de $\sin(a+b)$ :
\[
\sin(a+b) - \sin(a-b) = (\sin a \cos b + \cos a \sin b) - (\sin a \cos b - \cos a \sin b) = 2 \cos a \sin b
\]
Donc $\sin p - \sin q = 2 \cos\frac{p+q}{2} \sin\frac{p-q}{2}$. \quad \textcolor{popgreen}{\textbf{OK}}
\end{experience}

\begin{popfigure}
\begin{tikzpicture}[scale=1.2, >=Stealth]
    % Cercle trigonométrique
    \draw[popline, thick] (0,0) circle (2.5cm);
    \draw[popmuted, ->] (-3,0) -- (3,0) node[right] {$x$};
    \draw[popmuted, ->] (0,-3) -- (0,3) node[above] {$y$};
    \node[below left] at (0,0) {$O$};
    
    % Angles p et q
    \def\p{60}
    \def\q{20}
    \def\rp{2.5}
    
    % Point P (angle p)
    \coordinate (P) at (\p:\rp);
    \draw[popblue, thick] (0,0) -- (P);
    \draw[popblue, thick] (0.5,0) arc (0:\p:0.5) node[midway, right=2pt] {$p$};
    \node[popblue, above right] at (P) {$P(\cos p, \sin p)$};
    
    % Point Q (angle q)
    \coordinate (Q) at (\q:\rp);
    \draw[poporange, thick] (0,0) -- (Q);
    \draw[poporange, thick] (0.3,0) arc (0:\q:0.3) node[midway, right=2pt] {$q$};
    \node[poporange, above right] at (Q) {$Q(\cos q, \sin q)$};
    
    % Demi-somme (p+q)/2
    \pgfmathsetmacro{\dsum}{(\p+\q)/2}
    \coordinate (M) at (\dsum:\rp);
    \draw[popgreen, thick, dashed] (0,0) -- (M);
    \draw[popgreen, thick] (0.7,0) arc (0:\dsum:0.7) node[midway, above left=2pt] {$\frac{p+q}{2}$};
    \node[popgreen, above left] at (M) {$M$};
    
    % Demi-différence (p-q)/2
    \pgfmathsetmacro{\ddiff}{(\p-\q)/2}
    \coordinate (N) at (\ddiff:\rp);
    \draw[poppurple, thick, dashed] (0,0) -- (N);
    \draw[poppurple, thick] (0.9,0) arc (0:\ddiff:0.9) node[midway, below right=2pt] {$\frac{p-q}{2}$};
    \node[poppurple, below right] at (N) {$N$};
    
    % Légende
    \begin{scope}[shift={(3.5,1.5)}]
        \draw[popblue, thick] (0,0.3) -- (0.5,0.3) node[right] {$p$};
        \draw[poporange, thick] (0,0) -- (0.5,0) node[right] {$q$};
        \draw[popgreen, thick, dashed] (0,-0.3) -- (0.5,-0.3) node[right] {$\frac{p+q}{2}$ (demi-somme)};
        \draw[poppurple, thick, dashed] (0,-0.6) -- (0.5,-0.6) node[right] {$\frac{p-q}{2}$ (demi-différence)};
    \end{scope}
\end{tikzpicture}
\legende{Représentation géométrique des angles $p$, $q$, de leur demi-somme $\frac{p+q}{2}$ et de leur demi-différence $\frac{p-q}{2}$ sur le cercle trigonométrique.}
\end{popfigure}

\courssub{Applications à la factorisation et à la résolution d'équations}

La factorisation transforme une équation somme $=0$ en un produit $=0$, ce qui permet d'appliquer la règle : $A \times B = 0 \iff A=0 \text{ ou } B=0$.

\begin{methode}[Résoudre une équation par factorisation somme $\to$ produit]
\begin{enumerate}
    \item Identifier la structure : $\sin p \pm \sin q = 0$ ou $\cos p \pm \cos q = 0$.
    \item Appliquer la formule adaptée pour factoriser l'expression.
    \item Mettre chaque facteur égal à zéro.
    \item Résoudre les équations élémentaires obtenues ($\cos X = 0$, $\sin X = 0$, $\cos X = \cos \alpha$, $\sin X = \sin \alpha$).
    \item Rassembler les solutions (souvent sur un intervalle donné comme $[0, 2\pi[$ ou $\mathbb{R}$).
\end{enumerate}
\end{methode}

\begin{exoresolu}[{Factoriser et résoudre $\sin 5x + \sin x = 0$ sur $[0, 2\pi[$}]
\textbf{Énoncé :} Factoriser l'expression $E = \sin 5x + \sin x$, puis résoudre l'équation $E = 0$ sur l'intervalle $[0, 2\pi[$.

\textbf{Solution détaillée :}

\noindent\textbf{1. Factorisation.}
On reconnaît une somme de sinus : $\sin p + \sin q$ avec $p = 5x$ et $q = x$.
On applique la formule (3) : $\sin p + \sin q = 2 \sin\frac{p+q}{2} \cos\frac{p-q}{2}$.
Calcul des demi-somme et demi-différence :
\[
\frac{p+q}{2} = \frac{5x+x}{2} = 3x \qquad \frac{p-q}{2} = \frac{5x-x}{2} = 2x
\]
Donc :
\[
E = \sin 5x + \sin x = 2 \sin 3x \cos 2x
\]

\noindent\textbf{2. Résolution de l'équation $E=0$.}
L'équation devient :
\[
2 \sin 3x \cos 2x = 0 \iff \sin 3x = 0 \quad \text{ou} \quad \cos 2x = 0
\]

\noindent\textbf{Cas 1 : $\sin 3x = 0$.}
$\sin \theta = 0 \iff \theta = k\pi,\ k \in \mathbb{Z}$.
Ici $\theta = 3x$, donc $3x = k\pi \iff x = \frac{k\pi}{3}$.
On cherche $x \in [0, 2\pi[$ :
$k=0 \to x=0$ ; $k=1 \to x=\frac{\pi}{3}$ ; $k=2 \to x=\frac{2\pi}{3}$ ; $k=3 \to x=\pi$ ;
$k=4 \to x=\frac{4\pi}{3}$ ; $k=5 \to x=\frac{5\pi}{3}$ ; $k=6 \to x=2\pi$ (exclu car intervalle ouvert à droite).
Solutions : $\left\{0, \frac{\pi}{3}, \frac{2\pi}{3}, \pi, \frac{4\pi}{3}, \frac{5\pi}{3}\right\}$.

\noindent\textbf{Cas 2 : $\cos 2x = 0$.}
$\cos \theta = 0 \iff \theta = \frac{\pi}{2} + k\pi,\ k \in \mathbb{Z}$.
Ici $\theta = 2x$, donc $2x = \frac{\pi}{2} + k\pi \iff x = \frac{\pi}{4} + k\frac{\pi}{2}$.
On cherche $x \in [0, 2\pi[$ :
$k=0 \to x=\frac{\pi}{4}$ ; $k=1 \to x=\frac{3\pi}{4}$ ; $k=2 \to x=\frac{5\pi}{4}$ ; $k=3 \to x=\frac{7\pi}{4}$ ; $k=4 \to x=\frac{9\pi}{4} > 2\pi$.
Solutions : $\left\{\frac{\pi}{4}, \frac{3\pi}{4}, \frac{5\pi}{4}, \frac{7\pi}{4}\right\}$.

\noindent\textbf{Conclusion :}
L'ensemble des solutions sur $[0, 2\pi[$ est :
\[
S = \left\{0, \frac{\pi}{4}, \frac{\pi}{3}, \frac{3\pi}{4}, \frac{2\pi}{3}, \pi, \frac{4\pi}{3}, \frac{5\pi}{4}, \frac{5\pi}{3}, \frac{7\pi}{4}\right\}
\]
\end{exoresolu}

\begin{popfigure}
\begin{tikzpicture}
\begin{axis}[
    width=\linewidth, height=7cm,
    axis lines=middle,
    xlabel=$x$, ylabel=$y$,
    domain=0:6.28, samples=500,
    xmin=0, xmax=6.5, ymin=-2.5, ymax=2.5,
    xtick={0,1.57,3.14,4.71,6.28},
    xticklabels={$0$, $\frac{\pi}{2}$, $\pi$, $\frac{3\pi}{2}$, $2\pi$},
    ytick={-2,-1,1,2},
    grid=major, grid style={popline!30},
    legend style={at={(0.98,0.98)}, anchor=north east, fill=white, fill opacity=0.8},
]
% Courbe y = sin(5x) + sin(x) -- x en radians, pgfplots sin attend degrés -> deg()
\addplot[popblue, thick, domain=0:6.28, samples=500] {sin(deg(5*x)) + sin(deg(x))};
\addlegendentry{$y = \sin 5x + \sin x$}

% Courbe factorisée : 2 sin(3x) cos(2x)
\addplot[poporange, thick, dashed, domain=0:6.28, samples=500] {2*sin(deg(3*x))*cos(deg(2*x))};
\addlegendentry{$y = 2\sin 3x \cos 2x$ (factorisée)}

% Enveloppes
\addplot[popgreen, thick, domain=0:6.28, samples=200] {2*abs(cos(deg(2*x)))};
\addplot[popgreen, thick, domain=0:6.28, samples=200] {-2*abs(cos(deg(2*x)))};
\addlegendentry{Enveloppes $\pm 2|\cos 2x|$}

% Zéros
\addplot[only marks, mark=*, mark size=2, popred] coordinates {
    (0,0) (0.785,0) (1.047,0) (2.094,0) (2.356,0) (3.142,0) 
    (3.927,0) (4.189,0) (5.236,0) (5.498,0)
};
\addlegendentry{Zéros (solutions)}
\end{axis}
\end{tikzpicture}
\legende{Graphiques de $y = \sin 5x + \sin x$ (bleu) et de sa forme factorisée $y = 2\sin 3x \cos 2x$ (orange pointillé, superposée). Les zéros correspondent aux intersections avec l'axe des abscisses. Les courbes vertes représentent l'enveloppe $\pm 2|\cos 2x|$ qui module l'amplitude de l'oscillation rapide $\sin 3x$.}
\end{popfigure}

\courssub{Application à l'étude du signe d'une expression trigonométrique}

La forme factorisée permet de construire un \textbf{tableau de signes} rigoureux, exactement comme pour les fonctions polynomiales ou rationnelles.

\begin{exemplebox}[{Étude du signe de $f(x) = \cos 3x - \cos x$ sur $[0, \pi]$}]
\textbf{1. Factorisation.}
Différence de cosinus : $\cos p - \cos q = -2 \sin\frac{p+q}{2} \sin\frac{p-q}{2}$.
$p=3x, q=x \Rightarrow \frac{p+q}{2}=2x, \frac{p-q}{2}=x$.
\[
f(x) = -2 \sin 2x \sin x
\]

\textbf{2. Zéros des facteurs sur $[0, \pi]$.}
$\sin x = 0 \iff x = 0, \pi$.
$\sin 2x = 0 \iff 2x = k\pi \iff x = k\frac{\pi}{2}$. Sur $[0, \pi] : x = 0, \frac{\pi}{2}, \pi$.

\textbf{3. Tableau de signes.}
\begin{center}
\begin{tabular}{|c|ccccc|}\hline
$x$ & $0$ & & $\frac{\pi}{2}$ & & $\pi$ \\ \hline
$\sin x$ & $0$ & $+$ & $+$ & $+$ & $0$ \\
$\sin 2x$ & $0$ & $+$ & $0$ & $-$ & $0$ \\
$-2$ & $-$ & $-$ & $-$ & $-$ & $-$ \\ \hline
$f(x)$ & $0$ & $-$ & $0$ & $+$ & $0$ \\ \hline
\end{tabular}
\end{center}

\textbf{4. Conclusion.}
$f(x) \le 0$ sur $\left[0, \frac{\pi}{2}\right]$ et $f(x) \ge 0$ sur $\left[\frac{\pi}{2}, \pi\right]$.
\end{exemplebox}

\courssub{Erreurs fréquentes et pièges à éviter}

\begin{attention}[Pièges classiques en transformation sommes $\to$ produits]
\begin{enumerate}
    \item \textbf{Oublier le facteur 2.} \emph{Faux :} $\cos p + \cos q = \cos\frac{p+q}{2}\cos\frac{p-q}{2}$. \textbf{Vrai :} $2\cos\frac{p+q}{2}\cos\frac{p-q}{2}$.
    \item \textbf{Inverser demi-somme et demi-différence.} Vérifiez toujours : $\frac{p+q}{2}$ est la moyenne (au milieu), $\frac{p-q}{2}$ est l'écart. Si $p > q$, la demi-différence est positive.
    \item \textbf{Oublier le signe $-$ pour $\cos p - \cos q$.} C'est la \emph{seule} formule qui porte un signe $-$ devant le 2. Mnémotechnique : « Cosinus \emph{diminue} sur $[0,\pi]$, donc la différence donne un signe moins ».
    \item \textbf{Diviser l'angle par 2 au lieu de faire la demi-somme.} Pour $\sin 5x + \sin x$, on ne fait pas $\sin(5x/2)\dots$ mais $\sin(\frac{5x+x}{2}) = \sin 3x$.
    \item \textbf{Ne pas factoriser complètement avant un tableau de signes.} L'expression $\sin 3x \cos 2x$ est factorisée ; $\sin 5x + \sin x$ ne l'est pas. Le tableau de signes ne se fait \emph{que} sur la forme factorisée.
\end{enumerate}
\end{attention}

\courssub{Exercices d'application}

\begin{exercice}[Entraînement : Factorisation et résolution]
\begin{enumerate}
    \item Factoriser les expressions suivantes :
    \begin{itemize}
        \item[a)] $\cos 7x + \cos 3x$
        \item[b)] $\sin 4x - \sin 2x$
        \item[c)] $\cos x - \cos 5x$
    \end{itemize}
    \item Résoudre sur $[0, 2\pi[$ les équations :
    \begin{itemize}
        \item[a)] $\sin 3x + \sin x = 0$
        \item[b)] $\cos 4x - \cos 2x = 0$
        \item[c)] $\sin 5x - \sin x = \cos 3x$ \emph{(Indication : factoriser le membre de gauche, puis utiliser $\cos 3x = \sin(\frac{\pi}{2} - 3x)$ ou factoriser le membre de droite après passage à gauche).}
    \end{itemize}
    \item Étudier le signe de $g(x) = \sin 4x + \sin 2x$ sur l'intervalle $[0, \pi]$.
\end{enumerate}
\end{exercice}

\begin{corrige}[Exercice n°1]
\begin{enumerate}
    \item \textbf{Factorisations :}
    \begin{itemize}
        \item[a)] $\cos 7x + \cos 3x = 2 \cos\left(\frac{7x+3x}{2}\right) \cos\left(\frac{7x-3x}{2}\right) = \boxed{2 \cos 5x \cos 2x}$.
        \item[b)] $\sin 4x - \sin 2x = 2 \cos\left(\frac{4x+2x}{2}\right) \sin\left(\frac{4x-2x}{2}\right) = \boxed{2 \cos 3x \sin x}$.
        \item[c)] $\cos x - \cos 5x = -2 \sin\left(\frac{x+5x}{2}\right) \sin\left(\frac{x-5x}{2}\right) = -2 \sin 3x \sin(-2x) = \boxed{2 \sin 3x \sin 2x}$ (car $\sin(-2x)=-\sin 2x$).
    \end{itemize}
    \item \textbf{Résolutions sur $[0, 2\pi[$ :}
    \begin{itemize}
        \item[a)] $\sin 3x + \sin x = 2 \sin 2x \cos x = 0 \iff \sin 2x=0 \text{ ou } \cos x=0$.
        $\sin 2x=0 \Rightarrow 2x=k\pi \Rightarrow x=k\frac{\pi}{2} \in \{0, \frac{\pi}{2}, \pi, \frac{3\pi}{2}\}$.
        $\cos x=0 \Rightarrow x=\frac{\pi}{2}+k\pi \in \{\frac{\pi}{2}, \frac{3\pi}{2}\}$.
        Union : $\boxed{S = \{0, \frac{\pi}{2}, \pi, \frac{3\pi}{2}\}}$.
        \item[b)] $\cos 4x - \cos 2x = -2 \sin 3x \sin x = 0 \iff \sin 3x=0 \text{ ou } \sin x=0$.
        $\sin x=0 \Rightarrow x=k\pi \in \{0, \pi\}$.
        $\sin 3x=0 \Rightarrow 3x=k\pi \Rightarrow x=k\frac{\pi}{3} \in \{0, \frac{\pi}{3}, \frac{2\pi}{3}, \pi, \frac{4\pi}{3}, \frac{5\pi}{3}\}$.
        Union : $\boxed{S = \{0, \frac{\pi}{3}, \frac{2\pi}{3}, \pi, \frac{4\pi}{3}, \frac{5\pi}{3}\}}$.
        \item[c)] $\sin 5x - \sin x - \cos 3x = 0 \iff 2 \cos 3x \sin 2x - \cos 3x = 0 \iff \cos 3x (2 \sin 2x - 1) = 0$.
        \begin{itemize}
            \item $\cos 3x = 0 \Rightarrow 3x = \frac{\pi}{2}+k\pi \Rightarrow x = \frac{\pi}{6}+k\frac{\pi}{3}$.
            Sur $[0, 2\pi[$ : $k=0 \to \frac{\pi}{6}$ ; $k=1 \to \frac{\pi}{2}$ ; $k=2 \to \frac{5\pi}{6}$ ; $k=3 \to \frac{7\pi}{6}$ ; $k=4 \to \frac{3\pi}{2}$ ; $k=5 \to \frac{11\pi}{6}$.
            \item $2\sin 2x - 1 = 0 \iff \sin 2x = \frac{1}{2} \iff 2x = \frac{\pi}{6}+2k\pi \text{ ou } \frac{5\pi}{6}+2k\pi \iff x = \frac{\pi}{12}+k\pi \text{ ou } \frac{5\pi}{12}+k\pi$.
            Sur $[0, 2\pi[$ : $\frac{\pi}{12}, \frac{13\pi}{12}, \frac{5\pi}{12}, \frac{17\pi}{12}$.
        \end{itemize}
        Ensemble des solutions (trié) :
        \[
        \boxed{S = \left\{\frac{\pi}{12}, \frac{\pi}{6}, \frac{5\pi}{12}, \frac{\pi}{2}, \frac{5\pi}{6}, \frac{13\pi}{12}, \frac{7\pi}{6}, \frac{17\pi}{12}, \frac{3\pi}{2}, \frac{11\pi}{6}\right\}}
        \]
    \end{itemize}
    \item \textbf{Signe de $g(x) = \sin 4x + \sin 2x = 2 \sin 3x \cos x$ sur $[0, \pi]$.}
    Zéros : $\sin 3x=0 \Rightarrow x \in \{0, \frac{\pi}{3}, \frac{2\pi}{3}, \pi\}$ ; $\cos x=0 \Rightarrow x=\frac{\pi}{2}$.
    \begin{center}
    \begin{tabular}{|c|ccccccccc|}\hline
    $x$ & $0$ & & $\frac{\pi}{3}$ & & $\frac{\pi}{2}$ & & $\frac{2\pi}{3}$ & & $\pi$ \\ \hline
    $\sin 3x$ & $0$ & $+$ & $0$ & $-$ & $-$ & $-$ & $0$ & $+$ & $0$ \\
    $\cos x$ & $+$ & $+$ & $+$ & $+$ & $0$ & $-$ & $-$ & $-$ & $-$ \\
    $2$ & $+$ & $+$ & $+$ & $+$ & $+$ & $+$ & $+$ & $+$ & $+$ \\ \hline
    $g(x)$ & $0$ & $+$ & $0$ & $-$ & $0$ & $+$ & $0$ & $-$ & $0$ \\ \hline
    \end{tabular}
    \end{center}
    $g(x) \ge 0$ sur $\left[0, \frac{\pi}{3}\right] \cup \left[\frac{\pi}{2}, \frac{2\pi}{3}\right]$ et $g(x) \le 0$ sur $\left[\frac{\pi}{3}, \frac{\pi}{2}\right] \cup \left[\frac{2\pi}{3}, \pi\right]$.
\end{enumerate}
\end{corrige}

\begin{savaistu}[Applications physiques : Battements acoustiques]
La formule $\sin p + \sin q = 2 \sin\frac{p+q}{2} \cos\frac{p-q}{2}$ modélise le phénomène des \textbf{battements} en acoustique. Si deux diapasons émettent des sons de fréquences proches $f_1$ et $f_2$ (pulsations $\omega_1=2\pi f_1$, $\omega_2=2\pi f_2$), l'onde résultante est :
\[
y(t) = \sin(\omega_1 t) + \sin(\omega_2 t) = 2 \underbrace{\sin\left(\frac{\omega_1+\omega_2}{2}t\right)}_{\text{son "moyen" aigu}} \underbrace{\cos\left(\frac{\omega_1-\omega_2}{2}t\right)}_{\text{enveloppe lente}}
\]
L'oreille perçoit une note de fréquence moyenne $\frac{f_1+f_2}{2}$ dont l'intensité varie (battements) à la fréquence $\frac{|f_1-f_2|}{2}$. C'est la technique utilisée par les accordeurs de piano : on écoute la disparition des battements pour accorder l'unisson.
\end{savaistu}

\coursec{Équations trigonométriques élémentaires}

Dans la section précédente, nous avons appris à transformer des sommes en produits : une expression comme $\sin x + \sin 3x$ devient un produit $2\sin 2x \cos x$. Pourquoi un tel effort ? Parce qu'un produit est nul dès que l'un de ses facteurs est nul : transformer une somme en produit permet de ramener une équation compliquée à des équations très simples, du type $\cos x = a$, $\sin x = a$ ou $\tan x = a$. Ce sont ces équations, dites \cle{élémentaires}, que nous allons maintenant résoudre complètement. Elles interviennent partout dans la vie technique : réglage d'un angle de coupe en atelier, positionnement d'un bras de robot, calcul de phases en électricité (les tensions alternatives sont des sinusoïdes !).

\courssub{L'équation $\cos x = a$}

\textbf{Intuition.} Sur le cercle trigonométrique, $\cos x$ est l'\emph{abscisse} du point image de $x$. Chercher les $x$ tels que $\cos x = a$, c'est donc chercher les points du cercle dont l'abscisse vaut $a$. Géométriquement, on trace la droite verticale d'équation $X = a$ et on regarde où elle coupe le cercle. Si $|a| > 1$, la droite ne touche pas le cercle : pas de solution. Si $|a| \le 1$, elle le coupe en \emph{deux} points symétriques par rapport à l'axe des abscisses, correspondant à deux angles opposés $\alpha$ et $-\alpha$.

\begin{propriete}[Équation $\cos x = a$]
Soit $a$ un réel.
\begin{itemize}
\item Si $|a| > 1$, l'équation $\cos x = a$ n'a \textbf{aucune solution} dans $\mathbb{R}$.
\item Si $|a| \le 1$, il existe un unique réel $\alpha \in [0\,;\pi]$ tel que $\cos \alpha = a$, et l'ensemble des solutions est :
\[ x = \alpha + 2k\pi \quad \text{ou} \quad x = -\alpha + 2k\pi, \qquad k \in \mathbb{Z}. \]
\end{itemize}
Ce résultat est admis ; il se justifie par la lecture sur le cercle trigonométrique et par la $2\pi$-périodicité du cosinus.
\end{propriete}

\begin{popfigure}
\begin{center}
\begin{tikzpicture}[scale=2.4]
\draw[popline,->] (-1.4,0) -- (1.5,0) node[below right] {$\cos$};
\draw[popline,->] (0,-1.3) -- (0,1.4) node[above left] {$\sin$};
\draw[popblue,thick] (0,0) circle (1);
\coordinate (A) at (0.5,0.866);
\coordinate (B) at (0.5,-0.866);
\draw[poporange,thick] (0.5,-1.15) -- (0.5,1.15) node[above,popdark] {$X = a$};
\draw[popgreen,thick,->] (0,0) -- (A);
\draw[popgreen,thick,->] (0,0) -- (B);
\fill[popink] (A) circle (0.03) node[above right,popdark] {$\alpha$};
\fill[popink] (B) circle (0.03) node[below right,popdark] {$-\alpha$};
\fill[poporange] (0.5,0) circle (0.03) node[below left=1pt,popdark] {$a$};
\draw[popgold,thick,->] (0.35,0) arc (0:60:0.35) node[midway,right,popdark] {\small $\alpha$};
\draw[popgold,thick,->] (0.35,0) arc (0:-60:0.35);
\end{tikzpicture}
\end{center}
\legende{La droite $X = a$ coupe le cercle en deux points : les angles $\alpha$ et $-\alpha$ (modulo $2\pi$).}
\end{popfigure}

\begin{attention}[Ne pas oublier la deuxième famille]
L'erreur la plus fréquente au Probatoire est d'écrire seulement $x = \alpha + 2k\pi$. Pour le cosinus, il y a TOUJOURS deux familles : $x = \alpha + 2k\pi$ \textbf{et} $x = -\alpha + 2k\pi$. Pensez aux deux points d'intersection sur le cercle !
\end{attention}

\courssub{L'équation $\sin x = a$}

\textbf{Intuition.} Cette fois, $\sin x$ est l'\emph{ordonnée} du point image. On trace donc la droite \emph{horizontale} $Y = a$. Si $|a| \le 1$, elle coupe le cercle en deux points symétriques par rapport à l'axe des ordonnées : leurs angles sont $\alpha$ et $\pi - \alpha$.

\begin{propriete}[Équation $\sin x = a$]
Soit $a$ un réel.
\begin{itemize}
\item Si $|a| > 1$, l'équation $\sin x = a$ n'a \textbf{aucune solution}.
\item Si $|a| \le 1$, il existe un unique réel $\alpha \in \left[-\dfrac{\pi}{2}\,;\dfrac{\pi}{2}\right]$ tel que $\sin \alpha = a$, et l'ensemble des solutions est :
\[ x = \alpha + 2k\pi \quad \text{ou} \quad x = \pi - \alpha + 2k\pi, \qquad k \in \mathbb{Z}. \]
\end{itemize}
Résultat admis, justifié par le cercle trigonométrique et la périodicité du sinus.
\end{propriete}

\begin{popfigure}
\begin{center}
\begin{tikzpicture}[scale=2.4]
\draw[popline,->] (-1.4,0) -- (1.5,0) node[below right] {$\cos$};
\draw[popline,->] (0,-1.3) -- (0,1.4) node[above left] {$\sin$};
\draw[popblue,thick] (0,0) circle (1);
\coordinate (A) at (0.866,0.5);
\coordinate (B) at (-0.866,0.5);
\draw[poporange,thick] (-1.2,0.5) -- (1.25,0.5) node[right,popdark] {$Y = a$};
\draw[popgreen,thick,->] (0,0) -- (A);
\draw[popgreen,thick,->] (0,0) -- (B);
\fill[popink] (A) circle (0.03) node[above right,popdark] {$\alpha$};
\fill[popink] (B) circle (0.03) node[above left,popdark] {$\pi - \alpha$};
\fill[poporange] (0,0.5) circle (0.03) node[below right=1pt,popdark] {$a$};
\draw[popgold,thick,->] (0.4,0) arc (0:30:0.4) node[midway,right,popdark] {\small $\alpha$};
\draw[popgold,thick,->] (0.55,0) arc (0:150:0.55);
\end{tikzpicture}
\end{center}
\legende{La droite $Y = a$ coupe le cercle en deux points : les angles $\alpha$ et $\pi - \alpha$ (modulo $2\pi$).}
\end{popfigure}

\courssub{L'équation $\tan x = a$}

\textbf{Intuition.} La tangente se lit sur la tangente au cercle au point $I(1\,;0)$ : c'est l'ordonnée du point où la droite $(OM)$ coupe cette verticale. Pour toute valeur $a$, la droite passant par l'origine et le point $(1\,;a)$ coupe le cercle en \emph{deux points diamétralement opposés}, correspondant aux angles $\alpha$ et $\alpha + \pi$. Comme la tangente est $\pi$-périodique, ces deux points ne donnent qu'\textbf{une seule} famille de solutions.

\begin{propriete}[Équation $\tan x = a$]
Pour tout réel $a$, l'équation $\tan x = a$ admet des solutions. Si $\alpha$ est l'unique réel de $\left]-\dfrac{\pi}{2}\,;\dfrac{\pi}{2}\right[$ tel que $\tan \alpha = a$, alors :
\[ x = \alpha + k\pi, \qquad k \in \mathbb{Z}. \]
Attention au domaine : $\tan x$ n'existe que si $x \neq \dfrac{\pi}{2} + k\pi$.
\end{propriete}

\begin{popfigure}
\begin{center}
\begin{tikzpicture}[scale=2.2]
\draw[popline,->] (-1.5,0) -- (1.7,0) node[below right] {$\cos$};
\draw[popline,->] (0,-1.5) -- (0,1.7) node[above left] {$\sin$};
\draw[popblue,thick] (0,0) circle (1);
\draw[popline] (1,-1.5) -- (1,1.7) node[above,popdark] {\small axe des tangentes};
\draw[popgreen,thick] (-0.707,-0.707) -- (1.15,1.15);
\fill[popink] (0.707,0.707) circle (0.035) node[above left,popdark] {$\dfrac{\pi}{4}$};
\fill[popink] (-0.707,-0.707) circle (0.035) node[below right,popdark] {$\dfrac{\pi}{4}+\pi$};
\fill[poporange] (1,1) circle (0.035) node[right,popdark] {$1$};
\draw[popgold,thick,->] (0.4,0) arc (0:45:0.4) node[midway,right,popdark] {\small $\frac{\pi}{4}$};
\fill[popink] (1,0) circle (0.03) node[below right,popdark] {$I$};
\end{tikzpicture}
\end{center}
\legende{Solutions de $\tan x = 1$ : les points diamétralement opposés $\frac{\pi}{4}$ et $\frac{\pi}{4} + \pi$ forment une seule famille $x = \frac{\pi}{4} + k\pi$.}
\end{popfigure}

\courssub{Cas particuliers et équations $\cos x = \cos \alpha$, $\sin x = \sin \alpha$}

\begin{aretenir}[Cas particuliers à connaître par cœur]
\begin{itemize}
\item $\cos x = 0 \iff x = \dfrac{\pi}{2} + k\pi$ ;
\item $\sin x = 0 \iff x = k\pi$ ;
\item $\sin x = 1 \iff x = \dfrac{\pi}{2} + 2k\pi$ ; \quad $\sin x = -1 \iff x = -\dfrac{\pi}{2} + 2k\pi$ ;
\item $\cos x = 1 \iff x = 2k\pi$ ; \quad $\cos x = -1 \iff x = \pi + 2k\pi$.
\end{itemize}
Dans ces cas, les deux familles habituelles se confondent en une seule.
\end{aretenir}

\begin{propriete}[Égalité de deux cosinus ou de deux sinus]
Pour tous réels $x$ et $\alpha$ :
\[ \cos x = \cos \alpha \iff x = \alpha + 2k\pi \ \text{ou}\ x = -\alpha + 2k\pi, \quad k \in \mathbb{Z}. \]
\[ \sin x = \sin \alpha \iff x = \alpha + 2k\pi \ \text{ou}\ x = \pi - \alpha + 2k\pi, \quad k \in \mathbb{Z}. \]
\end{propriete}

C'est exactement la même chose que $\cos x = a$ avec $a = \cos \alpha$ : la solution particulière $\alpha$ est donnée directement, sans calculatrice !

\begin{methode}[Résoudre une équation trigonométrique élémentaire]
\begin{enumerate}
\item \textbf{Isoler} la fonction trigonométrique : ramener l'équation à $\cos x = a$, $\sin x = a$ ou $\tan x = a$.
\item \textbf{Vérifier la condition d'existence} : pour cos et sin, il faut $|a| \le 1$.
\item \textbf{Trouver une solution particulière $\alpha$} : valeurs remarquables ($\frac{\pi}{6}, \frac{\pi}{4}, \frac{\pi}{3}, \ldots$) ou calculatrice (touches $\cos^{-1}$, $\sin^{-1}$, $\tan^{-1}$, en mode radians).
\item \textbf{Écrire toutes les solutions} avec $k \in \mathbb{Z}$, sans oublier la deuxième famille.
\item \textbf{Si un intervalle est imposé} (par exemple $[0\,;2\pi]$), choisir les valeurs de $k$ qui donnent des solutions dans cet intervalle.
\end{enumerate}
\end{methode}

\begin{exoresolu}[{Résoudre $2\cos x - 1 = 0$ sur l'intervalle $[0\,;2\pi]$}]
\textbf{Énoncé.} Un technicien règle l'angle $x$ d'une came ; le cahier des charges impose $2\cos x - 1 = 0$ avec $x \in [0\,;2\pi]$. Déterminer tous les angles possibles.
\tcblower
\textbf{Solution détaillée.}
\begin{enumerate}
\item \textbf{Isoler le cosinus :} $2\cos x - 1 = 0 \iff \cos x = \dfrac{1}{2}$.
\item \textbf{Condition :} $\left|\dfrac{1}{2}\right| \le 1$ : il y a des solutions.
\item \textbf{Solution particulière :} valeur remarquable, $\cos \dfrac{\pi}{3} = \dfrac{1}{2}$, donc $\alpha = \dfrac{\pi}{3}$.
\item \textbf{Solutions générales sur $\mathbb{R}$ :}
\[ x = \frac{\pi}{3} + 2k\pi \quad \text{ou} \quad x = -\frac{\pi}{3} + 2k\pi, \quad k \in \mathbb{Z}. \]
\item \textbf{Sélection dans $[0\,;2\pi]$ :}
\begin{itemize}
\item Première famille : $k = 0$ donne $x = \dfrac{\pi}{3}$ ; $k = 1$ donne $\dfrac{7\pi}{3} > 2\pi$, trop grand.
\item Deuxième famille : $k = 0$ donne $-\dfrac{\pi}{3} \notin [0\,;2\pi]$ ; $k = 1$ donne $-\dfrac{\pi}{3} + 2\pi = \dfrac{5\pi}{3}$, qui convient.
\end{itemize}
\end{enumerate}
\textbf{Conclusion :} $\mathcal{S} = \left\{\dfrac{\pi}{3}\,;\ \dfrac{5\pi}{3}\right\}$. Les deux angles de réglage possibles sont $60^\circ$ et $300^\circ$.
\end{exoresolu}

\begin{exoresolu}[{Résoudre $\sin x = \dfrac{\sqrt{2}}{2}$ sur $[-\pi\,;\pi]$}]
\textbf{Énoncé.} Résoudre $\sin x = \dfrac{\sqrt{2}}{2}$ dans $\mathbb{R}$, puis donner les solutions appartenant à $[-\pi\,;\pi]$.
\tcblower
\textbf{Solution détaillée.}
\begin{enumerate}
\item Condition : $\left|\dfrac{\sqrt{2}}{2}\right| \le 1$, OK. Valeur remarquable : $\sin \dfrac{\pi}{4} = \dfrac{\sqrt{2}}{2}$, donc $\alpha = \dfrac{\pi}{4}$.
\item Solutions générales :
\[ x = \frac{\pi}{4} + 2k\pi \quad \text{ou} \quad x = \pi - \frac{\pi}{4} + 2k\pi = \frac{3\pi}{4} + 2k\pi, \quad k \in \mathbb{Z}. \]
\item Dans $[-\pi\,;\pi]$ : avec $k = 0$, on obtient $\dfrac{\pi}{4}$ et $\dfrac{3\pi}{4}$, toutes deux dans l'intervalle. Avec $k = -1$ : $\dfrac{\pi}{4} - 2\pi = -\dfrac{7\pi}{4} < -\pi$, rejeté ; de même $\dfrac{3\pi}{4} - 2\pi = -\dfrac{5\pi}{4} < -\pi$, rejeté.
\end{enumerate}
\textbf{Conclusion :} $\mathcal{S}_{[-\pi;\pi]} = \left\{\dfrac{\pi}{4}\,;\ \dfrac{3\pi}{4}\right\}$.
\end{exoresolu}

\begin{exemplebox}[Équation avec tangente]
Résoudre $\tan x = 1$ sur $[0\,;2\pi]$. \\
On sait que $\tan \dfrac{\pi}{4} = 1$. Les solutions générales sont $x = \dfrac{\pi}{4} + k\pi$, $k \in \mathbb{Z}$. \\
Dans $[0\,;2\pi]$ : $k = 0$ donne $\dfrac{\pi}{4}$ ; $k = 1$ donne $\dfrac{5\pi}{4}$ ; $k = 2$ donne $\dfrac{9\pi}{4} > 2\pi$. \\
$\mathcal{S} = \left\{\dfrac{\pi}{4}\,;\ \dfrac{5\pi}{4}\right\}$ : ce sont les deux points diamétralement opposés de la figure ci-dessus.
\end{exemplebox}

\begin{attention}[Pièges de rédaction au Probatoire]
\begin{itemize}
\item \textbf{Oublier la deuxième famille} : pour $\sin x = a$, la deuxième famille est $x = \pi - \alpha + 2k\pi$ (et non $-\alpha$ !) ; pour $\cos x = a$, c'est $x = -\alpha + 2k\pi$.
\item \textbf{Oublier $k \in \mathbb{Z}$} : écrire « $x = \frac{\pi}{3}$ » sans mentionner $2k\pi$ ne donne qu'une solution parmi une infinité.
\item \textbf{Ne pas vérifier $|a| \le 1$} : l'équation $\cos x = 2$ n'a aucune solution, il faut le dire explicitement.
\item \textbf{Intervalle de recherche} : quand on demande les solutions dans $[0\,;2\pi]$, il faut tester plusieurs valeurs entières de $k$ ($0$, $1$, $-1$) et conclure par un ensemble fini de solutions, pas par une formule avec $k$.
\item \textbf{Mode de la calculatrice} : vérifier qu'elle est en \emph{radians}, sinon $\cos^{-1}(0{,}5)$ affiche $60$ (degrés) au lieu de $\frac{\pi}{3}$.
\end{itemize}
\end{attention}

\begin{exercice}[À faire]
\begin{enumerate}
\item Résoudre dans $\mathbb{R}$ puis dans $[0\,;2\pi]$ : $\cos x = -\dfrac{1}{2}$.
\item Résoudre dans $[-\pi\,;\pi]$ : $\sin x = -\dfrac{\sqrt{3}}{2}$.
\item Résoudre dans $\mathbb{R}$ : $\cos x = \cos \dfrac{\pi}{5}$, puis $\sin 2x = \sin \dfrac{\pi}{3}$.
\end{enumerate}
\end{exercice}

\begin{corrige}[Exercice 1]
\textbf{1.} $\cos x = -\dfrac{1}{2}$ : solution particulière $\alpha = \dfrac{2\pi}{3}$ car $\cos \dfrac{2\pi}{3} = -\dfrac{1}{2}$. Solutions : $x = \dfrac{2\pi}{3} + 2k\pi$ ou $x = -\dfrac{2\pi}{3} + 2k\pi$. Dans $[0\,;2\pi]$ : $\dfrac{2\pi}{3}$ et $-\dfrac{2\pi}{3} + 2\pi = \dfrac{4\pi}{3}$. Donc $\mathcal{S} = \left\{\dfrac{2\pi}{3}\,;\dfrac{4\pi}{3}\right\}$.

\textbf{2.} $\sin x = -\dfrac{\sqrt{3}}{2}$ : $\alpha = -\dfrac{\pi}{3}$. Solutions : $x = -\dfrac{\pi}{3} + 2k\pi$ ou $x = \pi + \dfrac{\pi}{3} + 2k\pi = \dfrac{4\pi}{3} + 2k\pi$. Dans $[-\pi\,;\pi]$ : $-\dfrac{\pi}{3}$ convient ; $\dfrac{4\pi}{3} > \pi$ mais avec $k = -1$ : $\dfrac{4\pi}{3} - 2\pi = -\dfrac{2\pi}{3}$ convient. Donc $\mathcal{S} = \left\{-\dfrac{2\pi}{3}\,;-\dfrac{\pi}{3}\right\}$.

\textbf{3.} $\cos x = \cos \dfrac{\pi}{5} \iff x = \dfrac{\pi}{5} + 2k\pi$ ou $x = -\dfrac{\pi}{5} + 2k\pi$. Pour $\sin 2x = \sin \dfrac{\pi}{3}$ : $2x = \dfrac{\pi}{3} + 2k\pi$ ou $2x = \dfrac{2\pi}{3} + 2k\pi$, soit $x = \dfrac{\pi}{6} + k\pi$ ou $x = \dfrac{\pi}{3} + k\pi$, $k \in \mathbb{Z}$.
\end{corrige}

\begin{savaistu}[Les sinusoïdes de l'électricien]
La tension du réseau électrique camerounais est une fonction $u(t) = 220\sqrt{2}\sin(100\pi t)$. Chercher les instants où la tension vaut $220$ volts revient exactement à résoudre une équation $\sin x = a$ ! C'est pourquoi les électrotechniciens et les automaticiens utilisent quotidiennement les formules de cette leçon, souvent sans même s'en rendre compte, à travers leurs appareils de mesure.
\end{savaistu}

\coursec{Équations se ramenant aux formes élémentaires}

Dans la section précédente, nous avons appris à résoudre les équations \emph{élémentaires} : $\cos x = a$, $\sin x = a$, $\tan x = a$, ainsi que $\cos x = \cos\alpha$ et $\sin x = \sin\alpha$. Mais dans la pratique — en électricité (courants alternatifs), en mécanique (vibrations), en topographie — les équations rencontrées ne sont presque jamais sous cette forme simple. Toute l'habileté consiste alors à \cle{transformer} l'équation donnée pour la \emph{ramener} à une ou plusieurs équations élémentaires que l'on sait résoudre. C'est l'objet de cette section : apprendre à reconnaître les grandes familles d'équations et à choisir la bonne transformation.

\courssub{La démarche générale}

Quelle que soit l'équation, la stratégie est toujours la même et comporte quatre étapes.

\begin{methode}[Démarche générale de résolution]
\begin{enumerate}
\item \textbf{Ramener} l'équation à une ou plusieurs formes élémentaires ($\cos X = a$, $\sin X = a$, $\tan X = a$, $\cos X = \cos\alpha$, etc.) par changement d'inconnue, factorisation ou formule de transformation.
\item \textbf{Résoudre} chaque équation élémentaire obtenue (en $X$ ou en $x$).
\item \textbf{Revenir} à l'inconnue $x$ si un changement d'inconnue a été fait, en ajustant l'intervalle de recherche.
\item \textbf{Sélectionner} les solutions appartenant à l'intervalle demandé et \textbf{vérifier} les conditions d'existence (rejeter les solutions parasites).
\end{enumerate}
\end{methode}

\begin{popfigure}
\begin{tikzpicture}[
node distance=0.55cm,
etape/.style={rectangle, rounded corners, draw=popblue, fill=popblueL, thick, text width=8.2cm, align=center, minimum height=0.9cm, font=\small},
fleche/.style={-{Stealth[length=2.5mm]}, thick, popdark}]
\node[etape] (e1) {\textbf{Étape 1.} Transformer l'équation pour obtenir une forme élémentaire : $\cos X = a$, $\sin X = a$, $\tan X = a$};
\node[etape, below=of e1] (e2) {\textbf{Étape 2.} Résoudre l'équation élémentaire en $X$ : $X = \alpha + 2k\pi$ ou $X = -\alpha + 2k\pi$};
\node[etape, below=of e2] (e3) {\textbf{Étape 3.} Revenir à $x$ : si $X = ax + b$, alors $x = \dfrac{X - b}{a}$};
\node[etape, below=of e3] (e4) {\textbf{Étape 4.} Choisir les entiers $k$ tels que $x$ appartienne à l'intervalle imposé ; vérifier les conditions d'existence};
\draw[fleche] (e1) -- (e2);
\draw[fleche] (e2) -- (e3);
\draw[fleche] (e3) -- (e4);
\end{tikzpicture}
\legende{Organigramme de la démarche de résolution d'une équation trigonométrique.}
\end{popfigure}

\courssub{Équations du type $\cos(ax+b) = c$ ou $\sin(ax+b) = c$}

\textbf{Intuition.} L'argument n'est plus $x$ mais une expression $ax + b$. L'idée est de traiter cet argument comme une \emph{seule} inconnue : on pose $X = ax + b$, on résout en $X$, puis on revient à $x$. Le point délicat est l'intervalle : si $x$ parcourt $[0\,; 2\pi]$, alors $X = ax+b$ parcourt un intervalle \emph{plus grand}, qu'il faut calculer.

\begin{methode}[Changement d'inconnue $X = ax + b$]
\begin{enumerate}
\item Poser $X = ax + b$ et déterminer le nouvel intervalle : si $x \in [0\,; 2\pi]$ et $a > 0$, alors $X \in [b\,; 2a\pi + b]$.
\item Résoudre l'équation élémentaire en $X$.
\item En déduire $x = \dfrac{X - b}{a}$ pour chaque valeur de $X$ trouvée.
\item Ne garder que les $x$ appartenant à l'intervalle initial.
\end{enumerate}
\end{methode}

\begin{exoresolu}[Équation avec argument $2x + \frac{\pi}{6}$]
Résoudre dans $[0\,; 2\pi]$ l'équation $\cos\left(2x + \dfrac{\pi}{6}\right) = \dfrac{1}{2}$.
\tcblower
\textbf{Étape 1.} Posons $X = 2x + \dfrac{\pi}{6}$. Quand $x \in [0\,; 2\pi]$, on a $X \in \left[\dfrac{\pi}{6}\,; 4\pi + \dfrac{\pi}{6}\right]$.

\textbf{Étape 2.} On résout $\cos X = \dfrac{1}{2} = \cos\dfrac{\pi}{3}$ :
\[ X = \frac{\pi}{3} + 2k\pi \quad \text{ou} \quad X = -\frac{\pi}{3} + 2k\pi, \quad k \in \mathbb{Z}. \]

\textbf{Étape 3.} On revient à $x$ : $2x + \dfrac{\pi}{6} = X$, donc $x = \dfrac{X - \frac{\pi}{6}}{2}$.
\begin{itemize}
\item Premier cas : $x = \dfrac{\frac{\pi}{3} - \frac{\pi}{6} + 2k\pi}{2} = \dfrac{\pi}{12} + k\pi$.
\item Deuxième cas : $x = \dfrac{-\frac{\pi}{3} - \frac{\pi}{6} + 2k\pi}{2} = -\dfrac{\pi}{4} + k\pi$.
\end{itemize}

\textbf{Étape 4.} On sélectionne les solutions dans $[0\,; 2\pi]$ :
\begin{itemize}
\item $x = \dfrac{\pi}{12} + k\pi$ : pour $k = 0$, $x = \dfrac{\pi}{12}$ ; pour $k = 1$, $x = \dfrac{13\pi}{12}$.
\item $x = -\dfrac{\pi}{4} + k\pi$ : pour $k = 1$, $x = \dfrac{3\pi}{4}$ ; pour $k = 2$, $x = \dfrac{7\pi}{4}$.
\end{itemize}
\[ \boxed{S = \left\{ \frac{\pi}{12}\,; \frac{3\pi}{4}\,; \frac{13\pi}{12}\,; \frac{7\pi}{4} \right\}} \]
On obtient \emph{quatre} solutions : c'est normal, car l'argument $2x + \frac{\pi}{6}$ « tourne deux fois plus vite » que $x$.
\end{exoresolu}

\begin{attention}[Ne pas oublier d'ajuster l'intervalle]
Avec $X = ax + b$, le nombre de solutions est en général multiplié par $|a|$. Si l'on cherche $k$ dans le même intervalle que pour une équation en $x$, on \emph{perd} des solutions. Toujours tester plusieurs valeurs de $k$ ($-1, 0, 1, 2, \ldots$) et vérifier pour chacune que $x$ appartient bien à l'intervalle demandé.
\end{attention}

\courssub{Équations factorisables : $\sin 2x = \sin x$, $\cos 3x = \cos x$}

Deux stratégies équivalentes s'offrent à nous.

\textbf{Stratégie A : utiliser directement les formes $\cos p = \cos q$ et $\sin p = \sin q$.} On sait que :
\[ \cos p = \cos q \iff p = q + 2k\pi \ \text{ou}\ p = -q + 2k\pi ; \qquad \sin p = \sin q \iff p = q + 2k\pi \ \text{ou}\ p = \pi - q + 2k\pi. \]

\textbf{Stratégie B : factoriser grâce aux formules de transformation somme $\to$ produit.} On écrit l'équation sous la forme $\sin p - \sin q = 0$ ou $\cos p - \cos q = 0$, puis on applique :
\[ \sin p - \sin q = 2\cos\frac{p+q}{2}\sin\frac{p-q}{2} \qquad ; \qquad \cos p - \cos q = -2\sin\frac{p+q}{2}\sin\frac{p-q}{2}. \]
Un produit de facteurs est nul si et seulement si l'un des facteurs est nul : on retombe sur des équations élémentaires.

\begin{exoresolu}[Résolution de $\sin 3x = \sin x$ par factorisation]
Résoudre dans $[0\,; 2\pi]$ l'équation $\sin 3x = \sin x$.
\tcblower
\textbf{Étape 1.} On écrit $\sin 3x - \sin x = 0$ et on applique la formule avec $p = 3x$, $q = x$ :
\[ \sin 3x - \sin x = 2\cos\frac{3x+x}{2}\sin\frac{3x-x}{2} = 2\cos 2x \, \sin x. \]
L'équation devient : $2\cos 2x \, \sin x = 0$, c'est-à-dire $\cos 2x = 0$ \textbf{ou} $\sin x = 0$.

\textbf{Étape 2.} Résolution de chaque facteur.
\begin{itemize}
\item $\sin x = 0 \iff x = k\pi$. Dans $[0\,; 2\pi]$ : $x = 0$, $x = \pi$, $x = 2\pi$.
\item $\cos 2x = 0 \iff 2x = \dfrac{\pi}{2} + k\pi \iff x = \dfrac{\pi}{4} + \dfrac{k\pi}{2}$.
\end{itemize}

\textbf{Étape 3 et 4.} Pour $\cos 2x = 0$, on teste $k = 0, 1, 2, 3$ : $x = \dfrac{\pi}{4}$, $\dfrac{3\pi}{4}$, $\dfrac{5\pi}{4}$, $\dfrac{7\pi}{4}$, toutes dans $[0\,; 2\pi]$.

\[ \boxed{S = \left\{ 0\,; \frac{\pi}{4}\,; \frac{3\pi}{4}\,; \pi\,; \frac{5\pi}{4}\,; \frac{7\pi}{4}\,; 2\pi \right\}} \]
\end{exoresolu}

\begin{exemplebox}[Résolution de $\cos 3x = \cos x$]
$\cos 3x = \cos x \iff 3x = x + 2k\pi$ ou $3x = -x + 2k\pi$, c'est-à-dire $x = k\pi$ ou $x = \dfrac{k\pi}{2}$. La deuxième famille contient la première, donc $x = \dfrac{k\pi}{2}$. Dans $[0\,; 2\pi]$ : $S = \left\{0\,; \dfrac{\pi}{2}\,; \pi\,; \dfrac{3\pi}{2}\,; 2\pi\right\}$.
\end{exemplebox}

\courssub{Équations du type $a\cos x + b\sin x = 0$}

\textbf{Intuition.} Une telle équation mélange $\cos x$ et $\sin x$ au même degré. En divisant par $\cos x$ (lorsqu'il est non nul), on fait apparaître $\tan x$, et l'on retombe sur une équation élémentaire.

\begin{methode}[Résolution de $a\cos x + b\sin x = 0$]
\begin{enumerate}
\item Discuter le cas $\cos x = 0$ (c'est-à-dire $x = \frac{\pi}{2} + k\pi$) : pour ces valeurs, $\sin x = \pm 1$, l'équation devient $\pm b = 0$.
  \begin{itemize}
  \item Si $b = 0$, ces valeurs sont solutions (l'équation se réduit à $a\cos x = 0$).
  \item Si $b \neq 0$, ces valeurs ne sont pas solutions ; on peut alors diviser par $\cos x$ sans risque.
  \end{itemize}
\item Diviser par $\cos x$ (si autorisé) : $a + b\tan x = 0$, soit $\tan x = -\dfrac{a}{b}$.
\item Résoudre cette équation élémentaire : $x = \alpha + k\pi$ où $\alpha$ est une valeur telle que $\tan\alpha = -\frac{a}{b}$.
\end{enumerate}
\end{methode}

\begin{exemplebox}[Exemple chiffré]
Résoudre $\cos x - \sqrt{3}\sin x = 0$ dans $[0\,; 2\pi]$.
\begin{itemize}
\item Si $\cos x = 0$, l'équation donne $-\sqrt{3}\sin x = 0$, impossible car alors $\sin x = \pm 1$. On peut diviser par $\cos x$.
\item $1 - \sqrt{3}\tan x = 0$, donc $\tan x = \dfrac{1}{\sqrt{3}} = \tan\dfrac{\pi}{6}$.
\item $x = \dfrac{\pi}{6} + k\pi$. Dans $[0\,; 2\pi]$ : $S = \left\{\dfrac{\pi}{6}\,; \dfrac{7\pi}{6}\right\}$.
\end{itemize}
\end{exemplebox}

\begin{attention}[Diviser par $\sin x$ ou $\cos x$ : le piège classique]
Diviser une équation par $\sin x$ ou $\cos x$ \emph{sans vérifier} le cas où ce facteur est nul fait perdre des solutions — ou en introduire de fausses. Par exemple, pour $\sin 2x = \sin x$, si l'on écrit $2\sin x \cos x = \sin x$ puis que l'on divise par $\sin x$, on obtient $\cos x = \frac{1}{2}$ et l'on \textbf{perd} toutes les solutions $x = k\pi$ ! La bonne méthode : \emph{tout passer dans un membre et factoriser}, jamais simplifier par une expression qui peut s'annuler.
\end{attention}

\courssub{Équations du second degré en $\cos x$ ou $\sin x$}

\textbf{Intuition.} Certaines équations ressemblent à des équations du second degré, où le rôle de l'inconnue est joué par $\cos x$ ou $\sin x$. On pose alors $X = \cos x$ (ou $X = \sin x$), on résout l'équation du second degré en $X$, puis on revient à $x$ — en n'oubliant pas que $X$ doit appartenir à $[-1\,; 1]$.

\begin{methode}[Équation du second degré en $\cos x$ ou $\sin x$]
\begin{enumerate}
\item Poser $X = \cos x$ (ou $X = \sin x$) et préciser la condition $-1 \leq X \leq 1$.
\item Résoudre l'équation du second degré en $X$ (discriminant $\Delta$).
\item \textbf{Écarter} toute racine hors de $[-1\,; 1]$ (solution parasite).
\item Pour chaque racine retenue, résoudre l'équation élémentaire $\cos x = X$ (ou $\sin x = X$).
\item Sélectionner les solutions dans l'intervalle demandé.
\end{enumerate}
\end{methode}

\begin{exoresolu}[{Résolution de $2\sin^2 x - \sin x - 1 = 0$ sur $[0\,; 2\pi]$}]
Résoudre sur $[0\,; 2\pi]$ l'équation $2\sin^2 x - \sin x - 1 = 0$.
\tcblower
\textbf{Étape 1.} Posons $X = \sin x$, avec la condition $-1 \leq X \leq 1$. L'équation devient :
\[ 2X^2 - X - 1 = 0. \]

\textbf{Étape 2.} Discriminant : $\Delta = (-1)^2 - 4 \times 2 \times (-1) = 1 + 8 = 9 > 0$. Deux racines :
\[ X_1 = \frac{1 - 3}{4} = -\frac{1}{2} \qquad ; \qquad X_2 = \frac{1 + 3}{4} = 1. \]

\textbf{Étape 3.} Les deux racines appartiennent à $[-1\,; 1]$ : on les garde toutes les deux.

\textbf{Étape 4.} Retour à $x$.
\begin{itemize}
\item $\sin x = 1 \iff x = \dfrac{\pi}{2} + 2k\pi$. Dans $[0\,; 2\pi]$ : $x = \dfrac{\pi}{2}$.
\item $\sin x = -\dfrac{1}{2} = \sin\left(-\dfrac{\pi}{6}\right)$, donc $x = -\dfrac{\pi}{6} + 2k\pi$ ou $x = \pi + \dfrac{\pi}{6} + 2k\pi = \dfrac{7\pi}{6} + 2k\pi$. Dans $[0\,; 2\pi]$ : $x = \dfrac{11\pi}{6}$ et $x = \dfrac{7\pi}{6}$.
\end{itemize}

\textbf{Vérification.} Pour $x = \frac{7\pi}{6}$ : $\sin x = -\frac{1}{2}$, et $2\times\frac{1}{4} - \left(-\frac{1}{2}\right) - 1 = \frac{1}{2} + \frac{1}{2} - 1 = 0$. Correct.

\[ \boxed{S = \left\{ \frac{\pi}{2}\,; \frac{7\pi}{6}\,; \frac{11\pi}{6} \right\}} \]
\end{exoresolu}

\begin{popfigure}
\begin{tikzpicture}[scale=1.05]
\draw[thick, popdark] (0,0) circle (2);
\draw[-{Stealth}, popmuted] (-2.5,0) -- (2.5,0) node[below left, font=\small] {$\cos$};
\draw[-{Stealth}, popmuted] (0,-2.5) -- (0,2.5) node[below left, font=\small] {$\sin$};
% sin x = 1 : point pi/2
\fill[popgreen] (0,2) circle (2.5pt) node[above right, popgreen, font=\small] {$\frac{\pi}{2}$};
% sin x = -1/2 : droite horizontale y = -1
\draw[poporange, thick] (-1.732,-1) -- (1.732,-1);
\fill[poporange] (-1.732,-1) circle (2.5pt) node[below left, poporange, font=\small] {$\frac{7\pi}{6}$};
\fill[poporange] (1.732,-1) circle (2.5pt) node[below right, poporange, font=\small] {$\frac{11\pi}{6}$};
\draw[poporange, dashed] (0,-1) node[left, font=\small, poporange] {$-\frac{1}{2}$} -- (-1.732,-1);
\end{tikzpicture}
\legende{Les trois solutions de $2\sin^2 x - \sin x - 1 = 0$ sur le cercle trigonométrique : la droite $\sin x = -\frac{1}{2}$ coupe le cercle en deux points, et $\sin x = 1$ donne le point sommet.}
\end{popfigure}

\begin{attention}[Solutions parasites et conditions d'existence]
\begin{itemize}
\item Toute racine $X$ de l'équation du second degré telle que $|X| > 1$ doit être \textbf{rejetée} : $\cos x$ et $\sin x$ ne dépassent jamais $1$ en valeur absolue. Par exemple, $2\cos^2 x - 3\cos x - 2 = 0$ donne $X = 2$ ou $X = -\frac{1}{2}$ : seule $X = -\frac{1}{2}$ est recevable.
\item Si l'équation contient $\tan x$, penser à la condition d'existence : $x \neq \frac{\pi}{2} + k\pi$. Toute solution trouvée qui viole cette condition est parasite et doit être écartée.
\end{itemize}
\end{attention}

\courssub{Rédaction complète et justifiée : un modèle à imiter}

Au Probatoire et au Baccalauréat technique, la qualité de la rédaction compte autant que le résultat. Voici la structure attendue.

\begin{methode}[Rédaction type d'une résolution]
\begin{enumerate}
\item \textbf{Annoncer} l'ensemble de résolution : « On résout dans $[0\,; 2\pi]$ l'équation (E) : ... ».
\item \textbf{Justifier chaque transformation} : citer la formule utilisée (transformation somme $\to$ produit, changement d'inconnue, produit nul).
\item \textbf{Écrire les solutions générales} avec le paramètre $k \in \mathbb{Z}$ avant de sélectionner.
\item \textbf{Sélectionner explicitement} : indiquer pour quelles valeurs de $k$ la solution appartient à l'intervalle.
\item \textbf{Conclure} par l'ensemble $S$ encadré, et vérifier si possible une solution.
\end{enumerate}
\end{methode}

\begin{exercice}[Pour s'entraîner]
\begin{enumerate}
\item Résoudre dans $[0\,; 2\pi]$ : $\sin\left(2x - \dfrac{\pi}{3}\right) = \dfrac{\sqrt{3}}{2}$.
\item Résoudre dans $[0\,; 2\pi]$ : $\cos 2x = \cos\left(x + \dfrac{\pi}{4}\right)$.
\item Résoudre dans $[0\,; 2\pi]$ : $2\cos^2 x + \cos x - 1 = 0$.
\item Résoudre dans $[0\,; 2\pi]$ : $\sqrt{3}\cos x + \sin x = 0$.
\end{enumerate}
\end{exercice}

\begin{corrige}[Exercice 1]
\textbf{1.} Posons $X = 2x - \frac{\pi}{3}$. $\sin X = \frac{\sqrt{3}}{2} = \sin\frac{\pi}{3}$ donne $X = \frac{\pi}{3} + 2k\pi$ ou $X = \frac{2\pi}{3} + 2k\pi$. D'où $x = \frac{\pi}{3} + k\pi$ ou $x = \frac{\pi}{2} + k\pi$. Dans $[0\,; 2\pi]$ : $S = \left\{\frac{\pi}{3}\,; \frac{\pi}{2}\,; \frac{4\pi}{3}\,; \frac{3\pi}{2}\right\}$.

\textbf{2.} $\cos 2x = \cos\left(x + \frac{\pi}{4}\right) \iff 2x = x + \frac{\pi}{4} + 2k\pi$ ou $2x = -x - \frac{\pi}{4} + 2k\pi$, soit $x = \frac{\pi}{4} + 2k\pi$ ou $x = -\frac{\pi}{12} + \frac{2k\pi}{3}$. Dans $[0\,; 2\pi]$ : $S = \left\{\frac{\pi}{4}\,; \frac{7\pi}{12}\,; \frac{5\pi}{4}\,; \frac{23\pi}{12}\right\}$.

\textbf{3.} Posons $X = \cos x$ : $2X^2 + X - 1 = 0$, $\Delta = 9$, $X = \frac{1}{2}$ ou $X = -1$ (toutes deux recevables). $\cos x = \frac{1}{2}$ donne $x = \pm\frac{\pi}{3} + 2k\pi$ ; $\cos x = -1$ donne $x = \pi + 2k\pi$. Dans $[0\,; 2\pi]$ : $S = \left\{\frac{\pi}{3}\,; \pi\,; \frac{5\pi}{3}\right\}$.

\textbf{4.} Si $\cos x = 0$, l'équation donne $\sin x = 0$ : impossible. On divise par $\cos x$ : $\tan x = -\sqrt{3} = \tan\left(-\frac{\pi}{3}\right)$, donc $x = -\frac{\pi}{3} + k\pi$. Dans $[0\,; 2\pi]$ : $S = \left\{\frac{2\pi}{3}\,; \frac{5\pi}{3}\right\}$.
\end{corrige}

\begin{aretenir}[L'essentiel de la section]
\begin{itemize}
\item Toute équation trigonométrique se résout en la \cle{ramenant} à une ou plusieurs formes élémentaires.
\item Argument $ax + b$ : poser $X = ax + b$, résoudre en $X$, revenir à $x$ en testant plusieurs valeurs de $k$.
\item $\sin p = \sin q$ ou $\cos p = \cos q$ : factoriser par les formules de transformation (produit nul) ou appliquer directement les formes élémentaires.
\item $a\cos x + b\sin x = 0$ : vérifier que $\cos x = 0$ n'est pas solution (si $b \neq 0$), puis diviser pour obtenir $\tan x = -\frac{a}{b}$.
\item Second degré en $\cos x$ ou $\sin x$ : poser $X$, résoudre, \textbf{rejeter} toute racine hors de $[-1\,; 1]$.
\item Ne jamais diviser par $\sin x$ ou $\cos x$ sans discuter le cas où le facteur est nul.
\end{itemize}
\end{aretenir}

\coursec{Inéquations trigonométriques}

Dans la section précédente, nous avons appris à résoudre des \cle{équations} trigonométriques : nous cherchions les valeurs \emph{exactes} de $x$ pour lesquelles une égalité est vraie, par exemple $\cos x = \dfrac{1}{2}$. Dans la vie courante et dans les métiers techniques, on se pose souvent une question différente : \emph{pendant quelle période} une grandeur dépasse-t-elle un seuil ? Pendant quelle fraction du tour un angle garde-t-il son sinus supérieur à $\dfrac{1}{2}$ ? Pendant combien d'heures un panneau solaire reçoit-il un ensoleillement suffisant ? Ces questions conduisent à des \cle{inéquations} trigonométriques : les solutions ne sont plus des points isolés, mais des \emph{arcs} du cercle trigonométrique, c'est-à-dire des \emph{intervalles}.

\courssub{Principe général : lire les solutions sur le cercle trigonométrique}

\begin{definition}[Inéquation trigonométrique]
Une \cle{inéquation trigonométrique} d'inconnue $x$ est une inégalité dans laquelle figurent $\cos x$, $\sin x$ ou $\tan x$, par exemple $\sin x \geq \dfrac{1}{2}$ ou $2\cos x + 1 < 0$. La résoudre, c'est déterminer \textbf{tous les intervalles} de valeurs de $x$ (dans un ensemble donné, le plus souvent $[0\,;\,2\pi]$) qui rendent l'inégalité vraie.
\end{definition}

L'idée directrice est simple : le cercle trigonométrique est un \emph{tableau de lecture} visuel.
\begin{itemize}
\item $\cos x$ se lit sur l'\textbf{axe horizontal} (axe des cosinus) ;
\item $\sin x$ se lit sur l'\textbf{axe vertical} (axe des sinus).
\end{itemize}
Comparer $\cos x$ ou $\sin x$ à un nombre $a$, c'est donc comparer la position d'un point du cercle à une droite horizontale ou verticale. Les points du cercle qui conviennent forment un ou plusieurs \textbf{arcs}, que l'on traduit ensuite en intervalles.

\begin{methode}[{Résoudre $\cos x \geq a$, $\sin x \geq a$ (ou avec $<$, $\leq$) sur $[0\,;\,2\pi]$}]
\begin{enumerate}
\item \textbf{Placer le seuil} : sur l'axe des cosinus (droite verticale $x = a$) pour une inéquation en $\cos x$, sur l'axe des sinus (droite horizontale $y = a$) pour une inéquation en $\sin x$.
\item \textbf{Trouver les points frontières} : résoudre l'équation associée ($\cos x = a$ ou $\sin x = a$) sur $[0\,;\,2\pi]$ ; ce sont les extrémités de l'arc solution.
\item \textbf{Repérer l'arc solution} : pour $\sin x \geq a$, on garde les points du cercle \emph{au-dessus} de la droite horizontale ; pour $\cos x \geq a$, les points \emph{à droite} de la droite verticale ; l'inégalité inverse donne la partie complémentaire.
\item \textbf{Écrire l'intervalle} en respectant le \emph{sens de parcours} du cercle (sens trigonométrique, de la borne la plus petite vers la plus grande sur $[0\,;\,2\pi]$), avec des crochets fermés si l'inégalité est large ($\geq$, $\leq$), ouverts si elle est stricte ($>$, $<$).
\end{enumerate}
\end{methode}

\begin{attention}[L'erreur classique : l'intervalle écrit dans le mauvais sens]
Sur le cercle, l'arc solution de $\cos x < -\dfrac{\sqrt{2}}{2}$ passe « par $\pi$ » : ses extrémités sont $\dfrac{3\pi}{4}$ et $\dfrac{5\pi}{4}$. Il faut écrire $\left]\dfrac{3\pi}{4}\,;\,\dfrac{5\pi}{4}\right[$ et surtout \textbf{pas} $\left]\dfrac{5\pi}{4}\,;\,\dfrac{3\pi}{4}\right[$, qui est vide ! Règle de sécurité : sur $[0\,;\,2\pi]$, un intervalle s'écrit toujours avec la borne \emph{inférieure} à gauche. Quand l'arc solution « enjambe » $0$ (par exemple pour $\cos x \geq \dfrac{1}{2}$), la solution se coupe en \textbf{deux} intervalles : $\left[0\,;\,\dfrac{\pi}{3}\right] \cup \left[\dfrac{5\pi}{3}\,;\,2\pi\right]$.
\end{attention}

\courssub{Inéquations avec le sinus : exemple fondateur}

\begin{exoresolu}[{Résoudre $\sin x \geq \dfrac{1}{2}$ sur $[0\,;\,2\pi]$}]
\textbf{Énoncé.} Déterminer l'ensemble des réels $x$ de $[0\,;\,2\pi]$ tels que $\sin x \geq \dfrac{1}{2}$.
\tcblower
\textbf{Solution détaillée.}
\begin{enumerate}
\item \textbf{Équation associée} : $\sin x = \dfrac{1}{2}$. Sur $[0\,;\,2\pi]$, les solutions sont $x = \dfrac{\pi}{6}$ et $x = \pi - \dfrac{\pi}{6} = \dfrac{5\pi}{6}$.
\item \textbf{Lecture sur le cercle} : on trace la droite horizontale d'ordonnée $\dfrac{1}{2}$. Les points du cercle dont le sinus est \emph{supérieur} à $\dfrac{1}{2}$ sont ceux situés \emph{au-dessus} de cette droite : c'est l'arc allant de $\dfrac{\pi}{6}$ à $\dfrac{5\pi}{6}$ en passant par $\dfrac{\pi}{2}$.
\item \textbf{Conclusion} : l'inégalité est large, donc les bornes sont incluses :
\[ S = \left[\dfrac{\pi}{6}\,;\,\dfrac{5\pi}{6}\right]. \]
\end{enumerate}
\textbf{Vérification} : $\sin\dfrac{\pi}{2} = 1 \geq \dfrac{1}{2}$ (dedans, correct) ; $\sin\pi = 0 < \dfrac{1}{2}$ (dehors, correct).
\end{exoresolu}

\begin{popfigure}
\begin{center}
\begin{tikzpicture}[scale=2.6]
\draw[popline] (-1.3,0) -- (1.3,0) node[below right] {\footnotesize cosinus};
\draw[popline] (0,-1.3) -- (0,1.3) node[above left] {\footnotesize sinus};
\draw[popdark,thick] (0,0) circle (1);
% arc solution hachuré : de 30° à 150°
\fill[pattern=north east lines, pattern color=popblue] (0.866,0.5) arc (30:150:1) -- (0,0) -- cycle;
\draw[popblue,very thick] (0.866,0.5) arc (30:150:1);
% droite seuil y = 1/2
\draw[poporange,dashed,thick] (-1.25,0.5) -- (1.25,0.5) node[right, poporange] {\footnotesize $y=\frac{1}{2}$};
% points frontières
\fill[poporange] (0.866,0.5) circle (0.02);
\fill[poporange] (-0.866,0.5) circle (0.02);
\draw[popline] (0,0) -- (0.866,0.5);
\draw[popline] (0,0) -- (-0.866,0.5);
\node[below right, popdark] at (0.55,0.28) {\footnotesize $\frac{\pi}{6}$};
\node[below left, popdark] at (-0.55,0.28) {\footnotesize $\frac{5\pi}{6}$};
\node[popblue] at (0,0.85) {\footnotesize \textbf{arc solution}};
\fill[popdark] (1,0) circle (0.015) node[below right] {\footnotesize $0$};
\fill[popdark] (0,1) circle (0.015) node[above right] {\footnotesize $\frac{\pi}{2}$};
\end{tikzpicture}
\end{center}
\legende{Cercle trigonométrique : les points dont le sinus est supérieur ou égal à $\frac{1}{2}$ forment l'arc hachuré, de $\frac{\pi}{6}$ à $\frac{5\pi}{6}$.}
\end{popfigure}

On peut confirmer cette lecture graphiquement avec la courbe de la fonction sinus : les solutions sont les abscisses pour lesquelles la courbe est \emph{au-dessus} de la droite $y = \dfrac{1}{2}$.

\begin{popfigure}
\begin{center}
\begin{tikzpicture}
\begin{axis}[
width=0.95\linewidth, height=5.2cm,
axis lines=middle,
xlabel={$x$}, ylabel={$y$},
xmin=-0.3, xmax=6.6, ymin=-1.3, ymax=1.4,
xtick={0.5236, 1.5708, 2.618, 3.1416, 4.7124, 6.2832},
xticklabels={$\frac{\pi}{6}$, $\frac{\pi}{2}$, $\frac{5\pi}{6}$, $\pi$, $\frac{3\pi}{2}$, $2\pi$},
ytick={-1, 0.5, 1},
yticklabels={$-1$, $\frac{1}{2}$, $1$},
tick label style={font=\footnotesize},
]
\addplot[poporange, dashed, thick, domain=0:6.2832] {0.5};
\addplot[popcurve, domain=0:6.2832, samples=200] {sin(deg(x))};
\addplot[popblue!25, draw=none, domain=0.5236:2.618, samples=60] {sin(deg(x))} \closedcycle;
\node[poporange, font=\footnotesize] at (axis cs:4.6,0.62) {$y=\frac{1}{2}$};
\node[popblue, font=\footnotesize] at (axis cs:1.5708,1.22) {zone solution};
\end{axis}
\end{tikzpicture}
\end{center}
\legende{Courbe $y=\sin x$ et droite $y=\frac{1}{2}$ : la zone colorée correspond à $\sin x \geq \frac{1}{2}$, soit $x\in\left[\frac{\pi}{6}\,;\,\frac{5\pi}{6}\right]$ (puis l'arc se répète par périodicité).}
\end{popfigure}

\courssub{Inéquations avec le cosinus}

Pour $\cos x$, le seuil se lit sur l'axe \emph{horizontal} : on trace la droite \emph{verticale} d'abscisse $a$.
\begin{itemize}
\item $\cos x \geq a$ : on garde les points du cercle \emph{à droite} de cette droite (arc passant par $0$) ;
\item $\cos x < a$ : on garde les points \emph{à gauche} (arc passant par $\pi$).
\end{itemize}

\begin{exoresolu}[{Résoudre $2\cos x + 1 < 0$ sur $[0\,;\,2\pi]$}]
\textbf{Énoncé.} Déterminer l'ensemble des réels $x$ de $[0\,;\,2\pi]$ tels que $2\cos x + 1 < 0$.
\tcblower
\textbf{Solution détaillée.}
\begin{enumerate}
\item \textbf{Ramener à la forme simple} : $2\cos x + 1 < 0 \iff \cos x < -\dfrac{1}{2}$.
\item \textbf{Équation associée} : $\cos x = -\dfrac{1}{2}$. L'angle de référence est $\dfrac{\pi}{3}$ (car $\cos\dfrac{\pi}{3} = \dfrac{1}{2}$). Le cosinus est négatif dans les quadrants II et III :
\[ x = \pi - \dfrac{\pi}{3} = \dfrac{2\pi}{3} \qquad \text{et} \qquad x = \pi + \dfrac{\pi}{3} = \dfrac{4\pi}{3}. \]
\item \textbf{Lecture sur le cercle} : la droite verticale d'abscisse $-\dfrac{1}{2}$ partage le cercle. Les points dont le cosinus est \emph{strictement inférieur} à $-\dfrac{1}{2}$ sont \emph{à gauche} de cette droite : c'est l'arc passant par $\pi$, de $\dfrac{2\pi}{3}$ à $\dfrac{4\pi}{3}$.
\item \textbf{Conclusion} : l'inégalité est stricte, donc crochets ouverts :
\[ S = \left]\dfrac{2\pi}{3}\,;\,\dfrac{4\pi}{3}\right[. \]
\end{enumerate}
\textbf{Vérification} : $2\cos\pi + 1 = -1 < 0$ (dedans, correct) ; $2\cos 0 + 1 = 3 > 0$ (dehors, correct).
\end{exoresolu}

\begin{popfigure}
\begin{center}
\begin{tikzpicture}[scale=2.6]
\draw[popline] (-1.3,0) -- (1.3,0) node[below right] {\footnotesize cosinus};
\draw[popline] (0,-1.3) -- (0,1.3) node[above left] {\footnotesize sinus};
\draw[popdark,thick] (0,0) circle (1);
% arc solution pour cos x < -1/2 : de 120° à 240°
\fill[pattern=north west lines, pattern color=poppurple] (-0.5,0.866) arc (120:240:1) -- (0,0) -- cycle;
\draw[poppurple,very thick] (-0.5,0.866) arc (120:240:1);
% droite seuil x = -1/2
\draw[poporange,dashed,thick] (-0.5,-1.25) -- (-0.5,1.25) node[above, poporange] {\footnotesize $x=-\frac{1}{2}$};
\fill[poporange] (-0.5,0.866) circle (0.02);
\fill[poporange] (-0.5,-0.866) circle (0.02);
\draw[popline] (0,0) -- (-0.5,0.866);
\draw[popline] (0,0) -- (-0.5,-0.866);
\node[above left, popdark] at (-0.75,0.65) {\footnotesize $\frac{2\pi}{3}$};
\node[below left, popdark] at (-0.75,-0.65) {\footnotesize $\frac{4\pi}{3}$};
\node[poppurple] at (-1.05,0.05) {\footnotesize \textbf{arc}};
\node[poppurple] at (-1.05,-0.12) {\footnotesize \textbf{solution}};
\fill[popdark] (1,0) circle (0.015) node[below right] {\footnotesize $0$};
\fill[popdark] (-1,0) circle (0.015) node[above left] {\footnotesize $\pi$};
\end{tikzpicture}
\end{center}
\legende{Arc solution de $\cos x < -\frac{1}{2}$ sur $[0\,;\,2\pi]$ : les points à gauche de la droite $x=-\frac{1}{2}$, soit $x\in\left]\frac{2\pi}{3}\,;\,\frac{4\pi}{3}\right[$.}
\end{popfigure}

\courssub{Inéquations avec la tangente : les valeurs interdites}

La tangente se lit sur la \textbf{tangente au cercle au point $(1\,;\,0)$} : pour un point $M$ du cercle, $\tan x$ est l'ordonnée du point d'intersection de la droite $(OM)$ avec cette tangente verticale.

\begin{attention}[Valeurs interdites pour $\tan x$]
$\tan x = \dfrac{\sin x}{\cos x}$ n'existe que si $\cos x \neq 0$, c'est-à-dire $x \neq \dfrac{\pi}{2} + k\pi$ ($k\in\mathbb{Z}$). Sur $[0\,;\,2\pi]$, les valeurs $\dfrac{\pi}{2}$ et $\dfrac{3\pi}{2}$ sont donc \textbf{toujours exclues} de l'ensemble des solutions, même si l'inégalité est large. De plus, la tangente est \emph{croissante} sur chacun de ses intervalles de définition, mais elle « saute » de $+\infty$ à $-\infty$ en $\dfrac{\pi}{2}$ : une inéquation comme $\tan x \leq 1$ a pour solution sur $[0\,;\,2\pi]$ :
\[ S = \left[0\,;\,\dfrac{\pi}{4}\right] \cup \left]\dfrac{\pi}{2}\,;\,\dfrac{5\pi}{4}\right] \cup \left]\dfrac{3\pi}{2}\,;\,2\pi\right]. \]
Oublier d'exclure $\dfrac{\pi}{2}$ et $\dfrac{3\pi}{2}$ est l'erreur la plus fréquente au Probatoire.
\end{attention}

\begin{exemplebox}[{Résoudre $\tan x \geq 1$ sur $[0\,;\,2\pi]$}]
\textbf{Équation associée} : $\tan x = 1 \iff x = \dfrac{\pi}{4}$ ou $x = \dfrac{5\pi}{4}$ sur $[0\,;\,2\pi]$.\\
Sur $\left[0\,;\,\dfrac{\pi}{2}\right[$, la tangente croît de $0$ à $+\infty$ : $\tan x \geq 1$ pour $x \in \left[\dfrac{\pi}{4}\,;\,\dfrac{\pi}{2}\right[$.\\
Sur $\left]\dfrac{\pi}{2}\,;\,\dfrac{3\pi}{2}\right[$, la tangente croît de $-\infty$ à $+\infty$ : $\tan x \geq 1$ pour $x \in \left[\dfrac{5\pi}{4}\,;\,\dfrac{3\pi}{2}\right[$.\\
Sur $\left]\dfrac{3\pi}{2}\,;\,2\pi\right]$, la tangente est négative : aucune solution.
\[ S = \left[\dfrac{\pi}{4}\,;\,\dfrac{\pi}{2}\right[ \;\cup\; \left[\dfrac{5\pi}{4}\,;\,\dfrac{3\pi}{2}\right[. \]
Notez les crochets \emph{ouverts} en $\dfrac{\pi}{2}$ et $\dfrac{3\pi}{2}$ : valeurs interdites.
\end{exemplebox}

\courssub{Inéquations se ramenant aux formes simples}

Comme pour les équations, de nombreuses inéquations se ramènent aux formes $\sin x \lesseqgtr a$, $\cos x \lesseqgtr a$, $\tan x \lesseqgtr a$ par \textbf{transformation} (formules d'addition, de duplication, produits-sommes vus dans cette leçon) ou par \textbf{factorisation}.

\begin{methode}[Résolution par factorisation et tableau de signes]
\begin{enumerate}
\item Regrouper tous les termes d'un même côté pour obtenir une inégalité de la forme $E(x) \geq 0$ (ou $>$, $\leq$, $<$).
\item Factoriser $E(x)$ en produit de facteurs simples ($\cos x$, $\sin x - a$, $2\cos x + 1$, etc.).
\item Résoudre les \emph{équations} associées à chaque facteur : elles donnent les valeurs frontières sur $[0\,;\,2\pi]$.
\item Construire un \textbf{tableau de signes} : une ligne par facteur, une ligne pour le produit.
\item Lire les intervalles où le produit a le signe demandé, en respectant les crochets (ouverts si inégalité stricte ou valeur interdite).
\end{enumerate}
\end{methode}

\begin{exoresolu}[{Étude de signe : résoudre $\sin x\,(2\cos x - 1) \geq 0$ sur $[0\,;\,2\pi]$}]
\textbf{Énoncé.} Déterminer le signe du produit $P(x) = \sin x\,(2\cos x - 1)$ et résoudre $P(x) \geq 0$ sur $[0\,;\,2\pi]$.
\tcblower
\textbf{Solution détaillée.}
\begin{enumerate}
\item \textbf{Valeurs frontières} : $\sin x = 0 \iff x \in \{0\,;\,\pi\,;\,2\pi\}$ ; \quad $2\cos x - 1 = 0 \iff \cos x = \dfrac{1}{2} \iff x \in \left\{\dfrac{\pi}{3}\,;\,\dfrac{5\pi}{3}\right\}$.
\item \textbf{Signe de $\sin x$} : positif sur $]0\,;\,\pi[$, négatif sur $]\pi\,;\,2\pi[$.
\item \textbf{Signe de $2\cos x - 1$} : $\cos x \geq \dfrac{1}{2}$ sur $\left[0\,;\,\dfrac{\pi}{3}\right] \cup \left[\dfrac{5\pi}{3}\,;\,2\pi\right]$, donc le facteur est positif sur ces intervalles, négatif sur $\left]\dfrac{\pi}{3}\,;\,\dfrac{5\pi}{3}\right[$.
\item \textbf{Tableau de signes} :
\begin{center}
\begin{tabular}{|c|ccccccccc|}\hline
$x$ & $0$ & & $\frac{\pi}{3}$ & & $\pi$ & & $\frac{5\pi}{3}$ & & $2\pi$ \\ \hline
$\sin x$ & $0$ & $+$ & $+$ & $+$ & $0$ & $-$ & $-$ & $-$ & $0$ \\ \hline
$2\cos x - 1$ & $+$ & $+$ & $0$ & $-$ & $-$ & $-$ & $0$ & $+$ & $+$ \\ \hline
$P(x)$ & $0$ & $+$ & $0$ & $-$ & $0$ & $+$ & $0$ & $-$ & $0$ \\ \hline
\end{tabular}
\end{center}
\item \textbf{Conclusion} : $P(x) \geq 0$ (inégalité large, bornes incluses) :
\[ S = \left[0\,;\,\dfrac{\pi}{3}\right] \cup \left[\pi\,;\,\dfrac{5\pi}{3}\right] \cup \{2\pi\}. \]
\end{enumerate}
\end{exoresolu}

\begin{attention}[Deux pièges de rédaction]
\begin{itemize}
\item \textbf{Sens de l'inégalité} : quand on divise par un nombre \emph{négatif} (par exemple en passant de $-2\sin x > 1$ à $\sin x < -\dfrac{1}{2}$), le sens de l'inégalité \emph{change}. En revanche, on ne multiplie jamais par $\cos x$ ou $\sin x$ sans discuter le signe : on préfère factoriser.
\item \textbf{Inégalité stricte ou large} : $\geq$ donne des crochets fermés sur les frontières ; $>$ donne des crochets ouverts. Les valeurs interdites de $\tan x$ donnent \emph{toujours} des crochets ouverts.
\end{itemize}
\end{attention}

\courssub{Lien avec les situations concrètes : dépassement de seuil}

De nombreuses grandeurs techniques sont \emph{périodiques} et se modélisent par des fonctions sinusoïdales : tension alternative, hauteur d'une marée, ensoleillement, position d'un piston. Savoir résoudre une inéquation trigonométrique, c'est savoir répondre à la question : \emph{« pendant quelle fraction de la période la grandeur dépasse-t-elle le seuil ? »}

\begin{exemplebox}[Tension alternative dans un atelier]
La tension d'un circuit est modélisée par $u(t) = 220\sqrt{2}\sin(100\pi t)$ (en volts, $t$ en secondes). Un appareil ne fonctionne correctement que lorsque $u(t) \geq 220$. On résout :
\[ 220\sqrt{2}\sin(100\pi t) \geq 220 \iff \sin(100\pi t) \geq \dfrac{1}{\sqrt{2}} = \dfrac{\sqrt{2}}{2}. \]
Sur une période, en posant $x = 100\pi t$, on obtient $x \in \left[\dfrac{\pi}{4}\,;\,\dfrac{3\pi}{4}\right]$, c'est-à-dire la moitié de la demi-période positive : l'appareil est alimenté correctement pendant \textbf{un quart de la période totale}.
\end{exemplebox}

\begin{savaistu}[Les inéquations trigonométriques et l'énergie solaire au Cameroun]
L'ensoleillement reçu par un panneau solaire au cours d'une journée suit approximativement une courbe en sinus : nul au lever et au coucher du soleil, maximal vers midi. Pour dimensionner une installation (à Yaoundé, Douala ou Maroua), l'ingénieur résout une inéquation du type $\sin x \geq a$ pour déterminer les \textbf{plages horaires} où l'ensoleillement dépasse le seuil de rentabilité du panneau. C'est exactement la méthode de cette section : placer le seuil, repérer l'arc solution, convertir en durée. À Maroua, où l'ensoleillement est parmi les plus forts du pays, ces plages sont plus longues : les mathématiques expliquent pourquoi le Nord camerounais est privilégié pour les centrales solaires comme celle de Guider !
\end{savaistu}

\begin{aretenir}[L'essentiel de la section]
\begin{itemize}
\item Résoudre $\sin x \geq a$ ou $\cos x \geq a$ : placer le seuil sur l'axe des sinus (droite horizontale) ou des cosinus (droite verticale), repérer l'\textbf{arc solution} sur le cercle, écrire l'\textbf{intervalle} dans le bon sens.
\item $\sin x \geq \dfrac{1}{2}$ sur $[0\,;\,2\pi]$ : $S = \left[\dfrac{\pi}{6}\,;\,\dfrac{5\pi}{6}\right]$ ; \quad $2\cos x + 1 < 0$ : $S = \left]\dfrac{2\pi}{3}\,;\,\dfrac{4\pi}{3}\right[$.
\item Pour $\tan x$ : exclure \textbf{toujours} $\dfrac{\pi}{2}$ et $\dfrac{3\pi}{2}$ (valeurs interdites).
\item Inéquations complexes : factoriser, puis \textbf{tableau de signes} sur $[0\,;\,2\pi]$.
\item Inégalité large $\Rightarrow$ crochets fermés ; stricte $\Rightarrow$ crochets ouverts ; arc enjambant $0$ $\Rightarrow$ deux intervalles.
\end{itemize}
\end{aretenir}

\begin{exercice}[Pour s'entraîner]
\begin{enumerate}
\item Résoudre sur $[0\,;\,2\pi]$ : $\cos x \geq \dfrac{\sqrt{2}}{2}$.
\item Résoudre sur $[0\,;\,2\pi]$ : $2\sin x - \sqrt{3} < 0$.
\item Résoudre sur $[0\,;\,2\pi]$ : $\tan x < \sqrt{3}$.
\item À l'aide d'un tableau de signes, résoudre sur $[0\,;\,2\pi]$ : $\cos x\,(2\sin x - 1) \leq 0$.
\end{enumerate}
\end{exercice}

\begin{corrige}[Exercice : éléments de réponse]
\begin{enumerate}
\item $\cos x = \dfrac{\sqrt{2}}{2}$ pour $x = \dfrac{\pi}{4}$ et $x = \dfrac{7\pi}{4}$. L'arc solution enjambe $0$ : $S = \left[0\,;\,\dfrac{\pi}{4}\right] \cup \left[\dfrac{7\pi}{4}\,;\,2\pi\right]$.
\item $\sin x < \dfrac{\sqrt{3}}{2}$ ; frontières $\dfrac{\pi}{3}$ et $\dfrac{2\pi}{3}$. On garde le complémentaire de l'arc supérieur : $S = \left[0\,;\,\dfrac{\pi}{3}\right[ \cup \left]\dfrac{2\pi}{3}\,;\,2\pi\right]$.
\item Frontières $\dfrac{\pi}{3}$ et $\dfrac{4\pi}{3}$ ; valeurs interdites $\dfrac{\pi}{2}$, $\dfrac{3\pi}{2}$ exclues : $S = \left[0\,;\,\dfrac{\pi}{3}\right[ \cup \left]\dfrac{\pi}{2}\,;\,\dfrac{4\pi}{3}\right[ \cup \left]\dfrac{3\pi}{2}\,;\,2\pi\right]$.
\item Frontières : $\cos x = 0 \iff x \in \left\{\dfrac{\pi}{2}\,;\,\dfrac{3\pi}{2}\right\}$ ; $2\sin x - 1 = 0 \iff \sin x = \dfrac{1}{2} \iff x \in \left\{\dfrac{\pi}{6}\,;\,\dfrac{5\pi}{6}\right\}$.
Tableau de signes :
\begin{center}
\begin{tabular}{|c|ccccccccccc|}\hline
$x$ & $0$ & & $\frac{\pi}{6}$ & & $\frac{\pi}{2}$ & & $\frac{5\pi}{6}$ & & $\frac{3\pi}{2}$ & & $2\pi$ \\ \hline
$\cos x$ & $+$ & $+$ & $+$ & $+$ & $0$ & $-$ & $-$ & $-$ & $0$ & $+$ & $+$ \\ \hline
$2\sin x - 1$ & $-$ & $-$ & $0$ & $+$ & $+$ & $+$ & $0$ & $-$ & $-$ & $-$ & $-$ \\ \hline
$P(x)$ & $-$ & $-$ & $0$ & $+$ & $0$ & $-$ & $0$ & $+$ & $0$ & $-$ & $-$ \\ \hline
\end{tabular}
\end{center}
$P(x) \leq 0$ (inégalité large) : $S = \left[0\,;\,\dfrac{\pi}{6}\right] \cup \left[\dfrac{\pi}{2}\,;\,\dfrac{5\pi}{6}\right] \cup \left[\dfrac{3\pi}{2}\,;\,2\pi\right]$.
\end{enumerate}
\end{corrige}

\newpage

\coursec{Activité d'intégration}

\coursec{Trigonométrie : Formules de transformation, Équations et Inéquations}

\courssub{Rappels : Formules d'addition et de duplication}
\begin{definition}[Formules d'addition]
Pour tous réels $a$ et $b$ :
\begin{align*}
\cos(a+b) &= \cos a\cos b - \sin a\sin b \\
\cos(a-b) &= \cos a\cos b + \sin a\sin b \\
\sin(a+b) &= \sin a\cos b + \cos a\sin b \\
\sin(a-b) &= \sin a\cos b - \cos a\sin b
\end{align*}
\end{definition}

\begin{propriete}[Formules de duplication]
Pour tout réel $x$ :
\begin{align*}
\sin 2x &= 2\sin x\cos x \\
\cos 2x &= \cos^2 x - \sin^2 x = 2\cos^2 x - 1 = 1 - 2\sin^2 x \\
\tan 2x &= \frac{2\tan x}{1-\tan^2 x} \quad (\cos x \neq 0,\ \cos 2x \neq 0)
\end{align*}
\end{propriete}

\begin{exemplebox}[Calcul de $\sin 75^\circ$ et $\cos 75^\circ$]
$75^\circ = 45^\circ + 30^\circ$.
$\sin 75^\circ = \sin 45^\circ\cos 30^\circ + \cos 45^\circ\sin 30^\circ = \frac{\sqrt{2}}{2}\frac{\sqrt{3}}{2} + \frac{\sqrt{2}}{2}\frac{1}{2} = \frac{\sqrt{6}+\sqrt{2}}{4}$.
$\cos 75^\circ = \cos 45^\circ\cos 30^\circ - \sin 45^\circ\sin 30^\circ = \frac{\sqrt{2}}{2}\frac{\sqrt{3}}{2} - \frac{\sqrt{2}}{2}\frac{1}{2} = \frac{\sqrt{6}-\sqrt{2}}{4}$.
\end{exemplebox}

\courssub{Transformer des produits en sommes}
\begin{propriete}[Formules de transformation Produit $\to$ Somme]
Pour tous réels $p$ et $q$ :
\begin{align*}
\cos p \cos q &= \frac{1}{2}[\cos(p+q) + \cos(p-q)] \\
\sin p \sin q &= \frac{1}{2}[\cos(p-q) - \cos(p+q)] \\
\sin p \cos q &= \frac{1}{2}[\sin(p+q) + \sin(p-q)] \\
\cos p \sin q &= \frac{1}{2}[\sin(p+q) - \sin(p-q)]
\end{align*}
\end{propriete}

\begin{methode}[Transformer un produit en somme]
\begin{enumerate}
    \item Identifier le type de produit ($\cos\cos$, $\sin\sin$, $\sin\cos$ ou $\cos\sin$).
    \item Choisir la formule correspondante.
    \item Remplacer $p$ et $q$ par les arguments donnés.
    \item Simplifier l'expression obtenue.
\end{enumerate}
\end{methode}

\begin{exemplebox}[Exprimer $\sin 50^\circ \cos 10^\circ$ en somme]
$\sin 50^\circ \cos 10^\circ = \frac{1}{2}[\sin(50^\circ+10^\circ) + \sin(50^\circ-10^\circ)] = \frac{1}{2}[\sin 60^\circ + \sin 40^\circ] = \frac{1}{2}\left[\frac{\sqrt{3}}{2} + \sin 40^\circ\right]$.
\end{exemplebox}

\courssub{Transformer des sommes en produits}
\begin{propriete}[Formules de transformation Somme $\to$ Produit]
Pour tous réels $p$ et $q$ :
\begin{align*}
\sin p + \sin q &= 2\sin\!\left(\frac{p+q}{2}\right)\cos\!\left(\frac{p-q}{2}\right) \\
\sin p - \sin q &= 2\cos\!\left(\frac{p+q}{2}\right)\sin\!\left(\frac{p-q}{2}\right) \\
\cos p + \cos q &= 2\cos\!\left(\frac{p+q}{2}\right)\cos\!\left(\frac{p-q}{2}\right) \\
\cos p - \cos q &= -2\sin\!\left(\frac{p+q}{2}\right)\sin\!\left(\frac{p-q}{2}\right)
\end{align*}
\end{propriete}

\begin{attention}[Pièges fréquents]
\begin{itemize}
    \item Ne pas confondre l'ordre des termes dans $\cos p - \cos q$ : le signe $-$ devant le $2$ est obligatoire.
    \item Vérifier que les arguments des fonctions trigonométriques sont bien en radians (ou degrés) de manière cohérente.
    \item Pour $\sin p \pm \sin q$, le facteur $\cos\frac{p-q}{2}$ ou $\sin\frac{p-q}{2}$ ne doit pas être oublié.
\end{itemize}
\end{attention}

\begin{methode}[Transformer une somme en produit]
\begin{enumerate}
    \item Identifier la somme : $\sin\pm\sin$ ou $\cos\pm\cos$.
    \item Poser $p$ et $q$ (l'ordre n'a pas d'importance pour la somme, mais change le signe pour la différence).
    \item Appliquer la formule adaptée.
    \item Simplifier les angles $\frac{p+q}{2}$ et $\frac{p-q}{2}$.
\end{enumerate}
\end{methode}

\begin{exemplebox}[Factoriser $\sin 3x + \sin x$]
$\sin 3x + \sin x = 2\sin\!\left(\frac{3x+x}{2}\right)\cos\!\left(\frac{3x-x}{2}\right) = 2\sin 2x \cos x$.
\end{exemplebox}

\courssub{Équations trigonométriques élémentaires}
\begin{definition}[Équations élémentaires]
On appelle équations trigonométriques élémentaires les équations de la forme :
$\cos x = a$, $\sin x = a$, $\tan x = a$ avec $a \in \mathbb{R}$.
\end{definition}

\begin{propriete}[Résolution dans $\mathbb{R}$]
Soit $a \in [-1\,;\, 1]$.
\begin{itemize}
    \item $\cos x = a \Leftrightarrow x = \arccos a + 2k\pi$ ou $x = -\arccos a + 2k\pi$, $k\in\mathbb{Z}$.
    \item $\sin x = a \Leftrightarrow x = \arcsin a + 2k\pi$ ou $x = \pi - \arcsin a + 2k\pi$, $k\in\mathbb{Z}$.
    \item $\tan x = a \Leftrightarrow x = \arctan a + k\pi$, $k\in\mathbb{Z}$.
\end{itemize}
Si $a \notin [-1\,;\, 1]$, $\cos x = a$ et $\sin x = a$ n'ont pas de solution.
\end{propriete}

\begin{methode}[Résoudre une équation élémentaire sur un intervalle $I$]
\begin{enumerate}
    \item Trouver les solutions générales dans $\mathbb{R}$.
    \item Déterminer l'intervalle de variation de l'inconnue (ex: $x \in [0\,;\, 2\pi]$).
    \item Sélectionner les valeurs de $k$ qui donnent des solutions dans $I$.
    \item Présenter l'ensemble solution $S$.
\end{enumerate}
\end{methode}

\begin{exoresolu}[{Résoudre $\cos(2x) = \frac{1}{2}$ sur $[0\,;\, 2\pi]$}]
\underline{Énoncé} : Résoudre $\cos(2x) = \frac{1}{2}$ sur $[0\,;\, 2\pi]$.
\tcblower
\underline{Solution} :
Posons $X = 2x$. Comme $x \in [0\,;\, 2\pi]$, $X \in [0\,;\, 4\pi]$.
$\cos X = \frac{1}{2} \Leftrightarrow X = \frac{\pi}{3} + 2k\pi$ ou $X = -\frac{\pi}{3} + 2k\pi$, $k\in\mathbb{Z}$.
Dans $[0\,;\, 4\pi]$ :
$k=0$ : $X_1 = \frac{\pi}{3}$, $X_2 = -\frac{\pi}{3}$ (rejeté).
$k=1$ : $X_3 = \frac{\pi}{3}+2\pi = \frac{7\pi}{3}$, $X_4 = -\frac{\pi}{3}+2\pi = \frac{5\pi}{3}$.
$k=2$ : $X_5 = \frac{\pi}{3}+4\pi = \frac{13\pi}{3} > 4\pi$ (rejeté).
Donc $X \in \left\{\frac{\pi}{3}, \frac{5\pi}{3}, \frac{7\pi}{3}\right\}$.
$x = \frac{X}{2} \in \left\{\frac{\pi}{6}, \frac{5\pi}{6}, \frac{7\pi}{6}\right\}$.
\boxed{S = \left\{\frac{\pi}{6}, \frac{5\pi}{6}, \frac{7\pi}{6}\right\}}
\end{exoresolu}

\courssub{Équations se ramenant aux formes élémentaires}
\begin{methode}[Stratégies de résolution]
\begin{enumerate}
    \item \textbf{Factorisation :} Mettre sous forme produit = 0. Ex: $(2\cos x - 1)(\sin x + 1) = 0$.
    \item \textbf{Changement de variable :} $\cos^2 x = 1 - \sin^2 x$ ou $t = \tan\frac{x}{2}$ (formules d'Euler).
    \item \textbf{Formules de transformation :} Transformer somme en produit pour factoriser. Ex: $\sin 3x + \sin x = 0 \Rightarrow 2\sin 2x \cos x = 0$.
    \item \textbf{Équation homogène en $\sin x$ et $\cos x$ :} Diviser par $\cos^n x$ (si $\cos x \neq 0$) pour obtenir une équation en $\tan x$.
\end{enumerate}
\end{methode}

\begin{exoresolu}[{Résoudre $\sin 3x + \sin x = 0$ sur $[0\,;\, 2\pi]$}]
\underline{Énoncé} : Résoudre $\sin 3x + \sin x = 0$ sur $[0\,;\, 2\pi]$.
\tcblower
\underline{Solution} :
Transformation somme en produit :
$\sin 3x + \sin x = 2\sin\!\left(\frac{3x+x}{2}\right)\cos\!\left(\frac{3x-x}{2}\right) = 2\sin 2x \cos x$.
L'équation devient $2\sin 2x \cos x = 0 \Leftrightarrow \sin 2x = 0$ ou $\cos x = 0$.
\begin{itemize}
    \item $\sin 2x = 0 \Leftrightarrow 2x = k\pi \Leftrightarrow x = \frac{k\pi}{2}$, $k\in\mathbb{Z}$.
    Sur $[0\,;\, 2\pi]$ : $x \in \{0, \frac{\pi}{2}, \pi, \frac{3\pi}{2}, 2\pi\}$.
    \item $\cos x = 0 \Leftrightarrow x = \frac{\pi}{2} + k\pi$, $k\in\mathbb{Z}$.
    Sur $[0\,;\, 2\pi]$ : $x \in \{\frac{\pi}{2}, \frac{3\pi}{2}\}$.
\end{itemize}
Réunion : $S = \{0, \frac{\pi}{2}, \pi, \frac{3\pi}{2}, 2\pi\}$.
\end{exoresolu}

\courssub{Inéquations trigonométriques}
\begin{methode}[Résoudre $\sin x \geq m$ ou $\cos x \geq m$ sur un intervalle]
\begin{enumerate}
    \item Résoudre l'équation associée $\sin x = m$ (ou $\cos x = m$) pour trouver les bornes.
    \item Utiliser le \textbf{cercle trigonométrique} ou le \textbf{tableau de variations} de la fonction sur une période ($[-\frac{\pi}{2}\,;\, \frac{3\pi}{2}]$ pour $\sin$, $[0\,;\, 2\pi]$ pour $\cos$).
    \item Repérer les intervalles où l'inégalité est vérifiée.
    \item Tenir compte de la période $2\pi$ et de l'intervalle imposé par l'énoncé.
\end{enumerate}
\end{methode}

\begin{popfigure}
\begin{tikzpicture}[scale=1.2]
    % Circle
    \draw[popline, thick] (0,0) circle (2cm);
    % Axes
    \draw[->, popmuted] (-2.3,0) -- (2.3,0) node[right] {$x$};
    \draw[->, popmuted] (0,-2.3) -- (0,2.3) node[above] {$y$};
    % Angles pi/4 and 3pi/4
    \draw[popblue, thick] (0,0) -- (1.414,1.414) node[midway, above right] {$r$};
    \draw[popblue, thick] (0,0) -- (-1.414,1.414) node[midway, above left] {$r$};
    % Arc for solution
    \draw[popgreen, very thick, decoration={markings, mark=at position 0.5 with {\arrow{>}}}, postaction={decorate}] (45:2) arc (45:135:2) node[midway, left=3mm] {$\sin x \geq \frac{\sqrt{2}}{2}$};
    % Horizontal line y = sqrt(2)/2
    \draw[poporange, dashed] (-2,1.414) -- (2,1.414) node[right] {$y=\frac{\sqrt{2}}{2}$};
    % Points
    \fill[popblue] (1.414,1.414) circle (2pt) node[above right] {$A(\frac{\pi}{4})$};
    \fill[popblue] (-1.414,1.414) circle (2pt) node[above left] {$B(\frac{3\pi}{4})$};
    % Angle labels
    \draw (0.5,0) arc (0:45:0.5) node[midway, right=2mm] {$\frac{\pi}{4}$};
    \draw (0.5,0) arc (0:135:0.7) node[midway, left=2mm] {$\frac{3\pi}{4}$};
\end{tikzpicture}
\legende{Résolution de $\sin x \geq \frac{\sqrt{2}}{2}$ sur le cercle trigonométrique : $x \in [\frac{\pi}{4}\,;\, \frac{3\pi}{4}] + 2k\pi$.}
\end{popfigure}

\begin{attention}[Points de vigilance pour les inéquations]
\begin{itemize}
    \item Ne pas oublier la période $2\pi$ pour l'ensemble des solutions dans $\mathbb{R}$.
    \item Sur un intervalle borné (ex: $[0\,;\, 2\pi]$), ne donner que les solutions incluses dans cet intervalle.
    \item Vérifier le sens de l'inégalité (croissant/décroissant) sur le tableau de variations.
    \item Attention au changement de variable $X = \omega t + \varphi$ : reporter l'intervalle de $t$ sur $X$ avant de résoudre.
\end{itemize}
\end{attention}

\courssub{Activité d'intégration : Couplage de deux groupes électrogènes}
\begin{aretenir}[Objectifs de l'activité]
Mobiliser l'ensemble des savoir-faire de la leçon (formules de transformation, résolution d'équations et d'inéquations) dans une situation professionnelle authentique : l'analyse de la tension électrique issue du couplage de deux groupes électrogènes. L'élève devra modéliser, calculer, résoudre et interpréter ses résultats pour émettre un avis technique argumenté.
\end{aretenir}

\begin{exemplebox}[Situation professionnelle]
Dans un quartier périphérique de Douala, une petite entreprise de menuiserie-aluminium fonctionne grâce à deux groupes électrogènes mis en parallèle pour assurer une puissance suffisante aux machines à souder et aux compresseurs. Le technicien de maintenance doit vérifier la qualité de la tension délivrée au tableau général.

La tension (en volts) délivrée par le premier groupe est modélisée par la fonction :
\[ u(t) = 220\sqrt{2}\,\sin(100\pi t) \]
La tension délivrée par le second groupe, déphasée par rapport au premier, est modélisée par :
\[ v(t) = 220\sqrt{2}\,\sin\!\left(100\pi t + \frac{\pi}{3}\right) \]
où $t$ est le temps en secondes. La pulsation $\omega = 100\pi$ rad/s correspond à une fréquence de 50 Hz (période $T = 0,02$ s). La tension efficace nominale du réseau est de 220 V ; la tension instantanée maximale (amplitude) est donc $220\sqrt{2} \approx 311$ V.

Un appareil de commande numérique (CNC) sensible, branché sur ce réseau, nécessite une tension instantanée $\geq 220$ V pour fonctionner correctement sans risque de plantage. Le technicien doit déterminer les caractéristiques de la tension résultante $w(t) = u(t) + v(t)$ et les créneaux temporels utilisables sur une période.
\end{exemplebox}

\begin{methode}[Consignes]
\begin{enumerate}
    \item \textbf{Expression de la tension totale.} En utilisant la formule de transformation \emph{somme en produit} (ou \emph{somme de sinus}), écrire $w(t) = u(t) + v(t)$ sous la forme canonique $A\sin(100\pi t + \varphi)$. Préciser l'amplitude $A$ (en volts) et le déphasage $\varphi$ (en radians).
    \item \textbf{Instants d'annulation de la tension.} Résoudre l'équation $w(t) = 0$ sur l'intervalle $[0\,;\, 0,02]$ (une période). Interpréter physiquement ces instants.
    \item \textbf{Durées de tension suffisante.} Résoudre l'inéquation $u(t) \geq 220$ sur $[0\,;\, 0,02]$. En déduire la durée totale, sur une période, pendant laquelle la tension du premier groupe seul est suffisante pour alimenter l'appareil CNC.
    \item \textbf{Conclusion technique.} Rédiger un court rapport (5 à 7 lignes) à l'attention du chef d'atelier : la mise en parallèle des deux groupes améliore-t-elle l'amplitude ? Quels sont les créneaux de fonctionnement sûr pour l'appareil sensible sur la tension $u(t)$ ? Quelle recommandation formuler ?
\end{enumerate}
\end{methode}

\begin{aretenir}[Ressources à mobiliser]
\begin{itemize}
    \item \textbf{Formule de transformation somme en produit :} $\sin p + \sin q = 2\sin\!\left(\frac{p+q}{2}\right)\cos\!\left(\frac{p-q}{2}\right)$.
    \item \textbf{Forme canonique :} $a\sin x + b\cos x = A\sin(x+\varphi)$ avec $A = \sqrt{a^2+b^2}$, $\cos\varphi = \frac{a}{A}$, $\sin\varphi = \frac{b}{A}$.
    \item \textbf{Valeurs remarquables :} $\sin\frac{\pi}{3} = \frac{\sqrt{3}}{2}$, $\cos\frac{\pi}{3} = \frac{1}{2}$, $\cos\frac{\pi}{6} = \frac{\sqrt{3}}{2}$.
    \item \textbf{Résolution de $\sin X = m$ :} $X = \arcsin m + 2k\pi$ ou $X = \pi - \arcsin m + 2k\pi$, $k\in\mathbb{Z}$.
    \item \textbf{Résolution de $\sin X \geq m$ :} Utiliser le cercle trigonométrique ou le tableau de variations de la fonction sinus sur $[-\frac{\pi}{2}\,;\, \frac{3\pi}{2}]$.
\end{itemize}
\end{aretenir}

\begin{popfigure}
\begin{tikzpicture}
\begin{axis}[
    width=\linewidth, height=6cm,
    axis lines=middle,
    xlabel={$t$ (s)}, ylabel={$u, v, w$ (V)},
    domain=0:0.02, samples=400,
    xmin=0, xmax=0.022, ymin=-600, ymax=600,
    xtick={0,0.005,0.01,0.015,0.02},
    xticklabels={$0$, $5$ ms, $10$ ms, $15$ ms, $20$ ms},
    ytick={-538.8,-311,0,220,311,538.8},
    yticklabels={$-539$, $-311$, $0$, $220$, $311$, $539$},
    legend style={at={(0.98,0.98)}, anchor=north east, fill=popcream, draw=popline},
    grid=major, grid style={popline!30},
]
\addplot[popblue, thick] {220*sqrt(max(0,2))*sin(deg(100*pi*x))};
\addlegendentry{$u(t)$}
\addplot[poporange, thick] {220*sqrt(max(0,2))*sin(deg(100*pi*x + pi/3))};
\addlegendentry{$v(t)$}
\addplot[popgreen, very thick] {220*sqrt(max(0,6))*sin(deg(100*pi*x + pi/6))};
\addlegendentry{$w(t)$}
\addplot[poppink, dashed, thick] coordinates {(0,220) (0.02,220)};
\addlegendentry{Seuil 220 V}
\end{axis}
\end{tikzpicture}
\legende{Signaux de tension $u(t)$, $v(t)$ et résultante $w(t)$ sur une période (20 ms).}
\end{popfigure}

\begin{exoresolu}[Corrigé détaillé de l'activité d'intégration]
\underline{\textbf{1. Expression de la tension totale $w(t)$}}

On a $w(t) = 220\sqrt{2}\left[ \sin(100\pi t) + \sin\!\left(100\pi t + \frac{\pi}{3}\right) \right]$.
Posons $p = 100\pi t + \frac{\pi}{3}$ et $q = 100\pi t$.
On applique la formule $\sin p + \sin q = 2\sin\!\left(\frac{p+q}{2}\right)\cos\!\left(\frac{p-q}{2}\right)$ :
\begin{align*}
\frac{p+q}{2} &= \frac{200\pi t + \frac{\pi}{3}}{2} = 100\pi t + \frac{\pi}{6} \\
\frac{p-q}{2} &= \frac{\frac{\pi}{3}}{2} = \frac{\pi}{6}
\end{align*}
Donc :
\[ \sin(100\pi t) + \sin\!\left(100\pi t + \frac{\pi}{3}\right) = 2\sin\!\left(100\pi t + \frac{\pi}{6}\right)\cos\!\left(\frac{\pi}{6}\right) \]
Or $\cos\frac{\pi}{6} = \frac{\sqrt{3}}{2}$. Ainsi :
\[ w(t) = 220\sqrt{2} \times 2 \times \frac{\sqrt{3}}{2} \sin\!\left(100\pi t + \frac{\pi}{6}\right) = 220\sqrt{6}\,\sin\!\left(100\pi t + \frac{\pi}{6}\right) \]
\textbf{Résultat :} La tension totale s'écrit sous la forme $w(t) = A\sin(100\pi t + \varphi)$ avec :
\[ \boxed{A = 220\sqrt{6} \approx 538,8 \text{ V}} \quad \text{et} \quad \boxed{\varphi = \frac{\pi}{6} \text{ rad}} \]
\emph{Commentaire :} L'amplitude $220\sqrt{6}$ est supérieure à l'amplitude de chaque source prise séparément ($220\sqrt{2} \approx 311$ V). Le couplage en parallèle avec un déphasage de $\pi/3$ augmente la tension de crête d'un facteur $\sqrt{3}$.

\bigskip
\underline{\textbf{2. Instants d'annulation de la tension totale sur $[0\,;\, 0,02]$}}

On résout $w(t) = 0 \Leftrightarrow 220\sqrt{6}\,\sin\!\left(100\pi t + \frac{\pi}{6}\right) = 0$.
Comme $220\sqrt{6} \neq 0$, cela revient à $\sin\!\left(100\pi t + \frac{\pi}{6}\right) = 0$.
Le sinus s'annule pour $X = k\pi$, $k\in\mathbb{Z}$.
Posons $X = 100\pi t + \frac{\pi}{6}$. On a :
\[ 100\pi t + \frac{\pi}{6} = k\pi \quad \Leftrightarrow \quad 100\pi t = k\pi - \frac{\pi}{6} \quad \Leftrightarrow \quad t = \frac{k}{100} - \frac{1}{600} \]
On cherche $t \in [0\,;\, 0,02] = [0\,;\, \frac{1}{50}] = [0\,;\, \frac{2}{100}]$.
\begin{itemize}
    \item $k = 0$ : $t = -\frac{1}{600} < 0$ (rejeté).
    \item $k = 1$ : $t = \frac{1}{100} - \frac{1}{600} = \frac{6-1}{600} = \frac{5}{600} = \frac{1}{120} \approx 0,00833$ s. \textbf{Accepté}.
    \item $k = 2$ : $t = \frac{2}{100} - \frac{1}{600} = \frac{12-1}{600} = \frac{11}{600} \approx 0,01833$ s. \textbf{Accepté}.
    \item $k = 3$ : $t = \frac{3}{100} - \frac{1}{600} = \frac{17}{600} > 0,02$ (rejeté).
\end{itemize}
\boxed{\text{La tension totale s'annule aux instants } t_1 = \frac{1}{120} \text{ s et } t_2 = \frac{11}{600} \text{ s sur la période.}}
\emph{Interprétation :} Ce sont les passages par zéro de la tension alternative résultante (zéros croissants ou décroissants selon la dérivée).

\bigskip
\underline{\textbf{3. Durée où $u(t) \geq 220$ V sur $[0\,;\, 0,02]$}}

On résout $u(t) \geq 220 \Leftrightarrow 220\sqrt{2}\,\sin(100\pi t) \geq 220$.
On divise par $220$ (positif) : $\sqrt{2}\,\sin(100\pi t) \geq 1 \Leftrightarrow \sin(100\pi t) \geq \frac{1}{\sqrt{2}} = \frac{\sqrt{2}}{2}$.
Posons $X = 100\pi t$. Comme $t \in [0\,;\, 0,02]$, alors $X \in [0\,;\, 2\pi]$.
On résout $\sin X \geq \frac{\sqrt{2}}{2}$ sur $[0\,;\, 2\pi]$.
Sur le cercle trigonométrique, $\sin X = \frac{\sqrt{2}}{2}$ pour $X = \frac{\pi}{4}$ et $X = \pi - \frac{\pi}{4} = \frac{3\pi}{4}$.
Le sinus est $\geq \frac{\sqrt{2}}{2}$ sur l'intervalle $\left[ \frac{\pi}{4}\,;\, \frac{3\pi}{4} \right]$.
Donc : $\frac{\pi}{4} \leq 100\pi t \leq \frac{3\pi}{4}$.
On divise par $100\pi$ (positif) :
\[ \frac{1}{400} \leq t \leq \frac{3}{400} \]
Soit $0,0025 \leq t \leq 0,0075$.
\boxed{\text{L'inéquation est vérifiée pour } t \in \left[ \frac{1}{400}\,;\, \frac{3}{400} \right] \text{ soit } [0,0025\,;\, 0,0075] \text{ s.}}
\emph{Durée totale :} $\Delta t = \frac{3}{400} - \frac{1}{400} = \frac{2}{400} = \frac{1}{200} = 0,005$ s = \textbf{5 ms}.

\bigskip
\underline{\textbf{4. Conclusion technique (Rapport au chef d'atelier)}}

\begin{quote}
\textbf{Rapport de maintenance - Tension groupe parallèle}\\
Monsieur le Chef d'atelier,\\
L'analyse de la tension résultante $w(t)$ montre une amplitude de crête de $539$ V ($220\sqrt{6}$), soit un gain de $73\%$ par rapport à un groupe seul ($311$ V). Cependant, l'appareil CNC nécessite $u(t) \geq 220$ V. Sur la seule source $u(t)$, cette condition n'est remplie que pendant $5$ ms par période de $20$ ms (soit $25\%$ du temps), sur le créneau $[2,5\,;\, 7,5]$ ms. La tension totale $w(t)$ étant déphasée de $\pi/6$, ses créneaux utiles sont décalés.\\
\textbf{Recommandation :} Ne pas brancher l'appareil sensible directement sur $u(t)$ ni sur $w(t)$ sans régulation. Installer un onduleur (ASI) ou un stabilisateur de tension pour garantir une tension continue $\geq 220$ V efficace.
\end{quote}
\end{exoresolu}

\begin{aretenir}[Bilan de l'activité]
Cette activité a permis de mettre en évidence la puissance des formules de transformation pour simplifier l'analyse de signaux physiques réels (électricité alternative). Les élèves ont réinvesti :
\begin{itemize}
    \item La transformation d'une somme de sinus en produit pour retrouver l'amplitude et la phase d'un signal résultant (compétence de modélisation).
    \item La résolution rigoureuse d'une équation trigonométrique $\sin X = 0$ en gérant le changement de variable $X = \omega t + \varphi$ et la sélection des solutions dans un intervalle borné (compétence algébrique).
    \item La résolution d'une inéquation $\sin X \geq m$ via le cercle trigonométrique, et l'interprétation concrète du résultat en \emph{durée} (compétence d'analyse fonctionnelle).
    \item La rédaction d'une conclusion technique structurée, passant du calcul mathématique à la décision professionnelle (compétence de communication).
\end{itemize}
L'ancrage dans le contexte camerounais (groupes électrogènes, quartiers périphériques, matériel CNC) donne du sens aux abstractions trigonométriques et prépare aux épreuves du Probatoire technique où les situations professionnelles sont centrales.
\end{aretenir}

\newpage

\coursec{Méthodes et exercices résolus}

\coursec{Trigonométrie : Formules de transformation, Équations et Inéquations}

\begin{aretenir}[Rappels : formules d'addition et de duplication]
Les formules suivantes sont admises et servent de base aux transformations :
\begin{align*}
\cos(a+b) &= \cos a\cos b - \sin a\sin b \\
\cos(a-b) &= \cos a\cos b + \sin a\sin b \\
\sin(a+b) &= \sin a\cos b + \cos a\sin b \\
\sin(a-b) &= \sin a\cos b - \cos a\sin b
\end{align*}
\medskip
\noindent \textbf{Formules de duplication} (cas particulier $a=b$) :
\[ \cos 2a = \cos^2 a - \sin^2 a = 2\cos^2 a - 1 = 1 - 2\sin^2 a \qquad ; \qquad \sin 2a = 2\sin a\cos a \]
\end{aretenir}

\courssub{Transformer des produits en sommes}

\begin{methode}[Linéariser un produit de cosinus et de sinus]
Pour transformer un produit en somme, on utilise les trois \cle{formules de transformation produit-somme} (démonstration exigible au Probatoire par addition/soustraction des formules d'addition) :
\[ \cos a\,\cos b = \tfrac{1}{2}\big[\cos(a+b)+\cos(a-b)\big] \]
\[ \sin a\,\sin b = \tfrac{1}{2}\big[\cos(a-b)-\cos(a+b)\big] \]
\[ \sin a\,\cos b = \tfrac{1}{2}\big[\sin(a+b)+\sin(a-b)\big] \]
Démarche :
\begin{enumerate}
\item Identifier la nature du produit ($\cos\cos$, $\sin\sin$ ou $\sin\cos$).
\item Écrire la formule correspondante en faisant attention au signe entre les deux termes.
\item Calculer $a+b$ et $a-b$, puis simplifier.
\item Vérifier le facteur $\tfrac{1}{2}$ devant la somme.
\end{enumerate}
\end{methode}

\begin{exoresolu}[Linéarisation d'un produit]
Énoncé : transformer en somme les expressions suivantes :
\[ A = \cos 3x\,\cos x \qquad ; \qquad B = \sin 5x\,\sin 2x \qquad ; \qquad C = 2\sin 4x\,\cos x. \]
\tcblower
\textbf{Solution détaillée.}

\textbf{Pour $A$ :} c'est un produit $\cos a\,\cos b$ avec $a=3x$ et $b=x$. On applique :
\begin{align*}
A &= \tfrac{1}{2}\big[\cos(3x+x)+\cos(3x-x)\big] \\
  &= \tfrac{1}{2}\big[\cos 4x + \cos 2x\big].
\end{align*}

\textbf{Pour $B$ :} c'est un produit $\sin a\,\sin b$ avec $a=5x$ et $b=2x$. Attention au signe moins devant $\cos(a+b)$ :
\begin{align*}
B &= \tfrac{1}{2}\big[\cos(5x-2x)-\cos(5x+2x)\big] \\
  &= \tfrac{1}{2}\big[\cos 3x - \cos 7x\big].
\end{align*}

\textbf{Pour $C$ :} on commence par appliquer la formule $\sin a\,\cos b$ à $\sin 4x\,\cos x$ :
\begin{align*}
\sin 4x\,\cos x &= \tfrac{1}{2}\big[\sin(4x+x)+\sin(4x-x)\big] = \tfrac{1}{2}\big[\sin 5x + \sin 3x\big].
\end{align*}
Donc $C = 2\sin 4x\,\cos x = \sin 5x + \sin 3x$.

\textbf{Conclusion :} $A = \tfrac{1}{2}(\cos 4x + \cos 2x)$, $B = \tfrac{1}{2}(\cos 3x - \cos 7x)$ et $C = \sin 5x + \sin 3x$.
\end{exoresolu}

\courssub{Transformer des sommes en produits}

\begin{methode}[Factoriser une somme de cosinus ou de sinus]
On utilise les \cle{formules somme-produit}, obtenues en posant $p = a+b$ et $q = a-b$ dans les formules produit-somme :
\[ \cos p + \cos q = 2\cos\frac{p+q}{2}\,\cos\frac{p-q}{2} \qquad ; \qquad \cos p - \cos q = -2\sin\frac{p+q}{2}\,\sin\frac{p-q}{2} \]
\[ \sin p + \sin q = 2\sin\frac{p+q}{2}\,\cos\frac{p-q}{2} \qquad ; \qquad \sin p - \sin q = 2\cos\frac{p+q}{2}\,\sin\frac{p-q}{2} \]
Démarche :
\begin{enumerate}
\item Vérifier que la somme porte sur deux fonctions de même nature (deux cosinus ou deux sinus).
\item Calculer les demi-somme $\dfrac{p+q}{2}$ et demi-différence $\dfrac{p-q}{2}$.
\item Écrire le produit avec le facteur $2$ (ou $-2$ pour $\cos p - \cos q$).
\item Simplifier les angles obtenus.
\end{enumerate}
Réflexe : factoriser une somme permet de résoudre des équations par \emph{produit nul}.
\end{methode}

\begin{exoresolu}[Factorisation d'une somme]
Énoncé : factoriser les expressions suivantes :
\[ D = \cos 5x + \cos x \qquad ; \qquad E = \sin 3x - \sin x. \]
\tcblower
\textbf{Solution détaillée.}

\textbf{Pour $D$ :} somme de deux cosinus avec $p=5x$ et $q=x$. On calcule :
\[ \frac{p+q}{2} = \frac{5x+x}{2} = 3x \qquad ; \qquad \frac{p-q}{2} = \frac{5x-x}{2} = 2x. \]
On applique la formule :
\[ D = 2\cos 3x\,\cos 2x. \]

\textbf{Pour $E$ :} différence de deux sinus avec $p=3x$ et $q=x$. On calcule :
\[ \frac{p+q}{2} = \frac{3x+x}{2} = 2x \qquad ; \qquad \frac{p-q}{2} = \frac{3x-x}{2} = x. \]
On applique la formule $\sin p - \sin q = 2\cos\frac{p+q}{2}\,\sin\frac{p-q}{2}$ :
\[ E = 2\cos 2x\,\sin x. \]

\textbf{Conclusion :} $D = 2\cos 3x\,\cos 2x$ et $E = 2\cos 2x\,\sin x$. On peut vérifier mentalement : pour $x=0$, $D = 2$ et $2\cos 0\,\cos 0 = 2$ ; pour $x = \tfrac{\pi}{2}$, $E = \sin\frac{3\pi}{2} - \sin\frac{\pi}{2} = -1-1 = -2$ et $2\cos\pi\,\sin\frac{\pi}{2} = -2$. Les égalités sont cohérentes.
\end{exoresolu}

\courssub{Équations trigonométriques élémentaires}

\begin{methode}[Équations $\cos x = a$ et $\sin x = a$]
\begin{enumerate}
\item Vérifier la \cle{condition d'existence} : l'équation n'a de solution que si $-1 \leq a \leq 1$.
\item Trouver une solution particulière $\alpha$ (cercle trigonométrique ou valeurs remarquables) telle que $\cos\alpha = a$ (ou $\sin\alpha = a$).
\item Écrire l'ensemble des solutions :
\begin{itemize}
\item $\cos x = \cos\alpha \iff x = \alpha + 2k\pi$ \textbf{ou} $x = -\alpha + 2k\pi$, $k\in\mathbb{Z}$ ;
\item $\sin x = \sin\alpha \iff x = \alpha + 2k\pi$ \textbf{ou} $x = \pi - \alpha + 2k\pi$, $k\in\mathbb{Z}$.
\end{itemize}
\item Si l'on demande les solutions dans un intervalle donné (souvent $[0\,;\,2\pi[$), choisir les valeurs de $k$ convenables et lister toutes les solutions.
\end{enumerate}
Cas particuliers : $\cos x = 0 \iff x = \frac{\pi}{2} + k\pi$ ; $\sin x = 0 \iff x = k\pi$ ; $\tan x = \tan\alpha \iff x = \alpha + k\pi$.
\end{methode}

\begin{exoresolu}[Équation élémentaire avec cosinus]
Énoncé : résoudre dans $\mathbb{R}$ l'équation $\cos x = \dfrac{1}{2}$, puis donner les solutions dans $[0\,;\,2\pi[$.
\tcblower
\textbf{Solution détaillée.}

\textbf{Étape 1 :} $\dfrac{1}{2}\in[-1\,;\,1]$, l'équation admet des solutions.

\textbf{Étape 2 :} on connaît la valeur remarquable $\cos\dfrac{\pi}{3} = \dfrac{1}{2}$. Donc $\alpha = \dfrac{\pi}{3}$ est une solution particulière.

\textbf{Étape 3 :} l'équation s'écrit $\cos x = \cos\dfrac{\pi}{3}$, d'où :
\[ x = \frac{\pi}{3} + 2k\pi \quad \text{ou} \quad x = -\frac{\pi}{3} + 2k\pi, \quad k\in\mathbb{Z}. \]

\textbf{Étape 4 :} solutions dans $[0\,;\,2\pi[$ :
\begin{itemize}
\item Pour $x = \frac{\pi}{3} + 2k\pi$ : $k=0$ donne $x = \frac{\pi}{3}$ (convenable) ; $k=1$ donne $\frac{7\pi}{3} > 2\pi$ (rejeté).
\item Pour $x = -\frac{\pi}{3} + 2k\pi$ : $k=1$ donne $x = -\frac{\pi}{3} + 2\pi = \frac{5\pi}{3}$ (convenable) ; $k=0$ donne $-\frac{\pi}{3} < 0$ (rejeté).
\end{itemize}

\textbf{Conclusion :} dans $\mathbb{R}$, $S = \left\{\frac{\pi}{3} + 2k\pi \,;\, -\frac{\pi}{3} + 2k\pi \,,\ k\in\mathbb{Z}\right\}$ ; dans $[0\,;\,2\pi[$, $S = \left\{\dfrac{\pi}{3}\,;\,\dfrac{5\pi}{3}\right\}$.
\end{exoresolu}

\begin{exoresolu}[Équation élémentaire avec sinus]
Énoncé : résoudre dans $[0\,;\,2\pi]$ l'équation $\sin x = -\dfrac{\sqrt{2}}{2}$.
\tcblower
\textbf{Solution détaillée.}

\textbf{Étape 1 :} $-\dfrac{\sqrt{2}}{2}\in[-1\,;\,1]$ : il y a des solutions.

\textbf{Étape 2 :} on sait que $\sin\dfrac{\pi}{4} = \dfrac{\sqrt{2}}{2}$, donc $\sin\left(-\dfrac{\pi}{4}\right) = -\dfrac{\sqrt{2}}{2}$ (le sinus est impair). Une solution particulière est $\alpha = -\dfrac{\pi}{4}$.

\textbf{Étape 3 :} $\sin x = \sin\left(-\frac{\pi}{4}\right)$ équivaut à :
\[ x = -\frac{\pi}{4} + 2k\pi \quad \text{ou} \quad x = \pi - \left(-\frac{\pi}{4}\right) + 2k\pi = \frac{5\pi}{4} + 2k\pi, \quad k\in\mathbb{Z}. \]

\textbf{Étape 4 :} dans $[0\,;\,2\pi]$ :
\begin{itemize}
\item $x = -\frac{\pi}{4} + 2k\pi$ : $k=1$ donne $x = \frac{7\pi}{4}$ ;
\item $x = \frac{5\pi}{4} + 2k\pi$ : $k=0$ donne $x = \frac{5\pi}{4}$.
\end{itemize}
Vérification : $\sin\frac{5\pi}{4} = -\frac{\sqrt{2}}{2}$ et $\sin\frac{7\pi}{4} = -\frac{\sqrt{2}}{2}$. Correct.

\textbf{Conclusion :} $S = \left\{\dfrac{5\pi}{4}\,;\,\dfrac{7\pi}{4}\right\}$.
\end{exoresolu}

\courssub{Équations se ramenant aux formes élémentaires}

\begin{methode}[Équations avec un angle composé ou factorisables]
\textbf{Cas 1 : angle composé} (type $\cos(2x) = a$ ou $\sin\left(x - \frac{\pi}{6}\right) = a$) :
\begin{enumerate}
\item Poser $X = $ l'angle composé et résoudre l'équation élémentaire en $X$.
\item Remplacer $X$ par son expression et isoler $x$ (diviser tous les termes, y compris $2k\pi$).
\item Sélectionner les solutions dans l'intervalle demandé en testant plusieurs valeurs de $k$.
\end{enumerate}
\textbf{Cas 2 : équation factorisable} (type $\sin 3x = \sin x$ ou somme nulle) :
\begin{enumerate}
\item Utiliser $\sin A = \sin B \iff A = B + 2k\pi$ ou $A = \pi - B + 2k\pi$, ou transformer la somme en produit.
\item Se ramener à un \cle{produit nul} : $P\times Q = 0 \iff P = 0$ ou $Q = 0$.
\item Résoudre chaque équation élémentaire obtenue.
\end{enumerate}
\end{methode}

\begin{exoresolu}[Équation à angle composé]
Énoncé : résoudre dans $\mathbb{R}$ l'équation $\cos\left(2x - \dfrac{\pi}{4}\right) = \dfrac{\sqrt{3}}{2}$, puis donner les solutions dans $[0\,;\,2\pi[$.
\tcblower
\textbf{Solution détaillée.}

\textbf{Étape 1 :} posons $X = 2x - \dfrac{\pi}{4}$. L'équation devient $\cos X = \dfrac{\sqrt{3}}{2}$. Or $\cos\dfrac{\pi}{6} = \dfrac{\sqrt{3}}{2}$, donc :
\[ X = \frac{\pi}{6} + 2k\pi \quad \text{ou} \quad X = -\frac{\pi}{6} + 2k\pi, \quad k\in\mathbb{Z}. \]

\textbf{Étape 2 :} on revient à $x$ :
\begin{align*}
2x - \frac{\pi}{4} = \frac{\pi}{6} + 2k\pi &\iff 2x = \frac{\pi}{6} + \frac{\pi}{4} + 2k\pi = \frac{5\pi}{12} + 2k\pi \\
&\iff x = \frac{5\pi}{24} + k\pi.
\end{align*}
De même :
\begin{align*}
2x - \frac{\pi}{4} = -\frac{\pi}{6} + 2k\pi &\iff 2x = -\frac{\pi}{6} + \frac{\pi}{4} + 2k\pi = \frac{\pi}{12} + 2k\pi \\
&\iff x = \frac{\pi}{24} + k\pi.
\end{align*}

\textbf{Étape 3 :} solutions dans $[0\,;\,2\pi[$ :
\begin{itemize}
\item $x = \frac{\pi}{24} + k\pi$ : $k=0$ donne $\frac{\pi}{24}$ ; $k=1$ donne $\frac{25\pi}{24}$.
\item $x = \frac{5\pi}{24} + k\pi$ : $k=0$ donne $\frac{5\pi}{24}$ ; $k=1$ donne $\frac{29\pi}{24}$.
\end{itemize}
(Les valeurs $k=2$ dépassent $2\pi$ et $k=-1$ donne des valeurs négatives.)

\textbf{Conclusion :} dans $\mathbb{R}$, $S = \left\{\frac{\pi}{24} + k\pi \,;\, \frac{5\pi}{24} + k\pi \,,\ k\in\mathbb{Z}\right\}$ ; dans $[0\,;\,2\pi[$, $S = \left\{\dfrac{\pi}{24}\,;\,\dfrac{5\pi}{24}\,;\,\dfrac{25\pi}{24}\,;\,\dfrac{29\pi}{24}\right\}$.
\end{exoresolu}

\begin{exoresolu}[Équation du type $\sin A = \sin B$]
Énoncé : résoudre dans $[0\,;\,2\pi[$ l'équation $\sin 3x = \sin x$.
\tcblower
\textbf{Solution détaillée.}

On applique la propriété : $\sin A = \sin B \iff A = B + 2k\pi$ ou $A = \pi - B + 2k\pi$, $k\in\mathbb{Z}$.

\textbf{Premier cas :}
\begin{align*}
3x = x + 2k\pi &\iff 2x = 2k\pi \iff x = k\pi.
\end{align*}
Dans $[0\,;\,2\pi[$ : $k=0$ donne $x=0$ ; $k=1$ donne $x=\pi$ ; $k=2$ donne $2\pi$ (exclu).

\textbf{Deuxième cas :}
\begin{align*}
3x = \pi - x + 2k\pi &\iff 4x = \pi + 2k\pi \iff x = \frac{\pi}{4} + \frac{k\pi}{2}.
\end{align*}
Dans $[0\,;\,2\pi[$ :
\begin{itemize}
\item $k=0$ : $x = \frac{\pi}{4}$ ;
\item $k=1$ : $x = \frac{\pi}{4} + \frac{\pi}{2} = \frac{3\pi}{4}$ ;
\item $k=2$ : $x = \frac{\pi}{4} + \pi = \frac{5\pi}{4}$ ;
\item $k=3$ : $x = \frac{\pi}{4} + \frac{3\pi}{2} = \frac{7\pi}{4}$ ;
\item $k=4$ : $x = \frac{9\pi}{4} > 2\pi$ (rejeté).
\end{itemize}

\textbf{Vérification} pour $x = \frac{3\pi}{4}$ : $\sin\frac{9\pi}{4} = \sin\frac{\pi}{4} = \frac{\sqrt{2}}{2}$ et $\sin\frac{3\pi}{4} = \frac{\sqrt{2}}{2}$. Correct.

\textbf{Conclusion :} $S = \left\{0\,;\,\dfrac{\pi}{4}\,;\,\dfrac{3\pi}{4}\,;\,\pi\,;\,\dfrac{5\pi}{4}\,;\,\dfrac{7\pi}{4}\right\}$.
\end{exoresolu}

\courssub{Inéquations trigonométriques}

\begin{methode}[Inéquations $\cos x \geq a$ ou $\sin x < a$ sur un intervalle]
\begin{enumerate}
\item Résoudre d'abord l'\textbf{équation} associée ($\cos x = a$) dans l'intervalle donné : ce sont les \cle{bornes} des zones solutions.
\item Placer ces bornes sur le \cle{cercle trigonométrique}.
\item Lire graphiquement la zone solution :
\begin{itemize}
\item pour $\cos x \geq a$ : les points du cercle dont l'abscisse est supérieure ou égale à $a$ (zone à droite de la droite verticale d'abscisse $a$) ;
\item pour $\sin x \geq a$ : les points dont l'ordonnée est supérieure ou égale à $a$ (zone au-dessus de la droite horizontale d'ordonnée $a$).
\end{itemize}
\item Écrire la solution sous forme d'intervalle(s) en respectant les inégalités strictes ou larges (bornes incluses ou exclues).
\end{enumerate}
\end{methode}

\begin{exoresolu}[Inéquation avec cosinus]
Énoncé : résoudre dans $[0\,;\,2\pi]$ l'inéquation $\cos x \geq \dfrac{1}{2}$.
\tcblower
\textbf{Solution détaillée.}

\textbf{Étape 1 :} on résout l'équation associée $\cos x = \dfrac{1}{2}$ dans $[0\,;\,2\pi]$ : d'après un exercice précédent, $x = \dfrac{\pi}{3}$ ou $x = \dfrac{5\pi}{3}$.

\textbf{Étape 2 et 3 :} sur le cercle trigonométrique, $\cos x$ est l'\emph{abscisse} du point image. On cherche les points dont l'abscisse est supérieure ou égale à $\frac{1}{2}$ : c'est la partie du cercle située à droite de la droite verticale passant par $\frac{1}{2}$, c'est-à-dire l'arc allant de $-\frac{\pi}{3}$ à $\frac{\pi}{3}$ en passant par $0$.

\begin{popfigure}
\begin{tikzpicture}[scale=2.2]
\draw[popline] (0,0) circle (1);
\draw[->] (-1.2,0) -- (1.2,0) node[right] {$x$};
\draw[->] (0,-1.2) -- (0,1.2) node[above] {$y$};
\draw[popblue, thick] (0.5,-1.2) -- (0.5,1.2) node[above] {$x=\frac{1}{2}$};
\fill[popblue!20] (0,0) -- (0.5, {sqrt(3)/2}) arc (60:-60:1) -- cycle;
\draw[popblue, very thick] (0.5, {sqrt(max(0,3))/2}) arc (60:-60:1);
\draw[popblue, thick] (0,0) -- (0.5, {sqrt(max(0,3))/2});
\draw[popblue, thick] (0,0) -- (0.5, {-sqrt(max(0,3))/2});
\node[popblue] at (0.75, 0) {$\cos x \geq \frac{1}{2}$};
\foreach \a/\t in {60/$\frac{\pi}{3}$, 300/$\frac{5\pi}{3}$}
  \draw[popdark] (\a:1) -- ++(\a:0.15) node[shift={(\a:0.35)}] {\t};
\draw[popdark] (0:1) -- ++(0:0.15) node[shift={(0:0.35)}] {$0$};
\end{tikzpicture}
\legende{Inéquation $\cos x \geq \frac{1}{2}$ sur le cercle trigonométrique}
\end{popfigure}

Test de sécurité : pour $x = 0$, $\cos 0 = 1 \geq \frac{1}{2}$ : vrai, $0$ est bien solution. Pour $x = \pi$, $\cos\pi = -1 < \frac{1}{2}$ : $\pi$ n'est pas solution. La zone solution est donc bien autour de $0$.

\textbf{Étape 4 :} sur $[0\,;\,2\pi]$, cet arc correspond à deux morceaux : de $0$ à $\frac{\pi}{3}$, puis de $\frac{5\pi}{3}$ à $2\pi$. L'inégalité est large, donc les bornes $\frac{\pi}{3}$ et $\frac{5\pi}{3}$ sont incluses.

\textbf{Conclusion :} $S = \left[0\,;\,\dfrac{\pi}{3}\right] \cup \left[\dfrac{5\pi}{3}\,;\,2\pi\right]$.
\end{exoresolu}

\begin{exoresolu}[Inéquation avec sinus]
Énoncé : résoudre dans $[0\,;\,2\pi]$ l'inéquation $\sin x \geq \dfrac{\sqrt{3}}{2}$.
\tcblower
\textbf{Solution détaillée.}

\textbf{Étape 1 :} équation associée $\sin x = \dfrac{\sqrt{3}}{2}$. Solutions dans $[0\,;\,2\pi]$ : $x = \dfrac{\pi}{3}$ et $x = \dfrac{2\pi}{3}$.

\textbf{Étape 2 et 3 :} sur le cercle, $\sin x$ est l'\emph{ordonnée}. On cherche les points au-dessus de la droite horizontale $y = \frac{\sqrt{3}}{2}$. C'est l'arc allant de $\frac{\pi}{3}$ à $\frac{2\pi}{3}$ en passant par $\frac{\pi}{2}$.

\begin{popfigure}
\begin{tikzpicture}[scale=2.2]
\draw[popline] (0,0) circle (1);
\draw[->] (-1.2,0) -- (1.2,0) node[right] {$x$};
\draw[->] (0,-1.2) -- (0,1.2) node[above] {$y$};
\draw[poporange, thick] (-1.2, {sqrt(max(0,3))/2}) -- (1.2, {sqrt(max(0,3))/2}) node[right] {$y=\frac{\sqrt{3}}{2}$};
\fill[poporange!20] (0,0) -- ({0.5}, {sqrt(3)/2}) arc (60:120:1) -- cycle;
\draw[poporange, very thick] ({0.5}, {sqrt(max(0,3))/2}) arc (60:120:1);
\draw[poporange, thick] (0,0) -- ({0.5}, {sqrt(max(0,3))/2});
\draw[poporange, thick] (0,0) -- ({-0.5}, {sqrt(max(0,3))/2});
\node[poporange] at (0, 0.7) {$\sin x \geq \frac{\sqrt{3}}{2}$};
\foreach \a/\t in {60/$\frac{\pi}{3}$, 120/$\frac{2\pi}{3}$}
  \draw[popdark] (\a:1) -- ++(\a:0.15) node[shift={(\a:0.35)}] {\t};
\end{tikzpicture}
\legende{Inéquation $\sin x \geq \frac{\sqrt{3}}{2}$ sur le cercle trigonométrique}
\end{popfigure}

Test : $\sin\frac{\pi}{2} = 1 \geq \frac{\sqrt{3}}{2}$ : vrai. $\sin\pi = 0 < \frac{\sqrt{3}}{2}$ : $\pi$ n'est pas solution.

\textbf{Étape 4 :} l'inégalité est large, bornes incluses.

\textbf{Conclusion :} $S = \left[\dfrac{\pi}{3}\,;\,\dfrac{2\pi}{3}\right]$.
\end{exoresolu}

\courssub{Applications techniques}

\begin{methode}[Modéliser et justifier dans un problème technique]
\begin{enumerate}
\item \textbf{Lire} la situation et identifier la grandeur périodique ou l'angle en jeu (rotation d'une poulie, inclinaison d'un panneau solaire, hauteur d'un point d'une roue).
\item \textbf{Traduire} les données en équation ou inéquation trigonométrique.
\item \textbf{Résoudre} avec les méthodes précédentes, en respectant l'intervalle imposé par le contexte (un tour complet, une journée...).
\item \textbf{Interpréter} chaque solution dans la situation et \textbf{conclure par une phrase} rédigée en français, avec l'unité si nécessaire.
\end{enumerate}
Réflexe Probatoire : la réponse finale doit toujours répondre à la question posée, pas seulement donner $x$.
\end{methode}

\begin{exoresolu}[Hauteur d'un point d'une roue de machine]
Énoncé : dans un atelier de Douala, un point $M$ est fixé sur le bord d'une roue de rayon $1$~m tournant autour de son centre $O$. La hauteur (en mètres) de $M$ par rapport au centre est donnée par $h(x) = \sin x$, où $x$ est l'angle de rotation en radians, $x\in[0\,;\,2\pi]$. Déterminer les positions de la roue pour lesquelles le point $M$ se trouve à une hauteur supérieure ou égale à $\dfrac{\sqrt{3}}{2}$~m au-dessus du centre.
\tcblower
\textbf{Solution détaillée.}

\textbf{Étape 1 et 2 :} on cherche les angles $x\in[0\,;\,2\pi]$ tels que :
\[ \sin x \geq \frac{\sqrt{3}}{2}. \]

\textbf{Étape 3 :} équation associée : $\sin x = \dfrac{\sqrt{3}}{2}$. Or $\sin\dfrac{\pi}{3} = \dfrac{\sqrt{3}}{2}$, donc :
\[ x = \frac{\pi}{3} + 2k\pi \quad \text{ou} \quad x = \pi - \frac{\pi}{3} + 2k\pi = \frac{2\pi}{3} + 2k\pi, \quad k\in\mathbb{Z}. \]
Dans $[0\,;\,2\pi]$ : $x = \frac{\pi}{3}$ et $x = \frac{2\pi}{3}$.

Sur le cercle trigonométrique, $\sin x$ est l'\emph{ordonnée} du point image : on cherche les points situés au-dessus de la droite horizontale d'ordonnée $\frac{\sqrt{3}}{2}$. C'est l'arc allant de $\frac{\pi}{3}$ à $\frac{2\pi}{3}$ en passant par $\frac{\pi}{2}$.

\begin{popfigure}
\begin{tikzpicture}
\begin{axis}[
    width=\linewidth, height=6cm,
    axis lines=middle,
    xlabel={$x$ (rad)}, ylabel={$h(x)=\sin x$},
    xtick={0,1.5708,3.1415,4.7123,6.2831},
    xticklabels={$0$, $\frac{\pi}{2}$, $\pi$, $\frac{3\pi}{2}$, $2\pi$},
    ytick={-1,0,0.866,1},
    yticklabels={$-1$, $0$, $\frac{\sqrt{3}}{2}$, $1$},
    domain=0:6.2831, samples=200,
    xmin=-0.3, xmax=6.6, ymin=-1.2, ymax=1.2,
    popcurve/.style={popblue, thick}
]
\addplot[popcurve] {sin(deg(x))};
\addplot[poporange, dashed, thick] coordinates {(0,0.866) (6.2831,0.866)} node[right] {$y=\frac{\sqrt{3}}{2}$};
\addplot[poporange, very thick, domain=1.047:2.094, samples=40] {sin(deg(x))};
\node[poporange] at (1.57, 0.4) {$S$};
\end{axis}
\end{tikzpicture}
\legende{Résolution graphique de $\sin x \geq \frac{\sqrt{3}}{2}$ pour la roue (zone hachurée)}
\end{popfigure}

Test : $\sin\frac{\pi}{2} = 1 \geq \frac{\sqrt{3}}{2}$ : vrai. $\sin\pi = 0 < \frac{\sqrt{3}}{2}$ : $\pi$ n'est pas solution. La zone est correcte.

Donc $S = \left[\dfrac{\pi}{3}\,;\,\dfrac{2\pi}{3}\right]$.

\textbf{Étape 4 — Interprétation :} le point $M$ se trouve à au moins $\frac{\sqrt{3}}{2}$~m (environ \SI{0,87}{m}) au-dessus du centre lorsque l'angle de rotation est compris entre $\frac{\pi}{3}$ rad (soit $60\degres$) et $\frac{2\pi}{3}$ rad (soit $120\degres$), c'est-à-dire pendant un tiers de tour autour de la position verticale haute.
\end{exoresolu}

\begin{attention}[Les 5 erreurs qui coûtent des points au Probatoire]
\begin{enumerate}
\item \textbf{Oublier la moitié des solutions.} Pour $\cos x = \cos\alpha$, il y a DEUX familles : $x = \alpha + 2k\pi$ \emph{et} $x = -\alpha + 2k\pi$. Pour $\sin x = \sin\alpha$ : $x = \alpha + 2k\pi$ \emph{et} $x = \pi - \alpha + 2k\pi$. N'en écrire qu'une fait perdre la moitié des points.
\item \textbf{Oublier de diviser $2k\pi$} dans les équations à angle composé : de $2x = \frac{\pi}{3} + 2k\pi$, on tire $x = \frac{\pi}{6} + k\pi$ (et non $x = \frac{\pi}{6} + 2k\pi$). Cela change toutes les solutions dans l'intervalle.
\item \textbf{Inverser les formules de transformation.} Dans $\cos p - \cos q = -2\sin\frac{p+q}{2}\sin\frac{p-q}{2}$, le signe est MOINS et ce sont des sinus ; dans $\sin p - \sin q$, c'est $2\cos\frac{p+q}{2}\,\sin\frac{p-q}{2}$. Apprends-les séparément et vérifie avec une valeur simple ($x=0$).
\item \textbf{Confondre cosinus et sinus sur le cercle} pour les inéquations : le cosinus se lit sur l'axe horizontal (abscisse), le sinus sur l'axe vertical (ordonnée). Fais toujours un test avec une valeur simple ($x=0$ ou $x=\frac{\pi}{2}$) pour valider ta zone.
\item \textbf{Mal gérer les bornes et l'intervalle.} Respecte les inégalités strictes ($<$ : bornes exclues) ou larges ($\leq$ : bornes incluses), et vérifie que chaque solution appartient bien à l'intervalle demandé : $2\pi$ est exclu si l'intervalle est $[0\,;\,2\pi[$. Conclus toujours par une phrase donnant l'ensemble $S$.
\end{enumerate}
\end{attention}

\newpage

\coursec{Exercices}

\courssub{Je vérifie mes connaissances}

\begin{exercice}[QCM : formules de transformation]
Pour chaque question, une seule réponse est exacte. La recopier sur la copie.
\begin{enumerate}
\item $\cos a \cos b$ est égal à :
\begin{enumerate}
\item[(a)] $\dfrac{1}{2}\big[\cos(a+b)+\cos(a-b)\big]$
\item[(b)] $\dfrac{1}{2}\big[\cos(a+b)-\cos(a-b)\big]$
\item[(c)] $\dfrac{1}{2}\big[\sin(a+b)+\sin(a-b)\big]$
\end{enumerate}
\item $\sin p + \sin q$ est égal à :
\begin{enumerate}
\item[(a)] $2\cos\dfrac{p+q}{2}\sin\dfrac{p-q}{2}$
\item[(b)] $2\sin\dfrac{p+q}{2}\cos\dfrac{p-q}{2}$
\item[(c)] $\sin(p+q)$
\end{enumerate}
\item L'équation $\cos x = \cos\dfrac{\pi}{6}$ a pour solutions dans $\mathbb{R}$ :
\begin{enumerate}
\item[(a)] $x = \dfrac{\pi}{6} + 2k\pi$ uniquement
\item[(b)] $x = \dfrac{\pi}{6} + 2k\pi$ ou $x = -\dfrac{\pi}{6} + 2k\pi$, $k \in \mathbb{Z}$
\item[(c)] $x = \dfrac{5\pi}{6} + 2k\pi$, $k \in \mathbb{Z}$
\end{enumerate}
\item $\cos 2x$ est égal à :
\begin{enumerate}
\item[(a)] $2\cos x$
\item[(b)] $1 - 2\sin^2 x$
\item[(c)] $2\sin x \cos x$
\end{enumerate}
\end{enumerate}
\end{exercice}

\begin{exercice}[Vrai ou faux, en justifiant]
Dire si chacune des affirmations suivantes est vraie ou fausse, en justifiant la réponse.
\begin{enumerate}
\item Pour tout réel $x$, $\sin 2x = 2\sin x$.
\item Pour tous réels $a$ et $b$, $\sin a \sin b = \dfrac{1}{2}\big[\cos(a-b) - \cos(a+b)\big]$.
\item L'équation $\sin x = 2$ admet une solution dans $\mathbb{R}$.
\item Pour tout réel $x$, $\cos x + \cos\left(x + \dfrac{\pi}{2}\right) = \sqrt{2}\cos\left(x + \dfrac{\pi}{4}\right)$.
\end{enumerate}
\end{exercice}

\begin{exercice}[Compléter les formules]
Recopier et compléter les égalités suivantes, valables pour tous réels $a$, $b$, $p$ et $q$ :
\begin{enumerate}
\item $\cos(a - b) = \ldots$
\item $\cos p + \cos q = 2\cos\dfrac{\ldots}{2}\cos\dfrac{\ldots}{2}$
\item $\cos p - \cos q = \ldots$
\item $\sin a \cos b = \dfrac{1}{2}\big[\ldots + \ldots\big]$
\end{enumerate}
\end{exercice}

\begin{exercice}[Calculs directs]
\begin{enumerate}
\item En utilisant $\dfrac{\pi}{12} = \dfrac{\pi}{3} - \dfrac{\pi}{4}$, calculer les valeurs exactes de $\cos\dfrac{\pi}{12}$ et $\sin\dfrac{\pi}{12}$.
\item En déduire la valeur exacte de $\tan\dfrac{\pi}{12}$.
\item Calculer $\cos\dfrac{7\pi}{12}$ en remarquant que $\dfrac{7\pi}{12} = \dfrac{\pi}{3} + \dfrac{\pi}{4}$.
\end{enumerate}
\end{exercice}

\courssub{Je m'entraîne}

\begin{exercice}[Transformer des produits en sommes]
Transformer chacun des produits suivants en une somme ou une différence :
\begin{enumerate}
\item $A = \cos 3x \cos 2x$ ;
\item $B = \sin 5x \sin x$ ;
\item $C = 2\sin 4x \cos x$ ;
\item $D = \cos\dfrac{x}{2}\sin\dfrac{3x}{2}$.
\end{enumerate}
\end{exercice}

\begin{exercice}[Transformer des sommes en produits]
\begin{enumerate}
\item Factoriser $E = \cos 4x + \cos 2x$ en un produit.
\item Résoudre sur $[0\,;2\pi]$ l'équation $\cos 4x + \cos 2x = 0$.
\item Factoriser $F = \sin 3x - \sin x$, puis résoudre sur $[0\,;\pi]$ l'équation $\sin 3x = \sin x$ en utilisant cette factorisation.
\end{enumerate}
\end{exercice}

\begin{exercice}[Équations trigonométriques élémentaires]
Résoudre dans $\mathbb{R}$, puis donner les solutions appartenant à $[0\,;2\pi]$ :
\begin{enumerate}
\item $\cos x = \dfrac{1}{2}$ ;
\item $\sin x = -\dfrac{\sqrt{2}}{2}$ ;
\item $\tan x = \sqrt{3}$ ;
\item $\cos\left(2x - \dfrac{\pi}{3}\right) = \dfrac{\sqrt{3}}{2}$.
\end{enumerate}
\end{exercice}

\begin{exercice}[Équations et inéquations]
\begin{enumerate}
\item Résoudre sur $[0\,;\pi]$ l'équation $\tan(2x) = 1$, en précisant au préalable les conditions d'existence.
\item Résoudre sur $[-\pi\,;\pi]$ l'inéquation $2\cos x - \sqrt{3} \geq 0$ et représenter les solutions sur le cercle trigonométrique.
\item Résoudre sur $[0\,;2\pi]$ l'inéquation $2\sin x - 1 < 0$.
\end{enumerate}
\end{exercice}

\courssub{Je me prépare au Probatoire}

\begin{exercice}[Identité remarquable et équation]
\textbf{Partie A : établissement d'une identité.}
\begin{enumerate}
\item Rappeler la formule donnant $\cos(a+b)$ et la formule de duplication $\cos 2x$ en fonction de $\cos x$.
\item En écrivant $3x = 2x + x$, démontrer que pour tout réel $x$ :
\[ \cos 3x = 4\cos^3 x - 3\cos x. \]
\end{enumerate}
\textbf{Partie B : résolution d'une équation.}

On considère l'équation $(E)$ : $4\cos^3 x - 3\cos x = \dfrac{1}{2}$, pour $x \in [0\,;\pi]$.
\begin{enumerate}
\item À l'aide de la partie A, montrer que résoudre $(E)$ revient à résoudre $\cos 3x = \dfrac{1}{2}$.
\item Résoudre dans $\mathbb{R}$ l'équation $\cos 3x = \dfrac{1}{2}$.
\item En déduire toutes les solutions de $(E)$ sur $[0\,;\pi]$.
\end{enumerate}
\end{exercice}

\begin{exercice}[Équation du second degré en $\cos x$]
Un technicien étudie un signal dont l'amplitude est donnée, pour $x \in [0\,;2\pi]$, par
\[ f(x) = 2\sin^2 x + 3\cos x. \]
On cherche les valeurs de $x$ pour lesquelles le signal s'annule, c'est-à-dire les solutions de l'équation $f(x) = 0$.
\begin{enumerate}
\item Montrer que pour tout réel $x$, $2\sin^2 x + 3\cos x = -2\cos^2 x + 3\cos x + 2$.
\item On pose $X = \cos x$. Montrer que résoudre $f(x) = 0$ revient à résoudre $-2X^2 + 3X + 2 = 0$ avec $X \in [-1\,;1]$.
\item Résoudre dans $\mathbb{R}$ l'équation $-2X^2 + 3X + 2 = 0$, puis écarter les solutions qui ne conviennent pas.
\item En déduire les solutions de $f(x) = 0$ sur $[0\,;2\pi]$.
\item Représenter ces solutions sur le cercle trigonométrique.
\end{enumerate}
\end{exercice}

\begin{exercice}[Hauteur d'eau dans un port]
Dans un port côtier, la hauteur d'eau (en mètres) au bout de $t$ heures après minuit est modélisée par :
\[ h(t) = 3 + 2\cos\left(\dfrac{\pi t}{6}\right), \quad t \in [0\,;24]. \]
Un bateau dont le tirant d'eau est de \SI{4}{m} ne peut entrer dans le port que lorsque la hauteur d'eau est d'au moins \SI{4}{m}.
\begin{enumerate}
\item Calculer $h(0)$, $h(6)$ et $h(12)$. Interpréter concrètement ces résultats.
\item Montrer que le bateau peut entrer dans le port si et seulement si $\cos\left(\dfrac{\pi t}{6}\right) \geq \dfrac{1}{2}$.
\item On pose $\theta = \dfrac{\pi t}{6}$. Résoudre sur $[0\,;4\pi]$ l'inéquation $\cos\theta \geq \dfrac{1}{2}$.
\item En déduire les intervalles de temps, sur $[0\,;24]$, pendant lesquels le bateau peut entrer dans le port.
\item Le bateau doit effectuer une manœuvre d'entrée qui dure 2 heures. Le capitaine souhaite entrer le plus tôt possible dans la journée. À quelle heure peut-il commencer sa manœuvre au plus tôt ? Justifier.
\end{enumerate}
\end{exercice}

\newpage

\coursec{Corrigés des exercices}

\coursec{Corrigés des exercices — Trigonométrie (Leçon 10)}

\begin{corrige}[Exercice 1 — QCM : formules de transformation]
\begin{enumerate}
\item Réponse \textbf{(a)}. La formule de transformation d'un produit en somme est
\[ \cos a \cos b = \frac{1}{2}\big[\cos(a+b)+\cos(a-b)\big]. \]
\item Réponse \textbf{(b)}. La formule de transformation d'une somme en produit est
\[ \sin p + \sin q = 2\sin\frac{p+q}{2}\cos\frac{p-q}{2}. \]
\item Réponse \textbf{(b)}. $\cos x = \cos\alpha \iff x = \alpha + 2k\pi$ ou $x = -\alpha + 2k\pi$, $k \in \mathbb{Z}$. Ici $\alpha = \dfrac{\pi}{6}$.
\item Réponse \textbf{(b)}. Les formules de duplication sont $\cos 2x = 2\cos^2 x - 1 = 1 - 2\sin^2 x = \cos^2 x - \sin^2 x$. La (c) donne $\sin 2x$, et la (a) est fausse en général.
\end{enumerate}
\end{corrige}

\begin{corrige}[Exercice 2 — Vrai ou faux, en justifiant]
\begin{enumerate}
\item \textbf{Faux.} La formule de duplication est $\sin 2x = 2\sin x \cos x$. Contre-exemple : pour $x = \dfrac{\pi}{2}$, $\sin 2x = \sin\pi = 0$ alors que $2\sin x = 2$.
\item \textbf{Vrai.} C'est la formule de transformation d'un produit en somme :
\[ \sin a \sin b = \frac{1}{2}\big[\cos(a-b) - \cos(a+b)\big]. \]
\item \textbf{Faux.} Pour tout réel $x$, $-1 \leq \sin x \leq 1$, donc $\sin x = 2$ n'a aucune solution dans $\mathbb{R}$.
\item \textbf{Vrai.} Avec $p = x$ et $q = x + \dfrac{\pi}{2}$ :
\begin{align*}
\cos x + \cos\left(x + \frac{\pi}{2}\right) &= 2\cos\frac{2x+\frac{\pi}{2}}{2}\cos\frac{-\frac{\pi}{2}}{2} = 2\cos\left(x+\frac{\pi}{4}\right)\cos\frac{\pi}{4} = \sqrt{2}\cos\left(x+\frac{\pi}{4}\right).
\end{align*}
\end{enumerate}
\end{corrige}

\begin{corrige}[Exercice 3 — Compléter les formules]
\begin{enumerate}
\item $\cos(a-b) = \cos a \cos b + \sin a \sin b$.
\item $\cos p + \cos q = 2\cos\dfrac{p+q}{2}\cos\dfrac{p-q}{2}$.
\item $\cos p - \cos q = -2\sin\dfrac{p+q}{2}\sin\dfrac{p-q}{2}$.
\item $\sin a \cos b = \dfrac{1}{2}\big[\sin(a+b) + \sin(a-b)\big]$.
\end{enumerate}
\end{corrige}

\begin{corrige}[Exercice 4 — Calculs directs]
\begin{enumerate}
\item Avec $\dfrac{\pi}{12} = \dfrac{\pi}{3} - \dfrac{\pi}{4}$ et les formules d'addition :
\begin{align*}
\cos\frac{\pi}{12} &= \cos\frac{\pi}{3}\cos\frac{\pi}{4} + \sin\frac{\pi}{3}\sin\frac{\pi}{4} = \frac{1}{2}\cdot\frac{\sqrt{2}}{2} + \frac{\sqrt{3}}{2}\cdot\frac{\sqrt{2}}{2} = \frac{\sqrt{2}+\sqrt{6}}{4}.\\
\sin\frac{\pi}{12} &= \sin\frac{\pi}{3}\cos\frac{\pi}{4} - \cos\frac{\pi}{3}\sin\frac{\pi}{4} = \frac{\sqrt{3}}{2}\cdot\frac{\sqrt{2}}{2} - \frac{1}{2}\cdot\frac{\sqrt{2}}{2} = \frac{\sqrt{6}-\sqrt{2}}{4}.
\end{align*}
\item $\tan\dfrac{\pi}{12} = \dfrac{\sin\frac{\pi}{12}}{\cos\frac{\pi}{12}} = \dfrac{\sqrt{6}-\sqrt{2}}{\sqrt{6}+\sqrt{2}} = \dfrac{(\sqrt{6}-\sqrt{2})^2}{6-2} = \dfrac{8 - 4\sqrt{3}}{4} = 2 - \sqrt{3}$.
\item $\cos\dfrac{7\pi}{12} = \cos\left(\dfrac{\pi}{3}+\dfrac{\pi}{4}\right) = \cos\dfrac{\pi}{3}\cos\dfrac{\pi}{4} - \sin\dfrac{\pi}{3}\sin\dfrac{\pi}{4} = \dfrac{\sqrt{2}-\sqrt{6}}{4}$.
On note que $\dfrac{7\pi}{12} = \dfrac{\pi}{2} + \dfrac{\pi}{12}$, donc $\cos\dfrac{7\pi}{12} = -\sin\dfrac{\pi}{12} = \dfrac{\sqrt{2}-\sqrt{6}}{4}$, ce qui confirme le résultat.
\end{enumerate}
\end{corrige}

\begin{corrige}[Exercice 5 — Transformer des produits en sommes]
\begin{enumerate}
\item $A = \cos 3x \cos 2x = \dfrac{1}{2}\big[\cos 5x + \cos x\big]$.
\item $B = \sin 5x \sin x = \dfrac{1}{2}\big[\cos 4x - \cos 6x\big]$.
\item $C = 2\sin 4x \cos x = 2\times\dfrac{1}{2}\big[\sin 5x + \sin 3x\big] = \sin 5x + \sin 3x$.
\item $D = \cos\dfrac{x}{2}\sin\dfrac{3x}{2} = \sin\dfrac{3x}{2}\cos\dfrac{x}{2} = \dfrac{1}{2}\big[\sin 2x + \sin x\big]$.
\end{enumerate}
\textbf{Point de vigilance :} pour $D$, on réécrit le produit sous la forme $\sin a \cos b$ avant d'appliquer la formule, et $\sin\left(\dfrac{3x}{2}-\dfrac{x}{2}\right) = \sin x$.
\end{corrige}

\begin{corrige}[Exercice 6 — Transformer des sommes en produits]
\begin{enumerate}
\item $E = \cos 4x + \cos 2x = 2\cos\dfrac{4x+2x}{2}\cos\dfrac{4x-2x}{2} = 2\cos 3x \cos x$.
\item $\cos 4x + \cos 2x = 0 \iff 2\cos 3x \cos x = 0 \iff \cos 3x = 0$ ou $\cos x = 0$.
\begin{itemize}
\item $\cos x = 0 \iff x = \dfrac{\pi}{2} + k\pi$ : sur $[0\,;2\pi]$, $x = \dfrac{\pi}{2}$ ou $x = \dfrac{3\pi}{2}$.
\item $\cos 3x = 0 \iff 3x = \dfrac{\pi}{2} + k\pi \iff x = \dfrac{\pi}{6} + \dfrac{k\pi}{3}$. Sur $[0\,;2\pi]$ : $x \in \left\{\dfrac{\pi}{6}\,;\dfrac{\pi}{2}\,;\dfrac{5\pi}{6}\,;\dfrac{7\pi}{6}\,;\dfrac{3\pi}{2}\,;\dfrac{11\pi}{6}\right\}$.
\end{itemize}
En réunissant : $\mathcal{S} = \left\{\dfrac{\pi}{6}\,;\dfrac{\pi}{2}\,;\dfrac{5\pi}{6}\,;\dfrac{7\pi}{6}\,;\dfrac{3\pi}{2}\,;\dfrac{11\pi}{6}\right\}$.
\item $F = \sin 3x - \sin x = 2\cos\dfrac{3x+x}{2}\sin\dfrac{3x-x}{2} = 2\cos 2x \sin x$.\\
$\sin 3x = \sin x \iff 2\cos 2x \sin x = 0 \iff \cos 2x = 0$ ou $\sin x = 0$.
\begin{itemize}
\item $\sin x = 0$ : sur $[0\,;\pi]$, $x = 0$ ou $x = \pi$.
\item $\cos 2x = 0 \iff 2x = \dfrac{\pi}{2} + k\pi \iff x = \dfrac{\pi}{4} + \dfrac{k\pi}{2}$ : sur $[0\,;\pi]$, $x = \dfrac{\pi}{4}$ ou $x = \dfrac{3\pi}{4}$.
\end{itemize}
$\mathcal{S} = \left\{0\,;\dfrac{\pi}{4}\,;\dfrac{3\pi}{4}\,;\pi\right\}$.
\end{enumerate}
\end{corrige}

\begin{corrige}[Exercice 7 — Équations trigonométriques élémentaires]
\begin{enumerate}
\item $\cos x = \dfrac{1}{2} = \cos\dfrac{\pi}{3} \iff x = \dfrac{\pi}{3} + 2k\pi$ ou $x = -\dfrac{\pi}{3} + 2k\pi$, $k\in\mathbb{Z}$. Sur $[0\,;2\pi]$ : $x = \dfrac{\pi}{3}$ ou $x = \dfrac{5\pi}{3}$.
\item $\sin x = -\dfrac{\sqrt{2}}{2} = \sin\left(-\dfrac{\pi}{4}\right) \iff x = -\dfrac{\pi}{4} + 2k\pi$ ou $x = \pi + \dfrac{\pi}{4} + 2k\pi = \dfrac{5\pi}{4} + 2k\pi$. Sur $[0\,;2\pi]$ : $x = \dfrac{5\pi}{4}$ ou $x = \dfrac{7\pi}{4}$.
\item $\tan x = \sqrt{3} = \tan\dfrac{\pi}{3} \iff x = \dfrac{\pi}{3} + k\pi$, $k\in\mathbb{Z}$. Sur $[0\,;2\pi]$ : $x = \dfrac{\pi}{3}$ ou $x = \dfrac{4\pi}{3}$.
\item $\cos\left(2x - \dfrac{\pi}{3}\right) = \dfrac{\sqrt{3}}{2} = \cos\dfrac{\pi}{6}$ donc :
\[ 2x - \frac{\pi}{3} = \frac{\pi}{6} + 2k\pi \quad\text{ou}\quad 2x - \frac{\pi}{3} = -\frac{\pi}{6} + 2k\pi \]
\[ \iff x = \frac{\pi}{4} + k\pi \quad\text{ou}\quad x = \frac{\pi}{12} + k\pi, \quad k\in\mathbb{Z}. \]
Sur $[0\,;2\pi]$ : $x \in \left\{\dfrac{\pi}{12}\,;\dfrac{\pi}{4}\,;\dfrac{13\pi}{12}\,;\dfrac{5\pi}{4}\right\}$.
\end{enumerate}
\textbf{Point de vigilance :} pour l'équation en $\cos(2x-\frac{\pi}{3})$, on résout d'abord en l'angle $2x-\frac{\pi}{3}$, puis on divise par 2 \emph{y compris le terme} $2k\pi$.
\end{corrige}

\begin{corrige}[Exercice 8 — Équations et inéquations]
\begin{enumerate}
\item \textbf{Conditions d'existence :} $\tan(2x)$ existe si et seulement si $2x \neq \dfrac{\pi}{2} + k\pi$, c'est-à-dire $x \neq \dfrac{\pi}{4} + \dfrac{k\pi}{2}$. Sur $[0\,;\pi]$, les valeurs interdites sont $\dfrac{\pi}{4}$ et $\dfrac{3\pi}{4}$.\\
$\tan(2x) = 1 = \tan\dfrac{\pi}{4} \iff 2x = \dfrac{\pi}{4} + k\pi \iff x = \dfrac{\pi}{8} + \dfrac{k\pi}{2}$.\\
Sur $[0\,;\pi]$ : $x = \dfrac{\pi}{8}$, $\dfrac{5\pi}{8}$ (avec $k=0,1$) ; $k=2$ donne $\dfrac{9\pi}{8} > \pi$. Aucune n'est une valeur interdite. $\mathcal{S} = \left\{\dfrac{\pi}{8}\,;\dfrac{5\pi}{8}\right\}$.
\item $2\cos x - \sqrt{3} \geq 0 \iff \cos x \geq \dfrac{\sqrt{3}}{2}$. Sur le cercle trigonométrique, les points dont l'abscisse est supérieure ou égale à $\dfrac{\sqrt{3}}{2}$ correspondent aux angles entre $-\dfrac{\pi}{6}$ et $\dfrac{\pi}{6}$. Sur $[-\pi\,;\pi]$ : $\mathcal{S} = \left[-\dfrac{\pi}{6}\,;\dfrac{\pi}{6}\right]$.
\begin{popfigure}
\begin{tikzpicture}[scale=1.6]
\draw[popline] (-1.3,0) -- (1.3,0); \draw[popline] (0,-1.3) -- (0,1.3);
\draw[popdark] (0,0) circle (1);
\draw[poporange, very thick] (30:1) arc (30:-30:1);
\draw[dashed, popblue] (0.866,1.1) -- (0.866,-1.1);
\fill[popink] (30:1) circle (0.03) node[above right] {$\frac{\pi}{6}$};
\fill[popink] (-30:1) circle (0.03) node[below right] {$-\frac{\pi}{6}$};
\node[popblue] at (0.55,1.15) {$x=\frac{\sqrt{3}}{2}$};
\end{tikzpicture}
\legende{Cercle trigonométrique : ensemble des solutions de $\cos x \geq \frac{\sqrt{3}}{2}$ sur $[-\pi;\pi]$.}
\end{popfigure}
\item $2\sin x - 1 < 0 \iff \sin x < \dfrac{1}{2}$. Sur $[0\,;2\pi]$, $\sin x = \dfrac{1}{2}$ pour $x = \dfrac{\pi}{6}$ et $x = \dfrac{5\pi}{6}$, et $\sin x < \dfrac{1}{2}$ en dehors de l'arc supérieur :
\[ \mathcal{S} = \left[0\,;\frac{\pi}{6}\right[ \cup \left]\frac{5\pi}{6}\,;2\pi\right]. \]
\end{enumerate}
\textbf{Point de vigilance :} pour une inéquation stricte, les bornes où l'égalité est réalisée sont exclues (crochets ouverts).
\end{corrige}

\begin{corrige}[Exercice 9 — Identité remarquable et équation]
\textbf{Partie A.}
\begin{enumerate}
\item $\cos(a+b) = \cos a\cos b - \sin a\sin b$ et $\cos 2x = 2\cos^2 x - 1$.
\item Pour tout réel $x$ :
\begin{align*}
\cos 3x &= \cos(2x+x) = \cos 2x \cos x - \sin 2x \sin x\\
&= (2\cos^2 x - 1)\cos x - (2\sin x\cos x)\sin x\\
&= 2\cos^3 x - \cos x - 2\sin^2 x \cos x\\
&= 2\cos^3 x - \cos x - 2(1-\cos^2 x)\cos x \quad (\text{car } \sin^2 x = 1 - \cos^2 x)\\
&= 2\cos^3 x - \cos x - 2\cos x + 2\cos^3 x = 4\cos^3 x - 3\cos x.
\end{align*}
\end{enumerate}
\textbf{Partie B.}
\begin{enumerate}
\item D'après la partie A, $4\cos^3 x - 3\cos x = \cos 3x$, donc $(E) \iff \cos 3x = \dfrac{1}{2}$.
\item $\cos 3x = \dfrac{1}{2} = \cos\dfrac{\pi}{3} \iff 3x = \dfrac{\pi}{3} + 2k\pi$ ou $3x = -\dfrac{\pi}{3} + 2k\pi$, c'est-à-dire
\[ x = \frac{\pi}{9} + \frac{2k\pi}{3} \quad\text{ou}\quad x = -\frac{\pi}{9} + \frac{2k\pi}{3}, \quad k\in\mathbb{Z}. \]
\item Sur $[0\,;\pi]$ : pour la première famille, $k=0$ donne $\dfrac{\pi}{9}$, $k=1$ donne $\dfrac{7\pi}{9}$ ; pour la seconde, $k=1$ donne $\dfrac{5\pi}{9}$. Les autres valeurs de $k$ sortent de l'intervalle.
\[ \mathcal{S} = \left\{\frac{\pi}{9}\,;\frac{5\pi}{9}\,;\frac{7\pi}{9}\right\}. \]
\end{enumerate}
\textbf{Barème indicatif (sur 5) :} A1 : 0,5 ; A2 : 1,5 ; B1 : 0,5 ; B2 : 1,5 ; B3 : 1.\\
\textbf{Points de vigilance :} citer $\sin 2x = 2\sin x\cos x$ et $\sin^2 x = 1-\cos^2 x$ ; ne pas oublier la deuxième famille de solutions ; vérifier que chaque solution appartient bien à $[0\,;\pi]$.
\end{corrige}

\begin{corrige}[Exercice 10 — Équation du second degré en $\cos x$]
\begin{enumerate}
\item Pour tout réel $x$, $\sin^2 x = 1 - \cos^2 x$, donc
\[ 2\sin^2 x + 3\cos x = 2(1-\cos^2 x) + 3\cos x = -2\cos^2 x + 3\cos x + 2. \]
\item Avec $X = \cos x$, $f(x) = 0 \iff -2X^2 + 3X + 2 = 0$. Comme pour tout réel $x$, $\cos x \in [-1\,;1]$, on ne retient que les solutions $X \in [-1\,;1]$.
\item $\Delta = 3^2 - 4\times(-2)\times 2 = 9 + 16 = 25 > 0$, donc deux racines :
\[ X_1 = \frac{-3 - 5}{-4} = 2, \qquad X_2 = \frac{-3 + 5}{-4} = -\frac{1}{2}. \]
$X_1 = 2 \notin [-1\,;1]$ est à écarter ; on conserve $X = -\dfrac{1}{2}$.
\item $\cos x = -\dfrac{1}{2} = \cos\dfrac{2\pi}{3} \iff x = \dfrac{2\pi}{3} + 2k\pi$ ou $x = -\dfrac{2\pi}{3} + 2k\pi$. Sur $[0\,;2\pi]$ :
\[ \mathcal{S} = \left\{\frac{2\pi}{3}\,;\frac{4\pi}{3}\right\}. \]
\item Cercle trigonométrique :
\begin{popfigure}
\begin{tikzpicture}[scale=1.6]
\draw[popline] (-1.3,0) -- (1.3,0); \draw[popline] (0,-1.3) -- (0,1.3);
\draw[popdark] (0,0) circle (1);
\draw[dashed, popblue] (-0.5,1.1) -- (-0.5,-1.1);
\fill[poporange] (120:1) circle (0.04) node[above left] {$\frac{2\pi}{3}$};
\fill[poporange] (240:1) circle (0.04) node[below left] {$\frac{4\pi}{3}$};
\node[popblue] at (-0.95,1.15) {$x=-\frac{1}{2}$};
\end{tikzpicture}
\legende{Cercle trigonométrique : solutions de $\cos x = -\frac{1}{2}$ sur $[0;2\pi]$.}
\end{popfigure}
\end{enumerate}
\textbf{Barème indicatif (sur 5) :} 1 : 1 ; 2 : 0,5 ; 3 : 1,5 ; 4 : 1,5 ; 5 : 0,5.\\
\textbf{Points de vigilance :} exiger la condition $X \in [-1\,;1]$ (justifier le rejet de $X=2$) ; donner les deux familles de solutions de $\cos x = -\frac12$ ; conclure sur l'intervalle demandé.
\end{corrige}

\begin{corrige}[Exercice 11 — Hauteur d'eau dans un port]
\begin{enumerate}
\item $h(0) = 3 + 2\cos 0 = 5$ ; $h(6) = 3 + 2\cos\pi = 3 - 2 = 1$ ; $h(12) = 3 + 2\cos 2\pi = 5$.\\
Interprétation : à minuit et à midi, la hauteur d'eau est maximale (\SI{5}{m}, pleine mer) ; à 6 h, elle est minimale (\SI{1}{m}, basse mer). La période de la marée est de 12 h.
\item Le bateau peut entrer si et seulement si $h(t) \geq 4$ :
\[ 3 + 2\cos\left(\frac{\pi t}{6}\right) \geq 4 \iff 2\cos\left(\frac{\pi t}{6}\right) \geq 1 \iff \cos\left(\frac{\pi t}{6}\right) \geq \frac{1}{2}. \]
\item Pour $t \in [0\,;24]$, $\theta = \dfrac{\pi t}{6} \in [0\,;4\pi]$. Sur $[0\,;2\pi]$, $\cos\theta \geq \dfrac{1}{2} \iff \theta \in \left[0\,;\dfrac{\pi}{3}\right] \cup \left[\dfrac{5\pi}{3}\,;2\pi\right]$. Par périodicité $2\pi$, sur $[0\,;4\pi]$ :
\[ \theta \in \left[0\,;\frac{\pi}{3}\right] \cup \left[\frac{5\pi}{3}\,;\frac{7\pi}{3}\right] \cup \left[\frac{11\pi}{3}\,;4\pi\right]. \]
\item Comme $t = \dfrac{6\theta}{\pi}$, on multiplie les bornes par $\dfrac{6}{\pi}$ :
\[ t \in [0\,;2] \cup [10\,;14] \cup [22\,;24]. \]
Le bateau peut entrer entre 0 h et 2 h, entre 10 h et 14 h, et entre 22 h et minuit.
\item La manœuvre dure 2 heures et doit tenir entièrement dans un intervalle d'entrée possible. Le premier intervalle est $[0\,;2]$, de durée exactement 2 h : le capitaine peut commencer sa manœuvre \textbf{à minuit (0 h)} au plus tôt, pour terminer à 2 h.
\end{enumerate}
\textbf{Barème indicatif (sur 5) :} 1 : 1 ; 2 : 0,5 ; 3 : 1,5 ; 4 : 1 ; 5 : 1.\\
\textbf{Points de vigilance :} justifier que $\theta$ décrit $[0\,;4\pi]$ quand $t$ décrit $[0\,;24]$ ; ne pas oublier la périodicité (deux marées hautes par jour) ; pour la question 5, vérifier que la durée de la manœuvre tient dans l'intervalle choisi.
\end{corrige}

\newpage

\coursec{Fiche bilan}

\begin{popfigure}
\begin{center}
\begin{tikzpicture}[mindmap, grow cyclic, every node/.style={concept, text=white, font=\small},
root concept/.append style={concept color=popblue, font=\bfseries},
level 1/.append style={level distance=3.6cm, sibling angle=60, font=\scriptsize},
level 1 concept/.append style={concept color=poporange}]
\node[root concept]{Trigonométrie : transformations et équations}
child[concept color=popgreen]{ node{Formules d'addition et duplication} }
child[concept color=poporange]{ node{Produit $\to$ somme} }
child[concept color=poppurple]{ node{Somme $\to$ produit} }
child[concept color=popblue]{ node{$\cos x=a$, $\sin x=a$, $\tan x=a$} }
child[concept color=popgold]{ node{Équations $a\cos x+b\sin x=c$} }
child[concept color=poppink]{ node{Inéquations trigonométriques} };
\end{tikzpicture}
\end{center}
\legende{Carte mentale de la leçon : trigonométrie en Première technique.}
\end{popfigure}

\begin{bilan}
\textbf{Définitions et formules clés.}
\begin{itemize}
\item \cle{Formules d'addition} : $\cos(a+b)=\cos a\cos b-\sin a\sin b$ ; $\cos(a-b)=\cos a\cos b+\sin a\sin b$ ; $\sin(a+b)=\sin a\cos b+\cos a\sin b$ ; $\sin(a-b)=\sin a\cos b-\cos a\sin b$.
\item \cle{Duplication} : $\cos(2a)=\cos^2 a-\sin^2 a=2\cos^2 a-1=1-2\sin^2 a$ ; $\sin(2a)=2\sin a\cos a$ ; $\tan(2a)=\dfrac{2\tan a}{1-\tan^2 a}$.
\item \cle{Linéarisation} : $\cos^2 a=\dfrac{1+\cos 2a}{2}$ ; $\sin^2 a=\dfrac{1-\cos 2a}{2}$.
\end{itemize}

\textbf{Transformations produit $\leftrightarrow$ somme.}
\begin{center}
\begin{tabularx}{\linewidth}{|X|X|}
\hline
\rowcolor{popblueL} \textbf{Produit en somme} & \textbf{Somme en produit} \\
\hline
$\cos a\cos b=\dfrac{1}{2}[\cos(a+b)+\cos(a-b)]$ & $\cos p+\cos q=2\cos\dfrac{p+q}{2}\cos\dfrac{p-q}{2}$ \\
\hline
$\sin a\sin b=\dfrac{1}{2}[\cos(a-b)-\cos(a+b)]$ & $\cos p-\cos q=-2\sin\dfrac{p+q}{2}\sin\dfrac{p-q}{2}$ \\
\hline
$\sin a\cos b=\dfrac{1}{2}[\sin(a+b)+\sin(a-b)]$ & $\sin p+\sin q=2\sin\dfrac{p+q}{2}\cos\dfrac{p-q}{2}$ \\
\hline
 & $\sin p-\sin q=2\cos\dfrac{p+q}{2}\sin\dfrac{p-q}{2}$ \\
\hline
\end{tabularx}
\end{center}

\textbf{Équations élémentaires.}
\begin{itemize}
\item $\cos x=\cos a \iff x=a+2k\pi$ ou $x=-a+2k\pi$, $k\in\mathbb{Z}$.
\item $\sin x=\sin a \iff x=a+2k\pi$ ou $x=\pi-a+2k\pi$, $k\in\mathbb{Z}$.
\item $\tan x=\tan a \iff x=a+k\pi$, $k\in\mathbb{Z}$.
\item $\cos x=a$ (ou $\sin x=a$) : solutions seulement si $-1\leq a\leq 1$.
\end{itemize}

\textbf{Méthodes express.}
\begin{itemize}
\item Équation en $\cos x$ et $\sin x$ du premier degré : poser $X=\cos x$ ou $X=\sin x$, résoudre, puis revenir à $x$.
\item $a\cos x+b\sin x=c$ : écrire $a\cos x+b\sin x=\sqrt{a^2+b^2}\cos(x-\varphi)$ avec $\cos\varphi=\dfrac{a}{\sqrt{a^2+b^2}}$, $\sin\varphi=\dfrac{b}{\sqrt{a^2+b^2}}$.
\item Inéquation : placer les solutions sur le \cle{cercle trigonométrique}, repérer l'arc solution, puis écrire l'ensemble des solutions avec $k\in\mathbb{Z}$.
\item Toujours conclure dans l'\cle{intervalle demandé} (souvent $[0\,;2\pi[$ ou $[-\pi\,;\pi]$).
\end{itemize}
\end{bilan}

\begin{aretenir}[Checklist avant le Probatoire]
\begin{itemize}
\item[$\square$] Je connais par cœur les formules d'addition de $\cos$ et $\sin$.
\item[$\square$] Je sais écrire $\cos 2a$ sous ses trois formes.
\item[$\square$] Je sais transformer un produit $\cos a\cos b$, $\sin a\sin b$, $\sin a\cos b$ en somme.
\item[$\square$] Je sais transformer $\cos p+\cos q$ et $\sin p+\sin q$ en produits.
\item[$\square$] Je sais résoudre $\cos x=\cos a$, $\sin x=\sin a$, $\tan x=\tan a$ sans oublier le $k\in\mathbb{Z}$.
\item[$\square$] Je vérifie que $|a|\leq 1$ avant d'affirmer que $\cos x=a$ a des solutions.
\item[$\square$] Je sais factoriser $a\cos x+b\sin x$ en $\sqrt{a^2+b^2}\cos(x-\varphi)$.
\item[$\square$] Je sais utiliser le cercle trigonométrique pour résoudre une inéquation.
\item[$\square$] Je donne les solutions dans l'intervalle exigé par l'énoncé.
\item[$\square$] Je vérifie mes solutions en remplaçant dans l'équation de départ.
\end{itemize}
\end{aretenir}

\findecours

\end{document}
